%%%%%%%%%%%%%%%%%%%%%%%%%%%%%%%%%%%%%%%%%%%%%%%%%%%%%%%%%%%%%%%%%%%%%%%%
% Portada con estructura de artículo científico: cabecera institucional,
% título y autoría alineados a la izquierda, resumen, palabras clave y
% clasificación, y un emblema calculado numéricamente (líneas de campo y
% equipotenciales de un conductor cargado con punta, figuras/datos/punta_*).
% Todos los bloques tienen coordenadas fijas en la página A4 (21 x 29,7 cm)
% para que ningún elemento se superponga a otro.
% Requiere en main.tex: \Universidad, \Facultad, \Titulacion, \Asignatura,
% \NumTema, \TituloTema, \SubtituloTema, \tipotrabajo, \miNombre,
% \miProfesor, \miUbicacion, \Afiliacion y el texto \ResumenPortada.
%%%%%%%%%%%%%%%%%%%%%%%%%%%%%%%%%%%%%%%%%%%%%%%%%%%%%%%%%%%%%%%%%%%%%%%%

\begin{titlepage}
\thispagestyle{empty}
\tikzexternaldisable
\begin{tikzpicture}[remember picture, overlay]
  \coordinate (O) at (current page.south west);

  % --- Bandas superior e inferior en degradado --------------------------
  \shade[left color=azulprof, right color=azulcielo!70]
    ($(O)+(0,29.25)$) rectangle ($(O)+(21,29.7)$);
  \shade[left color=azulcielo!45, right color=white]
    ($(O)+(0,29.15)$) rectangle ($(O)+(21,29.25)$);
  \shade[left color=azulprof, right color=azulcielo!70]
    ($(O)+(0,0)$) rectangle ($(O)+(21,0.3)$);
  % velo muy tenue en la cabecera
  \shade[top color=azulfondo, bottom color=white]
    ($(O)+(0,24.9)$) rectangle ($(O)+(21,29.15)$);

  % --- Cabecera institucional -------------------------------------------
  \node[anchor=north west, inner sep=0] at ($(O)+(2.5,28.3)$)
    {\includegraphics[height=1.2cm]{logo_usal.png}};
  \node[anchor=north east, inner sep=0] at ($(O)+(18.5,28.3)$)
    {\includegraphics[height=1.2cm]{logo_facultad.png}};
  \node[anchor=north, align=center, inner sep=0, text width=7.2cm] at ($(O)+(10.5,27.95)$)
    {{\color{azulprof}\scshape\textls[110]{\Universidad}}\\[3pt]
     {\color{grispizarra}\footnotesize\scshape\textls[50]{\Facultad\enspace·\enspace\Titulacion}}};
  \draw[azulclaro, line width=0.45pt] ($(O)+(2.5,26.55)$) -- ($(O)+(18.5,26.55)$);

  % --- Asignatura, título y subtítulo -----------------------------------
  \node[anchor=north west, align=left, inner sep=0, text width=16cm] at ($(O)+(2.5,25.95)$)
    {{\color{azulmedio}\scshape\textls[140]{\Asignatura\enspace·\enspace\NumTema}}};
  \node[anchor=north west, align=left, inner sep=0, text width=16cm,
        font=\bfseries\fontsize{29}{35}\selectfont, text=azulprof] at ($(O)+(2.5,25.05)$)
    {\TituloTema};
  \shade[left color=azulmedio, right color=white]
    ($(O)+(2.5,22.38)$) rectangle ($(O)+(10.5,22.44)$);
  \node[anchor=north west, align=left, inner sep=0, text width=15.6cm,
        font=\itshape\large, text=black!78] at ($(O)+(2.5,22.05)$)
    {\SubtituloTema};

  % --- Autoría y afiliación ----------------------------------------------
  \node[anchor=north west, align=left, inner sep=0, text width=7.6cm, font=\large] at ($(O)+(2.5,20.75)$)
    {{\scshape\small\color{azulmedio}\textls[80]{Autor}}\\[2pt]
     {\color{black!88}\miNombre\textsuperscript{\,a}}};
  \node[anchor=north west, align=left, inner sep=0, text width=7.6cm, font=\large] at ($(O)+(10.9,20.75)$)
    {{\scshape\small\color{azulmedio}\textls[80]{Profesor}}\\[2pt]
     {\color{black!88}\miProfesor\textsuperscript{\,a}}};
  \node[anchor=north west, align=left, inner sep=0, text width=16cm, font=\footnotesize] at ($(O)+(2.5,19.55)$)
    {\color{grispizarra}\textsuperscript{a}\,\Afiliacion};

  % --- Resumen (izquierda) y emblema (derecha) ---------------------------
  \node[anchor=north west, inner sep=0] at ($(O)+(2.5,18.6)$) {%
    \begin{tcolorbox}[enhanced, width=10.1cm, sharp corners, frame hidden, boxrule=0pt,
        interior style={left color=azulfondo, right color=white},
        borderline west={2pt}{0pt}{azulmedio},
        left=9pt, right=6pt, top=6pt, bottom=6pt]
      {\small\scshape\color{azulprof}\textls[90]{Resumen}}\par\smallskip
      {\footnotesize\ResumenPortada}
    \end{tcolorbox}};

  \begin{scope}[shift={($(O)+(15.7,15.4)$)}, x=1.06cm, y=1.06cm]
    \begin{scope}
      \clip[rounded corners=3pt] (-2.6,-2.0) rectangle (2.6,2.0);
      \fill[inner color=white, outer color=azulfondo] (-2.6,-2.0) rectangle (2.6,2.0);
      \tikzset{numLC/.style={azulacero!85, line width=0.42pt},
               numEQ/.style={azulcielo!80, line width=0.3pt, densely dashed},
               numFL/.style={-{Stealth[length=3pt, width=2.5pt]}, azulacero, line width=0.42pt}}
      % Conductor con punta - equipotenciales
% Generado por scripts/generar_datos.py (no editar a mano).
\EQ{(2.341,-2.040) (2.370,-2.012) (2.400,-1.982) (2.430,-1.951) (2.459,-1.920) (2.486,-1.890) (2.513,-1.860) (2.535,-1.835) (2.564,-1.800) (2.589,-1.770) (2.610,-1.743) (2.636,-1.710) (2.640,-1.705)}
\EQ{(-2.640,-1.913) (-2.610,-1.949) (-2.583,-1.980) (-2.556,-2.010) (-2.535,-2.033) (-2.529,-2.040)}
\EQ{(2.640,1.705) (2.613,1.740) (2.589,1.770) (2.565,1.799) (2.539,1.830) (2.513,1.860) (2.490,1.886) (2.460,1.919) (2.431,1.950) (2.402,1.980) (2.372,2.010) (2.341,2.040)}
\EQ{(-2.529,2.040) (-2.550,2.017) (-2.580,1.983) (-2.609,1.950) (-2.634,1.920) (-2.640,1.913)}
\EQ{(1.657,-2.040) (1.695,-2.015) (1.724,-1.995) (1.755,-1.973) (1.785,-1.952) (1.815,-1.930) (1.845,-1.907) (1.875,-1.884) (1.905,-1.860) (1.935,-1.835) (1.959,-1.815) (1.994,-1.785) (2.025,-1.758) (2.055,-1.730) (2.085,-1.702) (2.108,-1.680) (2.130,-1.658) (2.153,-1.635) (2.175,-1.612) (2.196,-1.590) (2.220,-1.564) (2.250,-1.531) (2.277,-1.500) (2.302,-1.470) (2.325,-1.442) (2.351,-1.410) (2.370,-1.385) (2.397,-1.350) (2.415,-1.325) (2.440,-1.290) (2.460,-1.260) (2.480,-1.230) (2.500,-1.200) (2.518,-1.170) (2.535,-1.143) (2.554,-1.110) (2.571,-1.080) (2.588,-1.050) (2.604,-1.020) (2.619,-0.990) (2.634,-0.960) (2.640,-0.947)}
\EQ{(-2.640,-1.202) (-2.625,-1.230) (-2.610,-1.258) (-2.595,-1.284) (-2.574,-1.320) (-2.556,-1.350) (-2.537,-1.380) (-2.520,-1.406) (-2.497,-1.440) (-2.477,-1.470) (-2.455,-1.500) (-2.433,-1.530) (-2.415,-1.554) (-2.387,-1.590) (-2.363,-1.620) (-2.340,-1.648) (-2.312,-1.680) (-2.286,-1.710) (-2.265,-1.733) (-2.244,-1.755) (-2.220,-1.780) (-2.190,-1.810) (-2.160,-1.839) (-2.138,-1.860) (-2.115,-1.881) (-2.089,-1.905) (-2.055,-1.934) (-2.025,-1.960) (-2.000,-1.980) (-1.965,-2.008) (-1.935,-2.031) (-1.923,-2.040)}
\EQ{(2.640,0.947) (2.626,0.975) (2.611,1.005) (2.596,1.035) (2.580,1.064) (2.565,1.091) (2.550,1.117) (2.528,1.155) (2.509,1.185) (2.490,1.215) (2.470,1.245) (2.450,1.275) (2.430,1.304) (2.408,1.335) (2.385,1.365) (2.363,1.395) (2.340,1.424) (2.315,1.455) (2.295,1.479) (2.265,1.514) (2.238,1.545) (2.210,1.575) (2.182,1.605) (2.160,1.628) (2.138,1.650) (2.115,1.673) (2.092,1.695) (2.070,1.716) (2.044,1.740) (2.011,1.770) (1.980,1.797) (1.950,1.823) (1.923,1.845) (1.890,1.872) (1.860,1.895) (1.830,1.918) (1.800,1.941) (1.770,1.963) (1.746,1.980) (1.710,2.005) (1.680,2.025) (1.657,2.040)}
\EQ{(-1.923,2.040) (-1.950,2.019) (-1.980,1.996) (-2.010,1.972) (-2.037,1.950) (-2.070,1.921) (-2.100,1.895) (-2.130,1.868) (-2.154,1.845) (-2.185,1.815) (-2.215,1.785) (-2.244,1.755) (-2.265,1.733) (-2.286,1.710) (-2.310,1.683) (-2.338,1.650) (-2.363,1.620) (-2.385,1.593) (-2.411,1.560) (-2.430,1.535) (-2.455,1.500) (-2.475,1.472) (-2.497,1.440) (-2.517,1.410) (-2.535,1.383) (-2.556,1.350) (-2.574,1.320) (-2.592,1.290) (-2.609,1.260) (-2.625,1.230) (-2.640,1.202)}
\EQ{(0.464,-2.040) (0.495,-2.032) (0.525,-2.025) (0.555,-2.017) (0.585,-2.008) (0.615,-2.000) (0.645,-1.991) (0.675,-1.982) (0.705,-1.972) (0.735,-1.962) (0.765,-1.952) (0.795,-1.941) (0.825,-1.930) (0.855,-1.919) (0.885,-1.908) (0.915,-1.896) (0.945,-1.884) (0.975,-1.871) (1.005,-1.858) (1.034,-1.845) (1.065,-1.831) (1.095,-1.817) (1.125,-1.803) (1.155,-1.788) (1.185,-1.772) (1.215,-1.757) (1.245,-1.741) (1.274,-1.725) (1.300,-1.710) (1.326,-1.695) (1.365,-1.672) (1.395,-1.654) (1.425,-1.635) (1.455,-1.616) (1.485,-1.596) (1.515,-1.576) (1.545,-1.555) (1.575,-1.534) (1.600,-1.515) (1.635,-1.489) (1.659,-1.470) (1.695,-1.442) (1.725,-1.417) (1.751,-1.395) (1.785,-1.366) (1.815,-1.339) (1.845,-1.311) (1.867,-1.290) (1.890,-1.267) (1.912,-1.245) (1.935,-1.222) (1.956,-1.200) (1.980,-1.174) (2.010,-1.140) (2.036,-1.110) (2.061,-1.080) (2.085,-1.050) (2.108,-1.020) (2.130,-0.991) (2.152,-0.960) (2.173,-0.930) (2.190,-0.905) (2.213,-0.870) (2.231,-0.840) (2.249,-0.810) (2.265,-0.782) (2.280,-0.755) (2.295,-0.726) (2.310,-0.697) (2.325,-0.666) (2.340,-0.633) (2.354,-0.600) (2.367,-0.570) (2.378,-0.540) (2.389,-0.510) (2.400,-0.480) (2.409,-0.450) (2.419,-0.420) (2.427,-0.390) (2.435,-0.360) (2.442,-0.330) (2.448,-0.300) (2.454,-0.270) (2.460,-0.240) (2.464,-0.210) (2.468,-0.180) (2.472,-0.150) (2.475,-0.120) (2.477,-0.090) (2.478,-0.060) (2.479,-0.030) (2.479,-0.000) (2.479,0.030) (2.478,0.060) (2.477,0.090) (2.475,0.120) (2.472,0.150) (2.468,0.180) (2.464,0.210) (2.460,0.240) (2.454,0.270) (2.448,0.300) (2.442,0.330) (2.435,0.360) (2.427,0.390) (2.419,0.420) (2.409,0.450) (2.400,0.479) (2.389,0.510) (2.378,0.540) (2.367,0.570) (2.355,0.598) (2.341,0.630) (2.328,0.660) (2.313,0.690) (2.298,0.720) (2.283,0.750) (2.266,0.780) (2.250,0.808) (2.231,0.840) (2.213,0.870) (2.193,0.900) (2.175,0.927) (2.152,0.960) (2.131,0.990) (2.108,1.020) (2.085,1.050) (2.061,1.080) (2.040,1.105) (2.010,1.140) (1.983,1.170) (1.956,1.200) (1.935,1.222) (1.912,1.245) (1.890,1.267) (1.867,1.290) (1.845,1.311) (1.819,1.335) (1.786,1.365) (1.755,1.392) (1.725,1.417) (1.697,1.440) (1.665,1.466) (1.640,1.485) (1.605,1.511) (1.580,1.530) (1.545,1.555) (1.516,1.575) (1.485,1.596) (1.455,1.616) (1.425,1.635) (1.395,1.654) (1.365,1.672) (1.335,1.690) (1.305,1.707) (1.275,1.724) (1.246,1.740) (1.218,1.755) (1.190,1.770) (1.160,1.785) (1.130,1.800) (1.099,1.815) (1.067,1.830) (1.035,1.845) (1.005,1.858) (0.975,1.871) (0.945,1.884) (0.915,1.896) (0.885,1.908) (0.855,1.919) (0.825,1.930) (0.795,1.941) (0.765,1.952) (0.735,1.962) (0.705,1.972) (0.675,1.982) (0.645,1.991) (0.615,2.000) (0.585,2.008) (0.555,2.017) (0.525,2.025) (0.495,2.032) (0.465,2.040) (0.464,2.040)}
\EQ{(-0.923,2.040) (-0.960,2.030) (-0.990,2.021) (-1.020,2.012) (-1.050,2.003) (-1.080,1.993) (-1.110,1.982) (-1.140,1.971) (-1.170,1.960) (-1.200,1.948) (-1.230,1.936) (-1.260,1.923) (-1.290,1.910) (-1.320,1.896) (-1.350,1.881) (-1.380,1.866) (-1.410,1.851) (-1.440,1.835) (-1.470,1.818) (-1.500,1.801) (-1.526,1.785) (-1.560,1.764) (-1.590,1.745) (-1.619,1.725) (-1.650,1.704) (-1.680,1.682) (-1.710,1.660) (-1.740,1.637) (-1.770,1.613) (-1.797,1.590) (-1.830,1.562) (-1.860,1.535) (-1.890,1.507) (-1.913,1.485) (-1.935,1.463) (-1.957,1.440) (-1.980,1.416) (-2.010,1.383) (-2.039,1.350) (-2.064,1.320) (-2.085,1.294) (-2.111,1.260) (-2.130,1.235) (-2.155,1.200) (-2.175,1.171) (-2.196,1.140) (-2.215,1.110) (-2.234,1.080) (-2.250,1.052) (-2.265,1.026) (-2.280,0.999) (-2.295,0.971) (-2.310,0.942) (-2.324,0.915) (-2.338,0.885) (-2.352,0.855) (-2.365,0.825) (-2.378,0.795) (-2.390,0.765) (-2.401,0.735) (-2.412,0.705) (-2.423,0.675) (-2.433,0.645) (-2.442,0.615) (-2.451,0.585) (-2.459,0.555) (-2.468,0.525) (-2.475,0.495) (-2.482,0.465) (-2.489,0.435) (-2.495,0.405) (-2.500,0.375) (-2.506,0.345) (-2.511,0.315) (-2.515,0.285) (-2.519,0.255) (-2.522,0.225) (-2.525,0.195) (-2.528,0.165) (-2.530,0.135) (-2.532,0.105) (-2.533,0.075) (-2.534,0.045) (-2.534,0.015) (-2.534,-0.030) (-2.533,-0.060) (-2.532,-0.090) (-2.531,-0.120) (-2.529,-0.150) (-2.527,-0.180) (-2.524,-0.210) (-2.520,-0.240) (-2.517,-0.270) (-2.513,-0.300) (-2.508,-0.330) (-2.503,-0.360) (-2.498,-0.390) (-2.492,-0.420) (-2.485,-0.450) (-2.479,-0.480) (-2.471,-0.510) (-2.464,-0.540) (-2.455,-0.570) (-2.447,-0.600) (-2.437,-0.630) (-2.428,-0.660) (-2.417,-0.690) (-2.407,-0.720) (-2.395,-0.750) (-2.384,-0.780) (-2.371,-0.810) (-2.358,-0.840) (-2.345,-0.870) (-2.331,-0.900) (-2.316,-0.930) (-2.301,-0.960) (-2.285,-0.990) (-2.269,-1.020) (-2.251,-1.050) (-2.235,-1.078) (-2.215,-1.110) (-2.196,-1.140) (-2.176,-1.170) (-2.155,-1.200) (-2.133,-1.230) (-2.115,-1.255) (-2.088,-1.290) (-2.064,-1.320) (-2.040,-1.348) (-2.012,-1.380) (-1.985,-1.410) (-1.957,-1.440) (-1.935,-1.463) (-1.913,-1.485) (-1.890,-1.507) (-1.865,-1.530) (-1.832,-1.560) (-1.800,-1.588) (-1.770,-1.613) (-1.742,-1.635) (-1.710,-1.660) (-1.683,-1.680) (-1.650,-1.704) (-1.620,-1.725) (-1.590,-1.745) (-1.560,-1.764) (-1.530,-1.783) (-1.501,-1.800) (-1.470,-1.818) (-1.440,-1.835) (-1.410,-1.851) (-1.380,-1.866) (-1.350,-1.881) (-1.320,-1.896) (-1.290,-1.910) (-1.260,-1.923) (-1.232,-1.935) (-1.200,-1.948) (-1.170,-1.960) (-1.140,-1.971) (-1.110,-1.982) (-1.080,-1.993) (-1.050,-2.003) (-1.020,-2.012) (-0.990,-2.021) (-0.960,-2.030) (-0.930,-2.038) (-0.923,-2.040)}
\EQ{(-0.510,-1.726) (-0.480,-1.728) (-0.450,-1.730) (-0.420,-1.731) (-0.390,-1.732) (-0.360,-1.733) (-0.330,-1.733) (-0.300,-1.733) (-0.270,-1.733) (-0.240,-1.732) (-0.210,-1.731) (-0.180,-1.729) (-0.150,-1.727) (-0.120,-1.725) (-0.090,-1.723) (-0.060,-1.720) (-0.030,-1.716) (-0.000,-1.713) (0.030,-1.709) (0.060,-1.705) (0.090,-1.700) (0.120,-1.695) (0.150,-1.690) (0.180,-1.684) (0.210,-1.678) (0.240,-1.672) (0.270,-1.666) (0.300,-1.659) (0.330,-1.652) (0.360,-1.644) (0.390,-1.636) (0.420,-1.628) (0.449,-1.620) (0.480,-1.611) (0.510,-1.602) (0.540,-1.593) (0.570,-1.583) (0.600,-1.573) (0.630,-1.562) (0.660,-1.552) (0.690,-1.541) (0.718,-1.530) (0.750,-1.518) (0.780,-1.506) (0.810,-1.493) (0.840,-1.481) (0.870,-1.468) (0.898,-1.455) (0.930,-1.441) (0.960,-1.426) (0.990,-1.412) (1.020,-1.397) (1.050,-1.382) (1.080,-1.366) (1.110,-1.350) (1.138,-1.335) (1.164,-1.320) (1.191,-1.305) (1.230,-1.282) (1.260,-1.264) (1.290,-1.245) (1.320,-1.226) (1.350,-1.206) (1.380,-1.186) (1.410,-1.165) (1.440,-1.144) (1.466,-1.125) (1.500,-1.099) (1.525,-1.080) (1.560,-1.052) (1.590,-1.028) (1.617,-1.005) (1.650,-0.976) (1.680,-0.949) (1.710,-0.921) (1.740,-0.891) (1.770,-0.861) (1.800,-0.829) (1.830,-0.796) (1.857,-0.765) (1.882,-0.735) (1.905,-0.706) (1.929,-0.675) (1.950,-0.646) (1.971,-0.615) (1.991,-0.585) (2.010,-0.555) (2.028,-0.525) (2.044,-0.495) (2.060,-0.465) (2.075,-0.435) (2.089,-0.405) (2.101,-0.375) (2.113,-0.345) (2.124,-0.315) (2.134,-0.285) (2.143,-0.255) (2.151,-0.225) (2.158,-0.195) (2.163,-0.165) (2.168,-0.135) (2.172,-0.105) (2.175,-0.075) (2.177,-0.045) (2.178,-0.015) (2.178,0.030) (2.176,0.060) (2.174,0.090) (2.170,0.120) (2.166,0.150) (2.161,0.180) (2.154,0.210) (2.147,0.240) (2.138,0.270) (2.129,0.300) (2.119,0.330) (2.107,0.360) (2.095,0.390) (2.082,0.420) (2.068,0.450) (2.052,0.480) (2.036,0.510) (2.019,0.540) (2.001,0.570) (1.981,0.600) (1.961,0.630) (1.940,0.660) (1.920,0.687) (1.894,0.720) (1.875,0.744) (1.845,0.779) (1.818,0.810) (1.790,0.840) (1.761,0.870) (1.731,0.900) (1.700,0.930) (1.668,0.960) (1.635,0.989) (1.605,1.015) (1.581,1.035) (1.545,1.064) (1.515,1.088) (1.486,1.110) (1.455,1.133) (1.425,1.155) (1.395,1.176) (1.365,1.196) (1.337,1.215) (1.305,1.235) (1.275,1.254) (1.245,1.273) (1.216,1.290) (1.185,1.308) (1.155,1.325) (1.125,1.342) (1.095,1.358) (1.065,1.374) (1.035,1.390) (1.005,1.405) (0.975,1.419) (0.945,1.434) (0.915,1.447) (0.885,1.461) (0.855,1.474) (0.825,1.487) (0.795,1.500) (0.765,1.512) (0.735,1.524) (0.705,1.535) (0.675,1.546) (0.645,1.557) (0.615,1.568) (0.585,1.578) (0.555,1.588) (0.525,1.597) (0.495,1.607) (0.465,1.615) (0.435,1.624) (0.405,1.632) (0.375,1.640) (0.345,1.648) (0.315,1.655) (0.285,1.662) (0.255,1.669) (0.225,1.675) (0.195,1.681) (0.165,1.687) (0.135,1.693) (0.105,1.698) (0.075,1.702) (0.045,1.707) (0.015,1.711) (-0.015,1.715) (-0.045,1.718) (-0.075,1.721) (-0.105,1.724) (-0.135,1.726) (-0.165,1.728) (-0.195,1.730) (-0.225,1.731) (-0.255,1.732) (-0.285,1.733) (-0.315,1.733) (-0.345,1.733) (-0.375,1.733) (-0.405,1.732) (-0.435,1.731) (-0.465,1.729) (-0.495,1.727) (-0.525,1.724) (-0.555,1.722) (-0.585,1.718) (-0.615,1.715) (-0.645,1.710) (-0.675,1.706) (-0.705,1.701) (-0.735,1.695) (-0.765,1.690) (-0.795,1.683) (-0.825,1.676) (-0.855,1.669) (-0.885,1.661) (-0.915,1.653) (-0.945,1.644) (-0.975,1.635) (-1.005,1.625) (-1.035,1.615) (-1.065,1.604) (-1.095,1.592) (-1.125,1.580) (-1.155,1.567) (-1.185,1.554) (-1.215,1.540) (-1.245,1.525) (-1.275,1.510) (-1.305,1.494) (-1.335,1.477) (-1.365,1.459) (-1.395,1.441) (-1.425,1.421) (-1.455,1.401) (-1.485,1.380) (-1.515,1.358) (-1.545,1.335) (-1.575,1.311) (-1.601,1.290) (-1.635,1.260) (-1.665,1.232) (-1.695,1.203) (-1.725,1.173) (-1.755,1.141) (-1.783,1.110) (-1.808,1.080) (-1.830,1.053) (-1.856,1.020) (-1.875,0.994) (-1.899,0.960) (-1.919,0.930) (-1.939,0.900) (-1.957,0.870) (-1.975,0.840) (-1.992,0.810) (-2.008,0.780) (-2.023,0.750) (-2.037,0.720) (-2.051,0.690) (-2.064,0.660) (-2.076,0.630) (-2.088,0.600) (-2.099,0.570) (-2.110,0.540) (-2.119,0.510) (-2.129,0.480) (-2.137,0.450) (-2.145,0.420) (-2.152,0.390) (-2.159,0.360) (-2.166,0.330) (-2.171,0.300) (-2.176,0.270) (-2.181,0.240) (-2.185,0.210) (-2.188,0.180) (-2.191,0.150) (-2.194,0.120) (-2.196,0.090) (-2.197,0.060) (-2.198,0.030) (-2.198,-0.000) (-2.198,-0.030) (-2.197,-0.060) (-2.196,-0.090) (-2.194,-0.120) (-2.191,-0.150) (-2.188,-0.180) (-2.185,-0.210) (-2.181,-0.240) (-2.176,-0.270) (-2.171,-0.300) (-2.166,-0.330) (-2.159,-0.360) (-2.152,-0.390) (-2.145,-0.420) (-2.137,-0.450) (-2.129,-0.480) (-2.119,-0.510) (-2.110,-0.540) (-2.099,-0.570) (-2.088,-0.600) (-2.076,-0.630) (-2.064,-0.660) (-2.051,-0.690) (-2.037,-0.720) (-2.023,-0.750) (-2.008,-0.780) (-1.992,-0.810) (-1.975,-0.840) (-1.957,-0.870) (-1.939,-0.900) (-1.920,-0.929) (-1.899,-0.960) (-1.878,-0.990) (-1.860,-1.014) (-1.832,-1.050) (-1.808,-1.080) (-1.785,-1.107) (-1.756,-1.140) (-1.728,-1.170) (-1.698,-1.200) (-1.667,-1.230) (-1.635,-1.260) (-1.605,-1.286) (-1.575,-1.311) (-1.546,-1.335) (-1.515,-1.358) (-1.485,-1.380) (-1.455,-1.401) (-1.425,-1.421) (-1.396,-1.440) (-1.365,-1.459) (-1.335,-1.477) (-1.305,-1.494) (-1.275,-1.510) (-1.245,-1.525) (-1.215,-1.540) (-1.185,-1.554) (-1.155,-1.567) (-1.125,-1.580) (-1.095,-1.592) (-1.065,-1.604) (-1.035,-1.615) (-1.005,-1.625) (-0.975,-1.635) (-0.945,-1.644) (-0.915,-1.653) (-0.885,-1.661) (-0.855,-1.669) (-0.825,-1.676) (-0.795,-1.683) (-0.765,-1.690) (-0.735,-1.695) (-0.705,-1.701) (-0.675,-1.706) (-0.645,-1.710) (-0.615,-1.715) (-0.585,-1.718) (-0.555,-1.722) (-0.525,-1.724) (-0.510,-1.726)}
\EQ{(-0.525,-1.395) (-0.495,-1.397) (-0.465,-1.398) (-0.435,-1.398) (-0.405,-1.398) (-0.375,-1.398) (-0.345,-1.397) (-0.315,-1.396) (-0.285,-1.394) (-0.255,-1.392) (-0.225,-1.389) (-0.195,-1.386) (-0.165,-1.383) (-0.135,-1.379) (-0.105,-1.375) (-0.075,-1.370) (-0.045,-1.366) (-0.015,-1.360) (0.015,-1.355) (0.045,-1.349) (0.075,-1.342) (0.105,-1.336) (0.135,-1.329) (0.165,-1.321) (0.195,-1.314) (0.225,-1.306) (0.255,-1.298) (0.285,-1.289) (0.315,-1.280) (0.345,-1.271) (0.375,-1.262) (0.405,-1.252) (0.435,-1.242) (0.465,-1.232) (0.495,-1.221) (0.525,-1.210) (0.555,-1.199) (0.585,-1.188) (0.615,-1.176) (0.645,-1.164) (0.675,-1.152) (0.704,-1.140) (0.735,-1.127) (0.765,-1.113) (0.795,-1.100) (0.825,-1.086) (0.855,-1.072) (0.885,-1.058) (0.915,-1.043) (0.945,-1.028) (0.975,-1.013) (1.005,-0.998) (1.035,-0.982) (1.065,-0.965) (1.095,-0.949) (1.125,-0.932) (1.154,-0.915) (1.185,-0.897) (1.215,-0.878) (1.245,-0.860) (1.275,-0.841) (1.305,-0.821) (1.335,-0.801) (1.365,-0.781) (1.395,-0.760) (1.425,-0.738) (1.449,-0.720) (1.485,-0.693) (1.515,-0.669) (1.544,-0.645) (1.575,-0.619) (1.605,-0.593) (1.630,-0.570) (1.662,-0.540) (1.692,-0.510) (1.720,-0.480) (1.747,-0.450) (1.770,-0.424) (1.797,-0.390) (1.819,-0.360) (1.840,-0.330) (1.859,-0.300) (1.875,-0.273) (1.890,-0.245) (1.905,-0.213) (1.919,-0.180) (1.929,-0.150) (1.938,-0.120) (1.945,-0.090) (1.949,-0.060) (1.952,-0.030) (1.953,-0.000) (1.952,0.030) (1.949,0.060) (1.945,0.090) (1.938,0.120) (1.929,0.150) (1.919,0.180) (1.907,0.210) (1.892,0.240) (1.877,0.270) (1.860,0.299) (1.840,0.330) (1.819,0.360) (1.800,0.386) (1.773,0.420) (1.747,0.450) (1.725,0.475) (1.695,0.507) (1.665,0.537) (1.635,0.565) (1.605,0.593) (1.580,0.615) (1.545,0.645) (1.515,0.669) (1.489,0.690) (1.455,0.716) (1.429,0.735) (1.395,0.760) (1.366,0.780) (1.335,0.801) (1.305,0.821) (1.276,0.840) (1.245,0.860) (1.215,0.878) (1.185,0.897) (1.155,0.914) (1.128,0.930) (1.102,0.945) (1.075,0.960) (1.047,0.975) (1.020,0.990) (0.991,1.005) (0.962,1.020) (0.932,1.035) (0.902,1.050) (0.870,1.065) (0.840,1.079) (0.810,1.093) (0.780,1.107) (0.750,1.120) (0.720,1.133) (0.690,1.146) (0.660,1.158) (0.631,1.170) (0.600,1.182) (0.570,1.194) (0.540,1.205) (0.510,1.216) (0.480,1.227) (0.450,1.237) (0.420,1.247) (0.390,1.257) (0.360,1.267) (0.330,1.276) (0.300,1.285) (0.270,1.293) (0.240,1.302) (0.210,1.310) (0.180,1.318) (0.150,1.325) (0.120,1.332) (0.090,1.339) (0.060,1.346) (0.030,1.352) (-0.000,1.358) (-0.030,1.363) (-0.060,1.368) (-0.090,1.373) (-0.120,1.377) (-0.150,1.381) (-0.180,1.385) (-0.210,1.388) (-0.240,1.391) (-0.270,1.393) (-0.300,1.395) (-0.330,1.397) (-0.360,1.398) (-0.390,1.398) (-0.420,1.398) (-0.450,1.398) (-0.480,1.397) (-0.510,1.396) (-0.540,1.394) (-0.570,1.392) (-0.600,1.389) (-0.630,1.386) (-0.660,1.382) (-0.690,1.378) (-0.720,1.373) (-0.750,1.368) (-0.780,1.362) (-0.810,1.355) (-0.840,1.348) (-0.870,1.340) (-0.900,1.331) (-0.930,1.322) (-0.960,1.312) (-0.990,1.302) (-1.020,1.291) (-1.050,1.279) (-1.080,1.266) (-1.110,1.252) (-1.140,1.238) (-1.170,1.223) (-1.200,1.207) (-1.230,1.190) (-1.260,1.172) (-1.286,1.155) (-1.320,1.132) (-1.350,1.111) (-1.380,1.089) (-1.410,1.065) (-1.440,1.040) (-1.470,1.013) (-1.495,0.990) (-1.525,0.960) (-1.554,0.930) (-1.575,0.906) (-1.605,0.871) (-1.630,0.840) (-1.650,0.813) (-1.674,0.780) (-1.694,0.750) (-1.713,0.720) (-1.731,0.690) (-1.748,0.660) (-1.764,0.630) (-1.779,0.600) (-1.794,0.570) (-1.807,0.540) (-1.819,0.510) (-1.831,0.480) (-1.842,0.450) (-1.852,0.420) (-1.861,0.390) (-1.869,0.360) (-1.877,0.330) (-1.884,0.300) (-1.891,0.270) (-1.896,0.240) (-1.901,0.210) (-1.905,0.180) (-1.909,0.150) (-1.912,0.120) (-1.914,0.090) (-1.916,0.060) (-1.917,0.030) (-1.917,-0.000) (-1.917,-0.030) (-1.916,-0.060) (-1.914,-0.090) (-1.912,-0.120) (-1.909,-0.150) (-1.905,-0.180) (-1.901,-0.210) (-1.896,-0.240) (-1.891,-0.270) (-1.884,-0.300) (-1.877,-0.330) (-1.869,-0.360) (-1.861,-0.390) (-1.852,-0.420) (-1.842,-0.450) (-1.831,-0.480) (-1.819,-0.510) (-1.807,-0.540) (-1.794,-0.570) (-1.779,-0.600) (-1.764,-0.630) (-1.748,-0.660) (-1.731,-0.690) (-1.713,-0.720) (-1.695,-0.749) (-1.674,-0.780) (-1.653,-0.810) (-1.630,-0.840) (-1.606,-0.870) (-1.580,-0.900) (-1.560,-0.923) (-1.539,-0.945) (-1.515,-0.970) (-1.485,-0.999) (-1.462,-1.020) (-1.428,-1.050) (-1.395,-1.077) (-1.365,-1.100) (-1.335,-1.122) (-1.309,-1.140) (-1.275,-1.162) (-1.245,-1.181) (-1.215,-1.198) (-1.185,-1.215) (-1.156,-1.230) (-1.125,-1.245) (-1.095,-1.259) (-1.065,-1.272) (-1.035,-1.285) (-1.005,-1.296) (-0.975,-1.307) (-0.945,-1.317) (-0.915,-1.327) (-0.885,-1.336) (-0.855,-1.344) (-0.825,-1.351) (-0.795,-1.358) (-0.765,-1.365) (-0.735,-1.370) (-0.705,-1.376) (-0.675,-1.380) (-0.645,-1.384) (-0.615,-1.388) (-0.585,-1.391) (-0.555,-1.393) (-0.525,-1.395)}
\EQ{(-0.660,-1.111) (-0.630,-1.113) (-0.600,-1.115) (-0.570,-1.117) (-0.540,-1.118) (-0.510,-1.118) (-0.480,-1.117) (-0.450,-1.116) (-0.420,-1.115) (-0.390,-1.112) (-0.360,-1.109) (-0.330,-1.106) (-0.300,-1.102) (-0.270,-1.097) (-0.240,-1.092) (-0.210,-1.086) (-0.180,-1.080) (-0.150,-1.073) (-0.120,-1.066) (-0.090,-1.059) (-0.060,-1.051) (-0.030,-1.043) (-0.000,-1.035) (0.030,-1.026) (0.060,-1.017) (0.090,-1.008) (0.120,-0.998) (0.150,-0.988) (0.180,-0.978) (0.210,-0.968) (0.240,-0.958) (0.270,-0.947) (0.300,-0.936) (0.330,-0.925) (0.360,-0.914) (0.390,-0.903) (0.420,-0.891) (0.450,-0.880) (0.480,-0.868) (0.510,-0.856) (0.540,-0.844) (0.570,-0.831) (0.600,-0.819) (0.630,-0.806) (0.660,-0.793) (0.690,-0.780) (0.720,-0.767) (0.750,-0.753) (0.780,-0.740) (0.810,-0.726) (0.840,-0.712) (0.870,-0.698) (0.900,-0.684) (0.930,-0.669) (0.960,-0.654) (0.990,-0.640) (1.020,-0.625) (1.050,-0.609) (1.080,-0.594) (1.110,-0.578) (1.140,-0.562) (1.170,-0.546) (1.200,-0.529) (1.230,-0.513) (1.260,-0.496) (1.287,-0.480) (1.320,-0.460) (1.350,-0.442) (1.380,-0.424) (1.410,-0.405) (1.440,-0.386) (1.470,-0.366) (1.500,-0.346) (1.530,-0.325) (1.560,-0.303) (1.590,-0.280) (1.620,-0.256) (1.650,-0.231) (1.674,-0.210) (1.705,-0.180) (1.733,-0.150) (1.755,-0.123) (1.777,-0.090) (1.792,-0.060) (1.801,-0.030) (1.805,-0.000) (1.801,0.030) (1.792,0.060) (1.777,0.090) (1.757,0.120) (1.733,0.150) (1.710,0.175) (1.680,0.204) (1.657,0.225) (1.621,0.255) (1.590,0.280) (1.564,0.300) (1.530,0.325) (1.501,0.345) (1.470,0.366) (1.440,0.386) (1.410,0.405) (1.380,0.424) (1.350,0.442) (1.320,0.460) (1.290,0.478) (1.261,0.495) (1.235,0.510) (1.208,0.525) (1.181,0.540) (1.153,0.555) (1.125,0.570) (1.097,0.585) (1.068,0.600) (1.039,0.615) (1.009,0.630) (0.979,0.645) (0.949,0.660) (0.918,0.675) (0.887,0.690) (0.855,0.705) (0.825,0.719) (0.795,0.733) (0.765,0.746) (0.735,0.760) (0.705,0.773) (0.675,0.786) (0.645,0.799) (0.615,0.812) (0.585,0.825) (0.555,0.837) (0.525,0.850) (0.495,0.862) (0.465,0.874) (0.436,0.885) (0.405,0.897) (0.375,0.909) (0.345,0.920) (0.315,0.931) (0.285,0.942) (0.255,0.953) (0.225,0.963) (0.195,0.973) (0.165,0.983) (0.135,0.993) (0.105,1.003) (0.075,1.012) (0.045,1.021) (0.015,1.030) (-0.015,1.039) (-0.045,1.047) (-0.075,1.055) (-0.105,1.063) (-0.135,1.070) (-0.165,1.077) (-0.195,1.083) (-0.225,1.089) (-0.255,1.094) (-0.285,1.099) (-0.315,1.104) (-0.345,1.108) (-0.375,1.111) (-0.405,1.113) (-0.435,1.115) (-0.465,1.117) (-0.495,1.118) (-0.525,1.118) (-0.555,1.117) (-0.585,1.116) (-0.615,1.114) (-0.645,1.112) (-0.675,1.109) (-0.705,1.105) (-0.735,1.101) (-0.765,1.095) (-0.795,1.089) (-0.825,1.083) (-0.855,1.075) (-0.885,1.067) (-0.915,1.057) (-0.945,1.047) (-0.975,1.036) (-1.005,1.023) (-1.035,1.010) (-1.065,0.996) (-1.095,0.980) (-1.125,0.964) (-1.155,0.946) (-1.185,0.927) (-1.215,0.906) (-1.244,0.885) (-1.275,0.861) (-1.300,0.840) (-1.333,0.810) (-1.364,0.780) (-1.393,0.750) (-1.419,0.720) (-1.440,0.695) (-1.467,0.660) (-1.485,0.636) (-1.509,0.600) (-1.528,0.570) (-1.545,0.540) (-1.560,0.513) (-1.575,0.483) (-1.590,0.450) (-1.603,0.420) (-1.615,0.390) (-1.625,0.360) (-1.635,0.330) (-1.644,0.300) (-1.652,0.270) (-1.659,0.240) (-1.665,0.210) (-1.670,0.180) (-1.674,0.150) (-1.677,0.120) (-1.680,0.090) (-1.682,0.060) (-1.683,0.030) (-1.683,-0.000) (-1.683,-0.030) (-1.682,-0.060) (-1.680,-0.090) (-1.677,-0.120) (-1.674,-0.150) (-1.670,-0.180) (-1.665,-0.210) (-1.659,-0.240) (-1.652,-0.270) (-1.644,-0.300) (-1.635,-0.330) (-1.625,-0.360) (-1.615,-0.390) (-1.603,-0.420) (-1.590,-0.450) (-1.576,-0.480) (-1.561,-0.510) (-1.545,-0.540) (-1.530,-0.566) (-1.509,-0.600) (-1.489,-0.630) (-1.470,-0.656) (-1.444,-0.690) (-1.425,-0.713) (-1.395,-0.747) (-1.365,-0.779) (-1.335,-0.808) (-1.305,-0.835) (-1.282,-0.855) (-1.245,-0.884) (-1.215,-0.906) (-1.185,-0.927) (-1.157,-0.945) (-1.125,-0.964) (-1.095,-0.980) (-1.065,-0.996) (-1.035,-1.010) (-1.005,-1.023) (-0.976,-1.035) (-0.945,-1.047) (-0.915,-1.057) (-0.885,-1.067) (-0.855,-1.075) (-0.825,-1.083) (-0.795,-1.089) (-0.765,-1.095) (-0.735,-1.101) (-0.705,-1.105) (-0.675,-1.109) (-0.660,-1.111)}

      % Conductor con punta - líneas de campo
% Generado por scripts/generar_datos.py (no editar a mano).
\LC{(-0.388,0.905) (-0.381,0.941) (-0.375,0.976) (-0.370,1.012) (-0.365,1.048) (-0.360,1.083) (-0.356,1.119) (-0.353,1.155) (-0.349,1.191) (-0.347,1.227) (-0.344,1.263) (-0.342,1.298) (-0.340,1.334) (-0.338,1.370) (-0.337,1.406) (-0.335,1.442) (-0.334,1.478) (-0.333,1.514) (-0.333,1.550) (-0.332,1.586) (-0.332,1.622) (-0.332,1.658) (-0.332,1.694) (-0.332,1.730) (-0.332,1.766) (-0.332,1.802) (-0.333,1.838) (-0.333,1.874) (-0.334,1.910) (-0.335,1.946) (-0.336,1.982) (-0.336,2.000)}
\LC{(-0.651,0.929) (-0.653,0.965) (-0.655,1.000) (-0.658,1.036) (-0.661,1.072) (-0.665,1.108) (-0.669,1.144) (-0.673,1.180) (-0.678,1.215) (-0.683,1.251) (-0.688,1.287) (-0.693,1.322) (-0.698,1.358) (-0.704,1.393) (-0.710,1.429) (-0.716,1.464) (-0.722,1.500) (-0.729,1.535) (-0.735,1.571) (-0.742,1.606) (-0.749,1.641) (-0.756,1.677) (-0.763,1.712) (-0.770,1.747) (-0.777,1.782) (-0.785,1.818) (-0.792,1.853) (-0.800,1.888) (-0.808,1.923) (-0.815,1.958) (-0.823,1.993) (-0.826,2.000)}
\LC{(-0.888,0.884) (-0.899,0.918) (-0.911,0.953) (-0.923,0.987) (-0.934,1.021) (-0.947,1.054) (-0.959,1.088) (-0.971,1.122) (-0.984,1.156) (-0.997,1.189) (-1.010,1.223) (-1.023,1.257) (-1.036,1.290) (-1.049,1.323) (-1.063,1.357) (-1.077,1.390) (-1.091,1.423) (-1.104,1.457) (-1.118,1.490) (-1.133,1.523) (-1.147,1.556) (-1.161,1.589) (-1.176,1.622) (-1.190,1.655) (-1.205,1.688) (-1.219,1.721) (-1.234,1.753) (-1.249,1.786) (-1.264,1.819) (-1.279,1.852) (-1.294,1.884) (-1.309,1.917) (-1.324,1.950) (-1.339,1.983) (-1.349,2.000)}
\LC{(-1.097,0.786) (-1.117,0.816) (-1.136,0.846) (-1.156,0.876) (-1.176,0.907) (-1.196,0.937) (-1.216,0.966) (-1.236,0.996) (-1.256,1.026) (-1.276,1.056) (-1.297,1.085) (-1.317,1.115) (-1.338,1.145) (-1.359,1.174) (-1.379,1.203) (-1.400,1.233) (-1.421,1.262) (-1.442,1.291) (-1.463,1.320) (-1.485,1.350) (-1.506,1.379) (-1.527,1.408) (-1.548,1.437) (-1.570,1.466) (-1.591,1.494) (-1.613,1.523) (-1.635,1.552) (-1.656,1.581) (-1.678,1.609) (-1.700,1.638) (-1.722,1.667) (-1.743,1.695) (-1.765,1.724) (-1.787,1.753) (-1.809,1.781) (-1.831,1.809) (-1.853,1.838) (-1.875,1.866) (-1.898,1.895) (-1.920,1.923) (-1.942,1.951) (-1.964,1.980) (-1.987,2.000)}
\LC{(-1.274,0.641) (-1.300,0.666) (-1.326,0.690) (-1.352,0.715) (-1.379,0.740) (-1.405,0.764) (-1.432,0.788) (-1.458,0.813) (-1.485,0.837) (-1.512,0.861) (-1.538,0.885) (-1.565,0.909) (-1.592,0.933) (-1.619,0.957) (-1.646,0.981) (-1.673,1.005) (-1.700,1.028) (-1.727,1.052) (-1.754,1.076) (-1.781,1.099) (-1.809,1.123) (-1.836,1.146) (-1.863,1.170) (-1.891,1.193) (-1.918,1.216) (-1.946,1.240) (-1.973,1.263) (-2.001,1.286) (-2.028,1.309) (-2.056,1.332) (-2.084,1.355) (-2.112,1.378) (-2.139,1.401) (-2.167,1.424) (-2.195,1.447) (-2.223,1.469) (-2.251,1.492) (-2.279,1.515) (-2.307,1.537) (-2.335,1.560) (-2.363,1.582) (-2.391,1.605) (-2.419,1.627) (-2.447,1.650) (-2.475,1.672) (-2.504,1.695) (-2.532,1.717) (-2.560,1.739) (-2.589,1.761) (-2.600,1.776)}
\LC{(-1.407,0.462) (-1.438,0.481) (-1.469,0.498) (-1.501,0.516) (-1.532,0.534) (-1.563,0.551) (-1.595,0.569) (-1.626,0.586) (-1.658,0.604) (-1.689,0.621) (-1.721,0.638) (-1.753,0.656) (-1.784,0.673) (-1.816,0.690) (-1.847,0.707) (-1.879,0.724) (-1.911,0.741) (-1.943,0.758) (-1.975,0.775) (-2.006,0.792) (-2.038,0.809) (-2.070,0.825) (-2.102,0.842) (-2.134,0.859) (-2.166,0.875) (-2.198,0.892) (-2.230,0.908) (-2.262,0.925) (-2.294,0.941) (-2.326,0.957) (-2.358,0.974) (-2.390,0.990) (-2.422,1.006) (-2.454,1.022) (-2.487,1.038) (-2.519,1.055) (-2.551,1.071) (-2.583,1.086) (-2.600,1.097)}
\LC{(-1.493,0.261) (-1.527,0.271) (-1.562,0.280) (-1.597,0.290) (-1.631,0.300) (-1.666,0.310) (-1.701,0.319) (-1.735,0.329) (-1.770,0.339) (-1.805,0.348) (-1.839,0.358) (-1.874,0.367) (-1.909,0.377) (-1.943,0.386) (-1.978,0.396) (-2.013,0.405) (-2.048,0.415) (-2.082,0.424) (-2.117,0.434) (-2.152,0.443) (-2.187,0.452) (-2.222,0.462) (-2.256,0.471) (-2.291,0.480) (-2.326,0.489) (-2.361,0.498) (-2.396,0.507) (-2.430,0.517) (-2.465,0.526) (-2.500,0.535) (-2.535,0.544) (-2.570,0.553) (-2.600,0.562)}
\LC{(-1.529,0.044) (-1.565,0.045) (-1.601,0.046) (-1.637,0.048) (-1.673,0.049) (-1.709,0.050) (-1.745,0.052) (-1.781,0.053) (-1.817,0.055) (-1.853,0.057) (-1.889,0.058) (-1.925,0.060) (-1.961,0.061) (-1.997,0.063) (-2.033,0.064) (-2.068,0.066) (-2.104,0.067) (-2.140,0.069) (-2.176,0.070) (-2.212,0.072) (-2.248,0.073) (-2.284,0.075) (-2.320,0.076) (-2.356,0.078) (-2.392,0.079) (-2.428,0.081) (-2.464,0.082) (-2.500,0.084) (-2.536,0.085) (-2.572,0.087) (-2.600,0.088)}
\LC{(-1.512,-0.181) (-1.548,-0.187) (-1.583,-0.194) (-1.618,-0.200) (-1.654,-0.207) (-1.689,-0.213) (-1.725,-0.220) (-1.760,-0.226) (-1.795,-0.233) (-1.831,-0.240) (-1.866,-0.246) (-1.902,-0.253) (-1.937,-0.259) (-1.972,-0.266) (-2.008,-0.272) (-2.043,-0.279) (-2.079,-0.285) (-2.114,-0.292) (-2.150,-0.298) (-2.185,-0.305) (-2.220,-0.311) (-2.256,-0.317) (-2.291,-0.324) (-2.327,-0.330) (-2.362,-0.336) (-2.398,-0.343) (-2.433,-0.349) (-2.468,-0.355) (-2.504,-0.361) (-2.539,-0.367) (-2.575,-0.374) (-2.600,-0.380)}
\LC{(-1.445,-0.389) (-1.478,-0.404) (-1.511,-0.418) (-1.543,-0.433) (-1.576,-0.448) (-1.609,-0.462) (-1.642,-0.477) (-1.675,-0.492) (-1.708,-0.506) (-1.741,-0.521) (-1.774,-0.535) (-1.807,-0.549) (-1.840,-0.564) (-1.873,-0.578) (-1.906,-0.592) (-1.939,-0.607) (-1.972,-0.621) (-2.005,-0.635) (-2.038,-0.649) (-2.071,-0.663) (-2.105,-0.677) (-2.138,-0.691) (-2.171,-0.705) (-2.204,-0.719) (-2.237,-0.733) (-2.271,-0.747) (-2.304,-0.760) (-2.337,-0.774) (-2.370,-0.788) (-2.404,-0.801) (-2.437,-0.815) (-2.471,-0.828) (-2.504,-0.842) (-2.537,-0.855) (-2.571,-0.869) (-2.600,-0.882)}
\LC{(-1.330,-0.576) (-1.358,-0.599) (-1.386,-0.621) (-1.415,-0.643) (-1.443,-0.665) (-1.472,-0.687) (-1.500,-0.709) (-1.529,-0.731) (-1.558,-0.753) (-1.586,-0.774) (-1.615,-0.796) (-1.644,-0.818) (-1.673,-0.839) (-1.702,-0.861) (-1.731,-0.882) (-1.760,-0.903) (-1.789,-0.925) (-1.818,-0.946) (-1.847,-0.967) (-1.876,-0.988) (-1.905,-1.009) (-1.934,-1.030) (-1.964,-1.051) (-1.993,-1.072) (-2.022,-1.093) (-2.052,-1.114) (-2.081,-1.134) (-2.111,-1.155) (-2.140,-1.176) (-2.170,-1.196) (-2.199,-1.217) (-2.229,-1.237) (-2.258,-1.258) (-2.288,-1.278) (-2.318,-1.299) (-2.347,-1.319) (-2.377,-1.339) (-2.407,-1.359) (-2.437,-1.380) (-2.467,-1.400) (-2.497,-1.420) (-2.526,-1.440) (-2.556,-1.460) (-2.586,-1.480) (-2.600,-1.493)}
\LC{(-1.170,-0.735) (-1.192,-0.763) (-1.214,-0.791) (-1.237,-0.820) (-1.260,-0.848) (-1.282,-0.876) (-1.305,-0.904) (-1.328,-0.931) (-1.351,-0.959) (-1.374,-0.987) (-1.397,-1.015) (-1.420,-1.042) (-1.443,-1.070) (-1.466,-1.097) (-1.489,-1.125) (-1.513,-1.152) (-1.536,-1.179) (-1.560,-1.206) (-1.583,-1.234) (-1.607,-1.261) (-1.631,-1.288) (-1.655,-1.315) (-1.678,-1.342) (-1.702,-1.369) (-1.726,-1.396) (-1.750,-1.423) (-1.774,-1.449) (-1.798,-1.476) (-1.823,-1.503) (-1.847,-1.529) (-1.871,-1.556) (-1.895,-1.583) (-1.920,-1.609) (-1.944,-1.636) (-1.968,-1.662) (-1.993,-1.689) (-2.017,-1.715) (-2.042,-1.741) (-2.066,-1.768) (-2.091,-1.794) (-2.116,-1.820) (-2.140,-1.846) (-2.165,-1.873) (-2.190,-1.899) (-2.215,-1.925) (-2.239,-1.951) (-2.264,-1.977) (-2.289,-2.000)}
\LC{(-0.973,-0.852) (-0.988,-0.885) (-1.002,-0.918) (-1.017,-0.951) (-1.032,-0.983) (-1.047,-1.016) (-1.063,-1.048) (-1.078,-1.081) (-1.094,-1.113) (-1.110,-1.146) (-1.126,-1.178) (-1.142,-1.210) (-1.158,-1.242) (-1.174,-1.274) (-1.191,-1.306) (-1.207,-1.338) (-1.224,-1.370) (-1.241,-1.402) (-1.258,-1.434) (-1.275,-1.466) (-1.292,-1.497) (-1.309,-1.529) (-1.326,-1.561) (-1.343,-1.592) (-1.360,-1.624) (-1.378,-1.655) (-1.395,-1.687) (-1.413,-1.718) (-1.431,-1.750) (-1.448,-1.781) (-1.466,-1.812) (-1.484,-1.844) (-1.501,-1.875) (-1.519,-1.906) (-1.537,-1.937) (-1.555,-1.969) (-1.573,-2.000) (-1.579,-2.000)}
\LC{(-0.745,-0.919) (-0.751,-0.954) (-0.756,-0.990) (-0.763,-1.025) (-0.769,-1.061) (-0.776,-1.096) (-0.783,-1.131) (-0.790,-1.166) (-0.798,-1.202) (-0.806,-1.237) (-0.814,-1.272) (-0.822,-1.307) (-0.831,-1.342) (-0.839,-1.377) (-0.848,-1.412) (-0.857,-1.447) (-0.866,-1.481) (-0.876,-1.516) (-0.885,-1.551) (-0.895,-1.586) (-0.904,-1.620) (-0.914,-1.655) (-0.924,-1.690) (-0.934,-1.724) (-0.944,-1.759) (-0.954,-1.793) (-0.964,-1.828) (-0.975,-1.862) (-0.985,-1.897) (-0.996,-1.931) (-1.006,-1.965) (-1.017,-2.000) (-1.021,-2.000)}
\LC{(-0.487,-0.923) (-0.485,-0.959) (-0.482,-0.995) (-0.480,-1.031) (-0.479,-1.067) (-0.477,-1.103) (-0.476,-1.139) (-0.476,-1.175) (-0.475,-1.211) (-0.475,-1.247) (-0.476,-1.283) (-0.476,-1.319) (-0.477,-1.355) (-0.478,-1.391) (-0.479,-1.427) (-0.481,-1.463) (-0.482,-1.499) (-0.484,-1.535) (-0.486,-1.571) (-0.488,-1.606) (-0.491,-1.642) (-0.493,-1.678) (-0.496,-1.714) (-0.499,-1.750) (-0.501,-1.786) (-0.504,-1.822) (-0.507,-1.858) (-0.511,-1.894) (-0.514,-1.929) (-0.517,-1.965) (-0.521,-2.000)}
\LC{(-0.215,-0.847) (-0.203,-0.881) (-0.192,-0.915) (-0.182,-0.950) (-0.172,-0.985) (-0.163,-1.019) (-0.155,-1.054) (-0.147,-1.090) (-0.139,-1.125) (-0.132,-1.160) (-0.125,-1.195) (-0.119,-1.231) (-0.112,-1.266) (-0.107,-1.302) (-0.101,-1.337) (-0.095,-1.373) (-0.090,-1.409) (-0.085,-1.444) (-0.080,-1.480) (-0.076,-1.516) (-0.071,-1.551) (-0.067,-1.587) (-0.063,-1.623) (-0.059,-1.659) (-0.055,-1.694) (-0.051,-1.730) (-0.047,-1.766) (-0.044,-1.802) (-0.040,-1.838) (-0.037,-1.873) (-0.033,-1.909) (-0.030,-1.945) (-0.027,-1.981) (-0.025,-2.000)}
\LC{(0.096,-0.718) (0.109,-0.752) (0.122,-0.785) (0.135,-0.819) (0.147,-0.853) (0.160,-0.887) (0.172,-0.920) (0.184,-0.954) (0.195,-0.989) (0.206,-1.023) (0.218,-1.057) (0.229,-1.091) (0.239,-1.126) (0.250,-1.160) (0.260,-1.194) (0.271,-1.229) (0.281,-1.263) (0.291,-1.298) (0.301,-1.333) (0.311,-1.367) (0.321,-1.402) (0.330,-1.437) (0.340,-1.471) (0.349,-1.506) (0.359,-1.541) (0.368,-1.576) (0.377,-1.610) (0.386,-1.645) (0.396,-1.680) (0.405,-1.715) (0.414,-1.750) (0.423,-1.785) (0.432,-1.819) (0.441,-1.854) (0.449,-1.889) (0.458,-1.924) (0.467,-1.959) (0.476,-1.994) (0.479,-2.000)}
\LC{(0.432,-0.579) (0.446,-0.612) (0.460,-0.645) (0.474,-0.678) (0.488,-0.712) (0.501,-0.745) (0.515,-0.778) (0.529,-0.812) (0.542,-0.845) (0.556,-0.878) (0.570,-0.912) (0.583,-0.945) (0.597,-0.978) (0.611,-1.011) (0.624,-1.045) (0.638,-1.078) (0.652,-1.111) (0.665,-1.145) (0.679,-1.178) (0.693,-1.211) (0.706,-1.245) (0.720,-1.278) (0.734,-1.311) (0.748,-1.344) (0.761,-1.378) (0.775,-1.411) (0.789,-1.444) (0.803,-1.477) (0.817,-1.511) (0.831,-1.544) (0.845,-1.577) (0.858,-1.610) (0.872,-1.643) (0.886,-1.677) (0.900,-1.710) (0.914,-1.743) (0.928,-1.776) (0.942,-1.809) (0.956,-1.842) (0.970,-1.876) (0.984,-1.909) (0.998,-1.942) (1.013,-1.975) (1.027,-2.000)}
\LC{(0.775,-0.437) (0.789,-0.470) (0.804,-0.503) (0.818,-0.536) (0.833,-0.569) (0.848,-0.601) (0.863,-0.634) (0.878,-0.667) (0.894,-0.699) (0.909,-0.732) (0.925,-0.764) (0.941,-0.796) (0.957,-0.828) (0.974,-0.860) (0.990,-0.892) (1.007,-0.924) (1.024,-0.956) (1.041,-0.988) (1.058,-1.019) (1.075,-1.051) (1.092,-1.083) (1.110,-1.114) (1.128,-1.145) (1.146,-1.177) (1.164,-1.208) (1.182,-1.239) (1.200,-1.270) (1.218,-1.301) (1.236,-1.332) (1.255,-1.363) (1.273,-1.394) (1.292,-1.425) (1.311,-1.455) (1.330,-1.486) (1.349,-1.517) (1.368,-1.547) (1.387,-1.578) (1.406,-1.608) (1.425,-1.639) (1.444,-1.669) (1.464,-1.699) (1.483,-1.730) (1.503,-1.760) (1.522,-1.790) (1.542,-1.820) (1.561,-1.851) (1.581,-1.881) (1.601,-1.911) (1.621,-1.941) (1.640,-1.971) (1.660,-2.000)}
\LC{(1.098,-0.303) (1.113,-0.336) (1.128,-0.368) (1.143,-0.401) (1.159,-0.433) (1.176,-0.465) (1.192,-0.497) (1.210,-0.529) (1.228,-0.560) (1.246,-0.591) (1.264,-0.622) (1.283,-0.652) (1.303,-0.683) (1.322,-0.713) (1.342,-0.743) (1.363,-0.773) (1.383,-0.802) (1.404,-0.831) (1.425,-0.860) (1.447,-0.889) (1.468,-0.918) (1.490,-0.947) (1.512,-0.975) (1.535,-1.003) (1.557,-1.031) (1.580,-1.059) (1.603,-1.087) (1.626,-1.114) (1.650,-1.142) (1.673,-1.169) (1.697,-1.196) (1.720,-1.224) (1.744,-1.250) (1.768,-1.277) (1.792,-1.304) (1.817,-1.331) (1.841,-1.357) (1.865,-1.384) (1.890,-1.410) (1.915,-1.436) (1.939,-1.462) (1.964,-1.488) (1.989,-1.514) (2.014,-1.540) (2.039,-1.566) (2.065,-1.592) (2.090,-1.617) (2.115,-1.643) (2.141,-1.668) (2.166,-1.694) (2.192,-1.719) (2.217,-1.744) (2.243,-1.770) (2.269,-1.795) (2.294,-1.820) (2.320,-1.845) (2.346,-1.870) (2.372,-1.895) (2.398,-1.920) (2.424,-1.945) (2.450,-1.970) (2.476,-1.994) (2.485,-2.000)}
\LC{(1.378,-0.187) (1.393,-0.219) (1.410,-0.251) (1.427,-0.283) (1.445,-0.314) (1.465,-0.344) (1.485,-0.374) (1.505,-0.404) (1.527,-0.432) (1.549,-0.461) (1.572,-0.488) (1.596,-0.516) (1.620,-0.543) (1.644,-0.569) (1.669,-0.595) (1.694,-0.620) (1.720,-0.646) (1.746,-0.670) (1.773,-0.695) (1.800,-0.719) (1.827,-0.743) (1.854,-0.766) (1.881,-0.789) (1.909,-0.812) (1.937,-0.835) (1.965,-0.858) (1.993,-0.880) (2.022,-0.902) (2.050,-0.924) (2.079,-0.945) (2.108,-0.967) (2.137,-0.988) (2.166,-1.009) (2.195,-1.030) (2.225,-1.051) (2.254,-1.072) (2.284,-1.092) (2.313,-1.113) (2.343,-1.133) (2.373,-1.153) (2.403,-1.173) (2.433,-1.193) (2.463,-1.213) (2.493,-1.233) (2.523,-1.252) (2.554,-1.272) (2.584,-1.291) (2.600,-1.304)}
\LC{(1.600,-0.095) (1.617,-0.126) (1.637,-0.156) (1.660,-0.184) (1.684,-0.211) (1.709,-0.237) (1.736,-0.261) (1.764,-0.284) (1.792,-0.306) (1.821,-0.327) (1.851,-0.347) (1.881,-0.366) (1.912,-0.385) (1.943,-0.403) (1.975,-0.421) (2.006,-0.438) (2.038,-0.455) (2.070,-0.471) (2.102,-0.487) (2.135,-0.503) (2.167,-0.519) (2.200,-0.534) (2.233,-0.549) (2.265,-0.563) (2.298,-0.578) (2.331,-0.592) (2.365,-0.606) (2.398,-0.620) (2.431,-0.634) (2.464,-0.647) (2.498,-0.661) (2.531,-0.674) (2.565,-0.687) (2.598,-0.700) (2.600,-0.704)}
\LC{(1.739,-0.037) (1.769,-0.057) (1.801,-0.073) (1.835,-0.085) (1.869,-0.096) (1.904,-0.106) (1.939,-0.115) (1.974,-0.123) (2.009,-0.131) (2.044,-0.138) (2.080,-0.145) (2.115,-0.152) (2.150,-0.159) (2.186,-0.165) (2.221,-0.171) (2.257,-0.177) (2.292,-0.182) (2.328,-0.188) (2.364,-0.193) (2.399,-0.199) (2.435,-0.204) (2.470,-0.209) (2.506,-0.214) (2.542,-0.219) (2.577,-0.224) (2.600,-0.227)}
\LC{(1.657,0.071) (1.677,0.101) (1.700,0.128) (1.726,0.154) (1.754,0.177) (1.782,0.198) (1.812,0.218) (1.843,0.237) (1.875,0.254) (1.907,0.271) (1.939,0.287) (1.971,0.302) (2.004,0.317) (2.037,0.331) (2.071,0.345) (2.104,0.358) (2.138,0.371) (2.171,0.384) (2.205,0.397) (2.239,0.409) (2.273,0.421) (2.307,0.432) (2.341,0.444) (2.375,0.455) (2.409,0.466) (2.444,0.477) (2.478,0.488) (2.512,0.499) (2.547,0.509) (2.581,0.520) (2.600,0.527)}
\LC{(1.466,0.150) (1.482,0.182) (1.499,0.214) (1.518,0.245) (1.538,0.275) (1.559,0.304) (1.580,0.333) (1.603,0.361) (1.627,0.388) (1.651,0.414) (1.677,0.440) (1.702,0.465) (1.729,0.490) (1.755,0.514) (1.783,0.537) (1.810,0.561) (1.838,0.583) (1.866,0.605) (1.895,0.627) (1.924,0.649) (1.953,0.670) (1.982,0.691) (2.012,0.711) (2.042,0.732) (2.072,0.752) (2.102,0.772) (2.132,0.791) (2.162,0.811) (2.193,0.830) (2.223,0.849) (2.254,0.867) (2.285,0.886) (2.316,0.905) (2.347,0.923) (2.378,0.941) (2.409,0.959) (2.440,0.977) (2.472,0.995) (2.503,1.012) (2.534,1.030) (2.566,1.047) (2.598,1.064) (2.600,1.070)}
\LC{(1.206,0.258) (1.221,0.291) (1.237,0.323) (1.253,0.355) (1.270,0.387) (1.287,0.419) (1.305,0.450) (1.323,0.481) (1.342,0.512) (1.361,0.542) (1.381,0.572) (1.402,0.602) (1.422,0.631) (1.444,0.660) (1.465,0.689) (1.487,0.717) (1.510,0.746) (1.532,0.774) (1.555,0.801) (1.578,0.829) (1.602,0.856) (1.626,0.883) (1.650,0.910) (1.674,0.937) (1.698,0.963) (1.723,0.989) (1.748,1.015) (1.773,1.041) (1.798,1.067) (1.824,1.092) (1.849,1.118) (1.875,1.143) (1.901,1.168) (1.927,1.193) (1.953,1.218) (1.979,1.242) (2.005,1.267) (2.032,1.291) (2.058,1.316) (2.085,1.340) (2.112,1.364) (2.138,1.388) (2.165,1.412) (2.192,1.436) (2.219,1.460) (2.247,1.483) (2.274,1.507) (2.301,1.530) (2.328,1.554) (2.356,1.577) (2.383,1.600) (2.411,1.623) (2.439,1.646) (2.466,1.669) (2.494,1.692) (2.522,1.715) (2.550,1.738) (2.578,1.761) (2.600,1.783)}
\LC{(0.902,0.384) (0.916,0.417) (0.930,0.450) (0.945,0.483) (0.960,0.516) (0.975,0.548) (0.991,0.581) (1.007,0.613) (1.023,0.645) (1.039,0.677) (1.056,0.709) (1.073,0.741) (1.090,0.773) (1.108,0.804) (1.126,0.835) (1.144,0.867) (1.162,0.898) (1.180,0.929) (1.199,0.959) (1.217,0.990) (1.236,1.021) (1.255,1.051) (1.275,1.082) (1.294,1.112) (1.314,1.142) (1.334,1.172) (1.353,1.202) (1.373,1.232) (1.394,1.262) (1.414,1.292) (1.434,1.321) (1.455,1.351) (1.476,1.380) (1.496,1.410) (1.517,1.439) (1.538,1.468) (1.559,1.498) (1.580,1.527) (1.602,1.556) (1.623,1.585) (1.644,1.614) (1.666,1.643) (1.687,1.672) (1.709,1.700) (1.731,1.729) (1.752,1.758) (1.774,1.786) (1.796,1.815) (1.818,1.844) (1.840,1.872) (1.862,1.901) (1.884,1.929) (1.906,1.957) (1.928,1.986) (1.943,2.000)}
\LC{(0.565,0.524) (0.579,0.557) (0.593,0.590) (0.607,0.623) (0.621,0.656) (0.635,0.689) (0.650,0.722) (0.664,0.756) (0.678,0.789) (0.693,0.822) (0.707,0.854) (0.722,0.887) (0.736,0.920) (0.751,0.953) (0.766,0.986) (0.781,1.019) (0.796,1.052) (0.811,1.084) (0.826,1.117) (0.841,1.150) (0.856,1.182) (0.871,1.215) (0.886,1.248) (0.901,1.280) (0.917,1.313) (0.932,1.345) (0.948,1.378) (0.963,1.410) (0.979,1.443) (0.995,1.475) (1.010,1.508) (1.026,1.540) (1.042,1.572) (1.058,1.605) (1.073,1.637) (1.089,1.669) (1.105,1.702) (1.121,1.734) (1.137,1.766) (1.153,1.798) (1.169,1.830) (1.186,1.863) (1.202,1.895) (1.218,1.927) (1.234,1.959) (1.250,1.991) (1.256,2.000)}
\LC{(0.223,0.666) (0.236,0.699) (0.249,0.733) (0.262,0.766) (0.275,0.800) (0.288,0.833) (0.301,0.867) (0.314,0.901) (0.326,0.934) (0.338,0.968) (0.351,1.002) (0.363,1.036) (0.375,1.070) (0.387,1.104) (0.399,1.138) (0.410,1.172) (0.422,1.206) (0.433,1.240) (0.445,1.274) (0.456,1.308) (0.468,1.343) (0.479,1.377) (0.490,1.411) (0.501,1.445) (0.513,1.479) (0.524,1.514) (0.535,1.548) (0.546,1.582) (0.557,1.616) (0.568,1.651) (0.579,1.685) (0.590,1.719) (0.601,1.754) (0.611,1.788) (0.622,1.822) (0.633,1.857) (0.644,1.891) (0.655,1.925) (0.665,1.960) (0.676,1.994) (0.680,2.000)}
\LC{(-0.107,0.803) (-0.095,0.836) (-0.082,0.870) (-0.071,0.904) (-0.060,0.938) (-0.049,0.973) (-0.039,1.007) (-0.029,1.042) (-0.019,1.077) (-0.010,1.111) (-0.002,1.146) (0.007,1.181) (0.015,1.216) (0.023,1.252) (0.031,1.287) (0.038,1.322) (0.045,1.357) (0.052,1.393) (0.059,1.428) (0.066,1.463) (0.073,1.499) (0.079,1.534) (0.085,1.569) (0.091,1.605) (0.098,1.640) (0.104,1.676) (0.109,1.711) (0.115,1.747) (0.121,1.783) (0.126,1.818) (0.132,1.854) (0.137,1.889) (0.143,1.925) (0.148,1.960) (0.153,1.996) (0.155,2.000)}

      % Conductor con punta - flechas
% Generado por scripts/generar_datos.py (no editar a mano).
\FL{-0.332}{1.586}{89.3}
\FL{-0.742}{1.606}{100.9}
\FL{-1.161}{1.589}{113.5}
\FL{-1.635}{1.552}{126.9}
\FL{-2.093}{1.363}{140.4}
\FL{-2.144}{0.864}{152.5}
\FL{-2.187}{0.452}{165.0}
\FL{-2.200}{0.071}{177.6}
\FL{-2.197}{-0.307}{-169.8}
\FL{-2.160}{-0.700}{-157.2}
\FL{-2.111}{-1.155}{-144.9}
\FL{-1.855}{-1.538}{-132.3}
\FL{-1.337}{-1.582}{-118.7}
\FL{-0.898}{-1.597}{-105.6}
\FL{-0.488}{-1.595}{-93.5}
\FL{-0.070}{-1.563}{-83.0}
\FL{0.352}{-1.518}{-74.8}
\FL{0.798}{-1.466}{-67.4}
\FL{1.292}{-1.425}{-58.7}
\FL{1.890}{-1.410}{-46.8}
\FL{2.079}{-0.945}{-36.8}
\FL{2.189}{-0.529}{-24.8}
\FL{2.269}{-0.179}{-9.2}
\FL{2.216}{0.401}{19.8}
\FL{2.122}{0.785}{32.9}
\FL{2.005}{1.267}{42.9}
\FL{1.503}{1.420}{54.6}
\FL{0.984}{1.454}{64.3}
\FL{0.516}{1.491}{72.0}
\FL{0.081}{1.546}{79.9}

      % Contorno de la gota
% Generado por scripts/generar_datos.py (no editar a mano).
\def\formagota{(-0.259,0.833) -- (-0.262,0.834) -- (-0.265,0.835) -- (-0.268,0.836) -- (-0.270,0.837) -- (-0.273,0.839) -- (-0.276,0.840) -- (-0.279,0.841) -- (-0.281,0.842) -- (-0.284,0.843) -- (-0.287,0.844) -- (-0.290,0.845) -- (-0.293,0.846) -- (-0.295,0.847) -- (-0.298,0.848) -- (-0.301,0.849) -- (-0.304,0.850) -- (-0.306,0.851) -- (-0.309,0.852) -- (-0.312,0.853) -- (-0.315,0.854) -- (-0.318,0.855) -- (-0.320,0.855) -- (-0.323,0.856) -- (-0.326,0.857) -- (-0.329,0.858) -- (-0.332,0.859) -- (-0.335,0.860) -- (-0.337,0.861) -- (-0.340,0.862) -- (-0.343,0.863) -- (-0.346,0.863) -- (-0.349,0.864) -- (-0.352,0.865) -- (-0.354,0.866) -- (-0.357,0.867) -- (-0.360,0.867) -- (-0.363,0.868) -- (-0.366,0.869) -- (-0.369,0.870) -- (-0.371,0.870) -- (-0.374,0.871) -- (-0.377,0.872) -- (-0.380,0.873) -- (-0.383,0.873) -- (-0.386,0.874) -- (-0.389,0.875) -- (-0.392,0.876) -- (-0.394,0.876) -- (-0.397,0.877) -- (-0.400,0.878) -- (-0.403,0.878) -- (-0.406,0.879) -- (-0.409,0.879) -- (-0.412,0.880) -- (-0.415,0.881) -- (-0.418,0.881) -- (-0.420,0.882) -- (-0.423,0.882) -- (-0.426,0.883) -- (-0.429,0.884) -- (-0.432,0.884) -- (-0.435,0.885) -- (-0.438,0.885) -- (-0.441,0.886) -- (-0.444,0.886) -- (-0.447,0.887) -- (-0.449,0.887) -- (-0.452,0.888) -- (-0.455,0.888) -- (-0.458,0.889) -- (-0.461,0.889) -- (-0.464,0.890) -- (-0.467,0.890) -- (-0.470,0.891) -- (-0.473,0.891) -- (-0.476,0.891) -- (-0.479,0.892) -- (-0.482,0.892) -- (-0.485,0.893) -- (-0.487,0.893) -- (-0.490,0.893) -- (-0.493,0.894) -- (-0.496,0.894) -- (-0.499,0.894) -- (-0.502,0.895) -- (-0.505,0.895) -- (-0.508,0.895) -- (-0.511,0.896) -- (-0.514,0.896) -- (-0.517,0.896) -- (-0.520,0.896) -- (-0.523,0.897) -- (-0.526,0.897) -- (-0.529,0.897) -- (-0.532,0.897) -- (-0.535,0.898) -- (-0.537,0.898) -- (-0.540,0.898) -- (-0.543,0.898) -- (-0.546,0.898) -- (-0.549,0.899) -- (-0.552,0.899) -- (-0.555,0.899) -- (-0.558,0.899) -- (-0.561,0.899) -- (-0.564,0.899) -- (-0.567,0.899) -- (-0.570,0.899) -- (-0.573,0.900) -- (-0.576,0.900) -- (-0.579,0.900) -- (-0.582,0.900) -- (-0.585,0.900) -- (-0.588,0.900) -- (-0.591,0.900) -- (-0.594,0.900) -- (-0.597,0.900) -- (-0.599,0.900) -- (-0.602,0.900) -- (-0.605,0.900) -- (-0.608,0.900) -- (-0.611,0.900) -- (-0.614,0.900) -- (-0.617,0.900) -- (-0.620,0.900) -- (-0.623,0.900) -- (-0.626,0.900) -- (-0.629,0.900) -- (-0.632,0.899) -- (-0.635,0.899) -- (-0.638,0.899) -- (-0.641,0.899) -- (-0.644,0.899) -- (-0.647,0.899) -- (-0.650,0.899) -- (-0.653,0.898) -- (-0.656,0.898) -- (-0.659,0.898) -- (-0.661,0.898) -- (-0.664,0.898) -- (-0.667,0.897) -- (-0.670,0.897) -- (-0.673,0.897) -- (-0.676,0.897) -- (-0.679,0.896) -- (-0.682,0.896) -- (-0.685,0.896) -- (-0.688,0.896) -- (-0.691,0.895) -- (-0.694,0.895) -- (-0.697,0.895) -- (-0.700,0.894) -- (-0.703,0.894) -- (-0.706,0.894) -- (-0.709,0.893) -- (-0.711,0.893) -- (-0.714,0.893) -- (-0.717,0.892) -- (-0.720,0.892) -- (-0.723,0.892) -- (-0.726,0.891) -- (-0.729,0.891) -- (-0.732,0.890) -- (-0.735,0.890) -- (-0.738,0.889) -- (-0.741,0.889) -- (-0.744,0.888) -- (-0.747,0.888) -- (-0.750,0.887) -- (-0.752,0.887) -- (-0.755,0.886) -- (-0.758,0.886) -- (-0.761,0.885) -- (-0.764,0.885) -- (-0.767,0.884) -- (-0.770,0.884) -- (-0.773,0.883) -- (-0.776,0.883) -- (-0.779,0.882) -- (-0.781,0.881) -- (-0.784,0.881) -- (-0.787,0.880) -- (-0.790,0.880) -- (-0.793,0.879) -- (-0.796,0.878) -- (-0.799,0.878) -- (-0.802,0.877) -- (-0.805,0.876) -- (-0.807,0.876) -- (-0.810,0.875) -- (-0.813,0.874) -- (-0.816,0.874) -- (-0.819,0.873) -- (-0.822,0.872) -- (-0.825,0.871) -- (-0.828,0.871) -- (-0.830,0.870) -- (-0.833,0.869) -- (-0.836,0.868) -- (-0.839,0.868) -- (-0.842,0.867) -- (-0.845,0.866) -- (-0.847,0.865) -- (-0.850,0.864) -- (-0.853,0.864) -- (-0.856,0.863) -- (-0.859,0.862) -- (-0.862,0.861) -- (-0.864,0.860) -- (-0.867,0.859) -- (-0.870,0.858) -- (-0.873,0.858) -- (-0.876,0.857) -- (-0.879,0.856) -- (-0.881,0.855) -- (-0.884,0.854) -- (-0.887,0.853) -- (-0.890,0.852) -- (-0.893,0.851) -- (-0.895,0.850) -- (-0.898,0.849) -- (-0.901,0.848) -- (-0.904,0.847) -- (-0.907,0.846) -- (-0.909,0.845) -- (-0.912,0.844) -- (-0.915,0.843) -- (-0.918,0.842) -- (-0.920,0.841) -- (-0.923,0.840) -- (-0.926,0.839) -- (-0.929,0.838) -- (-0.931,0.837) -- (-0.934,0.836) -- (-0.937,0.835) -- (-0.940,0.833) -- (-0.942,0.832) -- (-0.945,0.831) -- (-0.948,0.830) -- (-0.951,0.829) -- (-0.953,0.828) -- (-0.956,0.827) -- (-0.959,0.825) -- (-0.961,0.824) -- (-0.964,0.823) -- (-0.967,0.822) -- (-0.969,0.821) -- (-0.972,0.819) -- (-0.975,0.818) -- (-0.978,0.817) -- (-0.980,0.816) -- (-0.983,0.814) -- (-0.986,0.813) -- (-0.988,0.812) -- (-0.991,0.811) -- (-0.994,0.809) -- (-0.996,0.808) -- (-0.999,0.807) -- (-1.002,0.805) -- (-1.004,0.804) -- (-1.007,0.803) -- (-1.009,0.801) -- (-1.012,0.800) -- (-1.015,0.799) -- (-1.017,0.797) -- (-1.020,0.796) -- (-1.023,0.795) -- (-1.025,0.793) -- (-1.028,0.792) -- (-1.030,0.790) -- (-1.033,0.789) -- (-1.036,0.788) -- (-1.038,0.786) -- (-1.041,0.785) -- (-1.043,0.783) -- (-1.046,0.782) -- (-1.048,0.780) -- (-1.051,0.779) -- (-1.054,0.777) -- (-1.056,0.776) -- (-1.059,0.774) -- (-1.061,0.773) -- (-1.064,0.771) -- (-1.066,0.770) -- (-1.069,0.768) -- (-1.071,0.767) -- (-1.074,0.765) -- (-1.076,0.764) -- (-1.079,0.762) -- (-1.081,0.760) -- (-1.084,0.759) -- (-1.086,0.757) -- (-1.089,0.756) -- (-1.091,0.754) -- (-1.094,0.752) -- (-1.096,0.751) -- (-1.099,0.749) -- (-1.101,0.748) -- (-1.104,0.746) -- (-1.106,0.744) -- (-1.108,0.743) -- (-1.111,0.741) -- (-1.113,0.739) -- (-1.116,0.738) -- (-1.118,0.736) -- (-1.121,0.734) -- (-1.123,0.732) -- (-1.125,0.731) -- (-1.128,0.729) -- (-1.130,0.727) -- (-1.133,0.725) -- (-1.135,0.724) -- (-1.137,0.722) -- (-1.140,0.720) -- (-1.142,0.718) -- (-1.144,0.717) -- (-1.147,0.715) -- (-1.149,0.713) -- (-1.151,0.711) -- (-1.154,0.709) -- (-1.156,0.708) -- (-1.158,0.706) -- (-1.161,0.704) -- (-1.163,0.702) -- (-1.165,0.700) -- (-1.168,0.698) -- (-1.170,0.697) -- (-1.172,0.695) -- (-1.174,0.693) -- (-1.177,0.691) -- (-1.179,0.689) -- (-1.181,0.687) -- (-1.184,0.685) -- (-1.186,0.683) -- (-1.188,0.681) -- (-1.190,0.679) -- (-1.192,0.677) -- (-1.195,0.675) -- (-1.197,0.674) -- (-1.199,0.672) -- (-1.201,0.670) -- (-1.204,0.668) -- (-1.206,0.666) -- (-1.208,0.664) -- (-1.210,0.662) -- (-1.212,0.660) -- (-1.214,0.658) -- (-1.217,0.656) -- (-1.219,0.654) -- (-1.221,0.652) -- (-1.223,0.649) -- (-1.225,0.647) -- (-1.227,0.645) -- (-1.229,0.643) -- (-1.231,0.641) -- (-1.234,0.639) -- (-1.236,0.637) -- (-1.238,0.635) -- (-1.240,0.633) -- (-1.242,0.631) -- (-1.244,0.629) -- (-1.246,0.627) -- (-1.248,0.624) -- (-1.250,0.622) -- (-1.252,0.620) -- (-1.254,0.618) -- (-1.256,0.616) -- (-1.258,0.614) -- (-1.260,0.612) -- (-1.262,0.609) -- (-1.264,0.607) -- (-1.266,0.605) -- (-1.268,0.603) -- (-1.270,0.601) -- (-1.272,0.598) -- (-1.274,0.596) -- (-1.276,0.594) -- (-1.278,0.592) -- (-1.280,0.590) -- (-1.282,0.587) -- (-1.284,0.585) -- (-1.286,0.583) -- (-1.288,0.581) -- (-1.290,0.578) -- (-1.291,0.576) -- (-1.293,0.574) -- (-1.295,0.571) -- (-1.297,0.569) -- (-1.299,0.567) -- (-1.301,0.565) -- (-1.303,0.562) -- (-1.305,0.560) -- (-1.306,0.558) -- (-1.308,0.555) -- (-1.310,0.553) -- (-1.312,0.551) -- (-1.314,0.548) -- (-1.315,0.546) -- (-1.317,0.544) -- (-1.319,0.541) -- (-1.321,0.539) -- (-1.323,0.537) -- (-1.324,0.534) -- (-1.326,0.532) -- (-1.328,0.529) -- (-1.330,0.527) -- (-1.331,0.525) -- (-1.333,0.522) -- (-1.335,0.520) -- (-1.336,0.517) -- (-1.338,0.515) -- (-1.340,0.513) -- (-1.341,0.510) -- (-1.343,0.508) -- (-1.345,0.505) -- (-1.346,0.503) -- (-1.348,0.500) -- (-1.350,0.498) -- (-1.351,0.495) -- (-1.353,0.493) -- (-1.355,0.490) -- (-1.356,0.488) -- (-1.358,0.485) -- (-1.359,0.483) -- (-1.361,0.480) -- (-1.363,0.478) -- (-1.364,0.475) -- (-1.366,0.473) -- (-1.367,0.470) -- (-1.369,0.468) -- (-1.370,0.465) -- (-1.372,0.463) -- (-1.373,0.460) -- (-1.375,0.458) -- (-1.376,0.455) -- (-1.378,0.453) -- (-1.379,0.450) -- (-1.381,0.448) -- (-1.382,0.445) -- (-1.384,0.442) -- (-1.385,0.440) -- (-1.387,0.437) -- (-1.388,0.435) -- (-1.389,0.432) -- (-1.391,0.430) -- (-1.392,0.427) -- (-1.394,0.424) -- (-1.395,0.422) -- (-1.396,0.419) -- (-1.398,0.416) -- (-1.399,0.414) -- (-1.401,0.411) -- (-1.402,0.409) -- (-1.403,0.406) -- (-1.405,0.403) -- (-1.406,0.401) -- (-1.407,0.398) -- (-1.408,0.395) -- (-1.410,0.393) -- (-1.411,0.390) -- (-1.412,0.387) -- (-1.414,0.385) -- (-1.415,0.382) -- (-1.416,0.379) -- (-1.417,0.377) -- (-1.419,0.374) -- (-1.420,0.371) -- (-1.421,0.369) -- (-1.422,0.366) -- (-1.423,0.363) -- (-1.425,0.361) -- (-1.426,0.358) -- (-1.427,0.355) -- (-1.428,0.352) -- (-1.429,0.350) -- (-1.430,0.347) -- (-1.432,0.344) -- (-1.433,0.341) -- (-1.434,0.339) -- (-1.435,0.336) -- (-1.436,0.333) -- (-1.437,0.330) -- (-1.438,0.328) -- (-1.439,0.325) -- (-1.440,0.322) -- (-1.441,0.319) -- (-1.442,0.317) -- (-1.443,0.314) -- (-1.444,0.311) -- (-1.445,0.308) -- (-1.447,0.306) -- (-1.448,0.303) -- (-1.448,0.300) -- (-1.449,0.297) -- (-1.450,0.294) -- (-1.451,0.292) -- (-1.452,0.289) -- (-1.453,0.286) -- (-1.454,0.283) -- (-1.455,0.280) -- (-1.456,0.278) -- (-1.457,0.275) -- (-1.458,0.272) -- (-1.459,0.269) -- (-1.460,0.266) -- (-1.461,0.264) -- (-1.461,0.261) -- (-1.462,0.258) -- (-1.463,0.255) -- (-1.464,0.252) -- (-1.465,0.249) -- (-1.466,0.247) -- (-1.466,0.244) -- (-1.467,0.241) -- (-1.468,0.238) -- (-1.469,0.235) -- (-1.469,0.232) -- (-1.470,0.229) -- (-1.471,0.227) -- (-1.472,0.224) -- (-1.472,0.221) -- (-1.473,0.218) -- (-1.474,0.215) -- (-1.475,0.212) -- (-1.475,0.209) -- (-1.476,0.207) -- (-1.477,0.204) -- (-1.477,0.201) -- (-1.478,0.198) -- (-1.479,0.195) -- (-1.479,0.192) -- (-1.480,0.189) -- (-1.480,0.186) -- (-1.481,0.183) -- (-1.482,0.181) -- (-1.482,0.178) -- (-1.483,0.175) -- (-1.483,0.172) -- (-1.484,0.169) -- (-1.485,0.166) -- (-1.485,0.163) -- (-1.486,0.160) -- (-1.486,0.157) -- (-1.487,0.154) -- (-1.487,0.151) -- (-1.488,0.149) -- (-1.488,0.146) -- (-1.489,0.143) -- (-1.489,0.140) -- (-1.490,0.137) -- (-1.490,0.134) -- (-1.490,0.131) -- (-1.491,0.128) -- (-1.491,0.125) -- (-1.492,0.122) -- (-1.492,0.119) -- (-1.492,0.116) -- (-1.493,0.113) -- (-1.493,0.111) -- (-1.494,0.108) -- (-1.494,0.105) -- (-1.494,0.102) -- (-1.495,0.099) -- (-1.495,0.096) -- (-1.495,0.093) -- (-1.495,0.090) -- (-1.496,0.087) -- (-1.496,0.084) -- (-1.496,0.081) -- (-1.497,0.078) -- (-1.497,0.075) -- (-1.497,0.072) -- (-1.497,0.069) -- (-1.498,0.066) -- (-1.498,0.063) -- (-1.498,0.061) -- (-1.498,0.058) -- (-1.498,0.055) -- (-1.498,0.052) -- (-1.499,0.049) -- (-1.499,0.046) -- (-1.499,0.043) -- (-1.499,0.040) -- (-1.499,0.037) -- (-1.499,0.034) -- (-1.499,0.031) -- (-1.500,0.028) -- (-1.500,0.025) -- (-1.500,0.022) -- (-1.500,0.019) -- (-1.500,0.016) -- (-1.500,0.013) -- (-1.500,0.010) -- (-1.500,0.007) -- (-1.500,0.004) -- (-1.500,0.001) -- (-1.500,-0.001) -- (-1.500,-0.004) -- (-1.500,-0.007) -- (-1.500,-0.010) -- (-1.500,-0.013) -- (-1.500,-0.016) -- (-1.500,-0.019) -- (-1.500,-0.022) -- (-1.500,-0.025) -- (-1.500,-0.028) -- (-1.499,-0.031) -- (-1.499,-0.034) -- (-1.499,-0.037) -- (-1.499,-0.040) -- (-1.499,-0.043) -- (-1.499,-0.046) -- (-1.499,-0.049) -- (-1.498,-0.052) -- (-1.498,-0.055) -- (-1.498,-0.058) -- (-1.498,-0.061) -- (-1.498,-0.063) -- (-1.498,-0.066) -- (-1.497,-0.069) -- (-1.497,-0.072) -- (-1.497,-0.075) -- (-1.497,-0.078) -- (-1.496,-0.081) -- (-1.496,-0.084) -- (-1.496,-0.087) -- (-1.495,-0.090) -- (-1.495,-0.093) -- (-1.495,-0.096) -- (-1.495,-0.099) -- (-1.494,-0.102) -- (-1.494,-0.105) -- (-1.494,-0.108) -- (-1.493,-0.111) -- (-1.493,-0.113) -- (-1.492,-0.116) -- (-1.492,-0.119) -- (-1.492,-0.122) -- (-1.491,-0.125) -- (-1.491,-0.128) -- (-1.490,-0.131) -- (-1.490,-0.134) -- (-1.490,-0.137) -- (-1.489,-0.140) -- (-1.489,-0.143) -- (-1.488,-0.146) -- (-1.488,-0.149) -- (-1.487,-0.151) -- (-1.487,-0.154) -- (-1.486,-0.157) -- (-1.486,-0.160) -- (-1.485,-0.163) -- (-1.485,-0.166) -- (-1.484,-0.169) -- (-1.483,-0.172) -- (-1.483,-0.175) -- (-1.482,-0.178) -- (-1.482,-0.181) -- (-1.481,-0.183) -- (-1.480,-0.186) -- (-1.480,-0.189) -- (-1.479,-0.192) -- (-1.479,-0.195) -- (-1.478,-0.198) -- (-1.477,-0.201) -- (-1.477,-0.204) -- (-1.476,-0.207) -- (-1.475,-0.209) -- (-1.475,-0.212) -- (-1.474,-0.215) -- (-1.473,-0.218) -- (-1.472,-0.221) -- (-1.472,-0.224) -- (-1.471,-0.227) -- (-1.470,-0.229) -- (-1.469,-0.232) -- (-1.469,-0.235) -- (-1.468,-0.238) -- (-1.467,-0.241) -- (-1.466,-0.244) -- (-1.466,-0.247) -- (-1.465,-0.249) -- (-1.464,-0.252) -- (-1.463,-0.255) -- (-1.462,-0.258) -- (-1.461,-0.261) -- (-1.461,-0.264) -- (-1.460,-0.266) -- (-1.459,-0.269) -- (-1.458,-0.272) -- (-1.457,-0.275) -- (-1.456,-0.278) -- (-1.455,-0.280) -- (-1.454,-0.283) -- (-1.453,-0.286) -- (-1.452,-0.289) -- (-1.451,-0.292) -- (-1.450,-0.294) -- (-1.449,-0.297) -- (-1.448,-0.300) -- (-1.448,-0.303) -- (-1.447,-0.306) -- (-1.445,-0.308) -- (-1.444,-0.311) -- (-1.443,-0.314) -- (-1.442,-0.317) -- (-1.441,-0.319) -- (-1.440,-0.322) -- (-1.439,-0.325) -- (-1.438,-0.328) -- (-1.437,-0.330) -- (-1.436,-0.333) -- (-1.435,-0.336) -- (-1.434,-0.339) -- (-1.433,-0.341) -- (-1.432,-0.344) -- (-1.430,-0.347) -- (-1.429,-0.350) -- (-1.428,-0.352) -- (-1.427,-0.355) -- (-1.426,-0.358) -- (-1.425,-0.361) -- (-1.423,-0.363) -- (-1.422,-0.366) -- (-1.421,-0.369) -- (-1.420,-0.371) -- (-1.419,-0.374) -- (-1.417,-0.377) -- (-1.416,-0.379) -- (-1.415,-0.382) -- (-1.414,-0.385) -- (-1.412,-0.387) -- (-1.411,-0.390) -- (-1.410,-0.393) -- (-1.408,-0.395) -- (-1.407,-0.398) -- (-1.406,-0.401) -- (-1.405,-0.403) -- (-1.403,-0.406) -- (-1.402,-0.409) -- (-1.401,-0.411) -- (-1.399,-0.414) -- (-1.398,-0.416) -- (-1.396,-0.419) -- (-1.395,-0.422) -- (-1.394,-0.424) -- (-1.392,-0.427) -- (-1.391,-0.430) -- (-1.389,-0.432) -- (-1.388,-0.435) -- (-1.387,-0.437) -- (-1.385,-0.440) -- (-1.384,-0.442) -- (-1.382,-0.445) -- (-1.381,-0.448) -- (-1.379,-0.450) -- (-1.378,-0.453) -- (-1.376,-0.455) -- (-1.375,-0.458) -- (-1.373,-0.460) -- (-1.372,-0.463) -- (-1.370,-0.465) -- (-1.369,-0.468) -- (-1.367,-0.470) -- (-1.366,-0.473) -- (-1.364,-0.475) -- (-1.363,-0.478) -- (-1.361,-0.480) -- (-1.359,-0.483) -- (-1.358,-0.485) -- (-1.356,-0.488) -- (-1.355,-0.490) -- (-1.353,-0.493) -- (-1.351,-0.495) -- (-1.350,-0.498) -- (-1.348,-0.500) -- (-1.346,-0.503) -- (-1.345,-0.505) -- (-1.343,-0.508) -- (-1.341,-0.510) -- (-1.340,-0.513) -- (-1.338,-0.515) -- (-1.336,-0.517) -- (-1.335,-0.520) -- (-1.333,-0.522) -- (-1.331,-0.525) -- (-1.330,-0.527) -- (-1.328,-0.529) -- (-1.326,-0.532) -- (-1.324,-0.534) -- (-1.323,-0.537) -- (-1.321,-0.539) -- (-1.319,-0.541) -- (-1.317,-0.544) -- (-1.315,-0.546) -- (-1.314,-0.548) -- (-1.312,-0.551) -- (-1.310,-0.553) -- (-1.308,-0.555) -- (-1.306,-0.558) -- (-1.305,-0.560) -- (-1.303,-0.562) -- (-1.301,-0.565) -- (-1.299,-0.567) -- (-1.297,-0.569) -- (-1.295,-0.571) -- (-1.293,-0.574) -- (-1.291,-0.576) -- (-1.290,-0.578) -- (-1.288,-0.581) -- (-1.286,-0.583) -- (-1.284,-0.585) -- (-1.282,-0.587) -- (-1.280,-0.590) -- (-1.278,-0.592) -- (-1.276,-0.594) -- (-1.274,-0.596) -- (-1.272,-0.598) -- (-1.270,-0.601) -- (-1.268,-0.603) -- (-1.266,-0.605) -- (-1.264,-0.607) -- (-1.262,-0.609) -- (-1.260,-0.612) -- (-1.258,-0.614) -- (-1.256,-0.616) -- (-1.254,-0.618) -- (-1.252,-0.620) -- (-1.250,-0.622) -- (-1.248,-0.624) -- (-1.246,-0.627) -- (-1.244,-0.629) -- (-1.242,-0.631) -- (-1.240,-0.633) -- (-1.238,-0.635) -- (-1.236,-0.637) -- (-1.234,-0.639) -- (-1.231,-0.641) -- (-1.229,-0.643) -- (-1.227,-0.645) -- (-1.225,-0.647) -- (-1.223,-0.649) -- (-1.221,-0.652) -- (-1.219,-0.654) -- (-1.217,-0.656) -- (-1.214,-0.658) -- (-1.212,-0.660) -- (-1.210,-0.662) -- (-1.208,-0.664) -- (-1.206,-0.666) -- (-1.204,-0.668) -- (-1.201,-0.670) -- (-1.199,-0.672) -- (-1.197,-0.674) -- (-1.195,-0.675) -- (-1.192,-0.677) -- (-1.190,-0.679) -- (-1.188,-0.681) -- (-1.186,-0.683) -- (-1.184,-0.685) -- (-1.181,-0.687) -- (-1.179,-0.689) -- (-1.177,-0.691) -- (-1.174,-0.693) -- (-1.172,-0.695) -- (-1.170,-0.697) -- (-1.168,-0.698) -- (-1.165,-0.700) -- (-1.163,-0.702) -- (-1.161,-0.704) -- (-1.158,-0.706) -- (-1.156,-0.708) -- (-1.154,-0.709) -- (-1.151,-0.711) -- (-1.149,-0.713) -- (-1.147,-0.715) -- (-1.144,-0.717) -- (-1.142,-0.718) -- (-1.140,-0.720) -- (-1.137,-0.722) -- (-1.135,-0.724) -- (-1.133,-0.725) -- (-1.130,-0.727) -- (-1.128,-0.729) -- (-1.125,-0.731) -- (-1.123,-0.732) -- (-1.121,-0.734) -- (-1.118,-0.736) -- (-1.116,-0.738) -- (-1.113,-0.739) -- (-1.111,-0.741) -- (-1.108,-0.743) -- (-1.106,-0.744) -- (-1.104,-0.746) -- (-1.101,-0.748) -- (-1.099,-0.749) -- (-1.096,-0.751) -- (-1.094,-0.752) -- (-1.091,-0.754) -- (-1.089,-0.756) -- (-1.086,-0.757) -- (-1.084,-0.759) -- (-1.081,-0.760) -- (-1.079,-0.762) -- (-1.076,-0.764) -- (-1.074,-0.765) -- (-1.071,-0.767) -- (-1.069,-0.768) -- (-1.066,-0.770) -- (-1.064,-0.771) -- (-1.061,-0.773) -- (-1.059,-0.774) -- (-1.056,-0.776) -- (-1.054,-0.777) -- (-1.051,-0.779) -- (-1.048,-0.780) -- (-1.046,-0.782) -- (-1.043,-0.783) -- (-1.041,-0.785) -- (-1.038,-0.786) -- (-1.036,-0.788) -- (-1.033,-0.789) -- (-1.030,-0.790) -- (-1.028,-0.792) -- (-1.025,-0.793) -- (-1.023,-0.795) -- (-1.020,-0.796) -- (-1.017,-0.797) -- (-1.015,-0.799) -- (-1.012,-0.800) -- (-1.009,-0.801) -- (-1.007,-0.803) -- (-1.004,-0.804) -- (-1.002,-0.805) -- (-0.999,-0.807) -- (-0.996,-0.808) -- (-0.994,-0.809) -- (-0.991,-0.811) -- (-0.988,-0.812) -- (-0.986,-0.813) -- (-0.983,-0.814) -- (-0.980,-0.816) -- (-0.978,-0.817) -- (-0.975,-0.818) -- (-0.972,-0.819) -- (-0.969,-0.821) -- (-0.967,-0.822) -- (-0.964,-0.823) -- (-0.961,-0.824) -- (-0.959,-0.825) -- (-0.956,-0.827) -- (-0.953,-0.828) -- (-0.951,-0.829) -- (-0.948,-0.830) -- (-0.945,-0.831) -- (-0.942,-0.832) -- (-0.940,-0.833) -- (-0.937,-0.835) -- (-0.934,-0.836) -- (-0.931,-0.837) -- (-0.929,-0.838) -- (-0.926,-0.839) -- (-0.923,-0.840) -- (-0.920,-0.841) -- (-0.918,-0.842) -- (-0.915,-0.843) -- (-0.912,-0.844) -- (-0.909,-0.845) -- (-0.907,-0.846) -- (-0.904,-0.847) -- (-0.901,-0.848) -- (-0.898,-0.849) -- (-0.895,-0.850) -- (-0.893,-0.851) -- (-0.890,-0.852) -- (-0.887,-0.853) -- (-0.884,-0.854) -- (-0.881,-0.855) -- (-0.879,-0.856) -- (-0.876,-0.857) -- (-0.873,-0.858) -- (-0.870,-0.858) -- (-0.867,-0.859) -- (-0.864,-0.860) -- (-0.862,-0.861) -- (-0.859,-0.862) -- (-0.856,-0.863) -- (-0.853,-0.864) -- (-0.850,-0.864) -- (-0.847,-0.865) -- (-0.845,-0.866) -- (-0.842,-0.867) -- (-0.839,-0.868) -- (-0.836,-0.868) -- (-0.833,-0.869) -- (-0.830,-0.870) -- (-0.828,-0.871) -- (-0.825,-0.871) -- (-0.822,-0.872) -- (-0.819,-0.873) -- (-0.816,-0.874) -- (-0.813,-0.874) -- (-0.810,-0.875) -- (-0.807,-0.876) -- (-0.805,-0.876) -- (-0.802,-0.877) -- (-0.799,-0.878) -- (-0.796,-0.878) -- (-0.793,-0.879) -- (-0.790,-0.880) -- (-0.787,-0.880) -- (-0.784,-0.881) -- (-0.781,-0.881) -- (-0.779,-0.882) -- (-0.776,-0.883) -- (-0.773,-0.883) -- (-0.770,-0.884) -- (-0.767,-0.884) -- (-0.764,-0.885) -- (-0.761,-0.885) -- (-0.758,-0.886) -- (-0.755,-0.886) -- (-0.752,-0.887) -- (-0.750,-0.887) -- (-0.747,-0.888) -- (-0.744,-0.888) -- (-0.741,-0.889) -- (-0.738,-0.889) -- (-0.735,-0.890) -- (-0.732,-0.890) -- (-0.729,-0.891) -- (-0.726,-0.891) -- (-0.723,-0.892) -- (-0.720,-0.892) -- (-0.717,-0.892) -- (-0.714,-0.893) -- (-0.711,-0.893) -- (-0.709,-0.893) -- (-0.706,-0.894) -- (-0.703,-0.894) -- (-0.700,-0.894) -- (-0.697,-0.895) -- (-0.694,-0.895) -- (-0.691,-0.895) -- (-0.688,-0.896) -- (-0.685,-0.896) -- (-0.682,-0.896) -- (-0.679,-0.896) -- (-0.676,-0.897) -- (-0.673,-0.897) -- (-0.670,-0.897) -- (-0.667,-0.897) -- (-0.664,-0.898) -- (-0.661,-0.898) -- (-0.659,-0.898) -- (-0.656,-0.898) -- (-0.653,-0.898) -- (-0.650,-0.899) -- (-0.647,-0.899) -- (-0.644,-0.899) -- (-0.641,-0.899) -- (-0.638,-0.899) -- (-0.635,-0.899) -- (-0.632,-0.899) -- (-0.629,-0.900) -- (-0.626,-0.900) -- (-0.623,-0.900) -- (-0.620,-0.900) -- (-0.617,-0.900) -- (-0.614,-0.900) -- (-0.611,-0.900) -- (-0.608,-0.900) -- (-0.605,-0.900) -- (-0.602,-0.900) -- (-0.599,-0.900) -- (-0.597,-0.900) -- (-0.594,-0.900) -- (-0.591,-0.900) -- (-0.588,-0.900) -- (-0.585,-0.900) -- (-0.582,-0.900) -- (-0.579,-0.900) -- (-0.576,-0.900) -- (-0.573,-0.900) -- (-0.570,-0.899) -- (-0.567,-0.899) -- (-0.564,-0.899) -- (-0.561,-0.899) -- (-0.558,-0.899) -- (-0.555,-0.899) -- (-0.552,-0.899) -- (-0.549,-0.899) -- (-0.546,-0.898) -- (-0.543,-0.898) -- (-0.540,-0.898) -- (-0.537,-0.898) -- (-0.535,-0.898) -- (-0.532,-0.897) -- (-0.529,-0.897) -- (-0.526,-0.897) -- (-0.523,-0.897) -- (-0.520,-0.896) -- (-0.517,-0.896) -- (-0.514,-0.896) -- (-0.511,-0.896) -- (-0.508,-0.895) -- (-0.505,-0.895) -- (-0.502,-0.895) -- (-0.499,-0.894) -- (-0.496,-0.894) -- (-0.493,-0.894) -- (-0.490,-0.893) -- (-0.487,-0.893) -- (-0.485,-0.893) -- (-0.482,-0.892) -- (-0.479,-0.892) -- (-0.476,-0.891) -- (-0.473,-0.891) -- (-0.470,-0.891) -- (-0.467,-0.890) -- (-0.464,-0.890) -- (-0.461,-0.889) -- (-0.458,-0.889) -- (-0.455,-0.888) -- (-0.452,-0.888) -- (-0.449,-0.887) -- (-0.447,-0.887) -- (-0.444,-0.886) -- (-0.441,-0.886) -- (-0.438,-0.885) -- (-0.435,-0.885) -- (-0.432,-0.884) -- (-0.429,-0.884) -- (-0.426,-0.883) -- (-0.423,-0.882) -- (-0.420,-0.882) -- (-0.418,-0.881) -- (-0.415,-0.881) -- (-0.412,-0.880) -- (-0.409,-0.879) -- (-0.406,-0.879) -- (-0.403,-0.878) -- (-0.400,-0.878) -- (-0.397,-0.877) -- (-0.394,-0.876) -- (-0.392,-0.876) -- (-0.389,-0.875) -- (-0.386,-0.874) -- (-0.383,-0.873) -- (-0.380,-0.873) -- (-0.377,-0.872) -- (-0.374,-0.871) -- (-0.371,-0.870) -- (-0.369,-0.870) -- (-0.366,-0.869) -- (-0.363,-0.868) -- (-0.360,-0.867) -- (-0.357,-0.867) -- (-0.354,-0.866) -- (-0.352,-0.865) -- (-0.349,-0.864) -- (-0.346,-0.863) -- (-0.343,-0.863) -- (-0.340,-0.862) -- (-0.337,-0.861) -- (-0.335,-0.860) -- (-0.332,-0.859) -- (-0.329,-0.858) -- (-0.326,-0.857) -- (-0.323,-0.856) -- (-0.320,-0.855) -- (-0.318,-0.855) -- (-0.315,-0.854) -- (-0.312,-0.853) -- (-0.309,-0.852) -- (-0.306,-0.851) -- (-0.304,-0.850) -- (-0.301,-0.849) -- (-0.298,-0.848) -- (-0.295,-0.847) -- (-0.293,-0.846) -- (-0.290,-0.845) -- (-0.287,-0.844) -- (-0.284,-0.843) -- (-0.281,-0.842) -- (-0.279,-0.841) -- (-0.276,-0.840) -- (-0.273,-0.839) -- (-0.270,-0.837) -- (-0.268,-0.836) -- (-0.265,-0.835) -- (-0.262,-0.834) -- (-0.259,-0.833) -- (-0.256,-0.832) -- (-0.251,-0.830) -- (-0.246,-0.827) -- (-0.239,-0.825) -- (-0.233,-0.822) -- (-0.227,-0.820) -- (-0.220,-0.817) -- (-0.214,-0.814) -- (-0.208,-0.812) -- (-0.201,-0.809) -- (-0.195,-0.806) -- (-0.189,-0.804) -- (-0.182,-0.801) -- (-0.176,-0.798) -- (-0.170,-0.796) -- (-0.163,-0.793) -- (-0.157,-0.791) -- (-0.151,-0.788) -- (-0.144,-0.785) -- (-0.138,-0.783) -- (-0.132,-0.780) -- (-0.125,-0.777) -- (-0.119,-0.775) -- (-0.113,-0.772) -- (-0.106,-0.770) -- (-0.100,-0.767) -- (-0.093,-0.764) -- (-0.087,-0.762) -- (-0.081,-0.759) -- (-0.074,-0.756) -- (-0.068,-0.754) -- (-0.062,-0.751) -- (-0.055,-0.749) -- (-0.049,-0.746) -- (-0.043,-0.743) -- (-0.036,-0.741) -- (-0.030,-0.738) -- (-0.024,-0.735) -- (-0.017,-0.733) -- (-0.011,-0.730) -- (-0.005,-0.727) -- (0.002,-0.725) -- (0.008,-0.722) -- (0.014,-0.720) -- (0.021,-0.717) -- (0.027,-0.714) -- (0.033,-0.712) -- (0.040,-0.709) -- (0.046,-0.706) -- (0.052,-0.704) -- (0.059,-0.701) -- (0.065,-0.699) -- (0.071,-0.696) -- (0.078,-0.693) -- (0.084,-0.691) -- (0.091,-0.688) -- (0.097,-0.685) -- (0.103,-0.683) -- (0.110,-0.680) -- (0.116,-0.677) -- (0.122,-0.675) -- (0.129,-0.672) -- (0.135,-0.670) -- (0.141,-0.667) -- (0.148,-0.664) -- (0.154,-0.662) -- (0.160,-0.659) -- (0.167,-0.656) -- (0.173,-0.654) -- (0.179,-0.651) -- (0.186,-0.649) -- (0.192,-0.646) -- (0.198,-0.643) -- (0.205,-0.641) -- (0.211,-0.638) -- (0.217,-0.635) -- (0.224,-0.633) -- (0.230,-0.630) -- (0.236,-0.627) -- (0.243,-0.625) -- (0.249,-0.622) -- (0.256,-0.620) -- (0.262,-0.617) -- (0.268,-0.614) -- (0.275,-0.612) -- (0.281,-0.609) -- (0.287,-0.606) -- (0.294,-0.604) -- (0.300,-0.601) -- (0.306,-0.599) -- (0.313,-0.596) -- (0.319,-0.593) -- (0.325,-0.591) -- (0.332,-0.588) -- (0.338,-0.585) -- (0.344,-0.583) -- (0.351,-0.580) -- (0.357,-0.577) -- (0.363,-0.575) -- (0.370,-0.572) -- (0.376,-0.570) -- (0.382,-0.567) -- (0.389,-0.564) -- (0.395,-0.562) -- (0.401,-0.559) -- (0.408,-0.556) -- (0.414,-0.554) -- (0.421,-0.551) -- (0.427,-0.549) -- (0.433,-0.546) -- (0.440,-0.543) -- (0.446,-0.541) -- (0.452,-0.538) -- (0.459,-0.535) -- (0.465,-0.533) -- (0.471,-0.530) -- (0.478,-0.528) -- (0.484,-0.525) -- (0.490,-0.522) -- (0.497,-0.520) -- (0.503,-0.517) -- (0.509,-0.514) -- (0.516,-0.512) -- (0.522,-0.509) -- (0.528,-0.506) -- (0.535,-0.504) -- (0.541,-0.501) -- (0.547,-0.499) -- (0.554,-0.496) -- (0.560,-0.493) -- (0.566,-0.491) -- (0.573,-0.488) -- (0.579,-0.485) -- (0.586,-0.483) -- (0.592,-0.480) -- (0.598,-0.478) -- (0.605,-0.475) -- (0.611,-0.472) -- (0.617,-0.470) -- (0.624,-0.467) -- (0.630,-0.464) -- (0.636,-0.462) -- (0.643,-0.459) -- (0.649,-0.456) -- (0.655,-0.454) -- (0.662,-0.451) -- (0.668,-0.449) -- (0.674,-0.446) -- (0.681,-0.443) -- (0.687,-0.441) -- (0.693,-0.438) -- (0.700,-0.435) -- (0.706,-0.433) -- (0.712,-0.430) -- (0.719,-0.428) -- (0.725,-0.425) -- (0.731,-0.422) -- (0.738,-0.420) -- (0.744,-0.417) -- (0.751,-0.414) -- (0.757,-0.412) -- (0.763,-0.409) -- (0.770,-0.406) -- (0.776,-0.404) -- (0.782,-0.401) -- (0.789,-0.399) -- (0.795,-0.396) -- (0.801,-0.393) -- (0.808,-0.391) -- (0.814,-0.388) -- (0.820,-0.385) -- (0.827,-0.383) -- (0.833,-0.380) -- (0.839,-0.378) -- (0.846,-0.375) -- (0.852,-0.372) -- (0.858,-0.370) -- (0.865,-0.367) -- (0.871,-0.364) -- (0.877,-0.362) -- (0.884,-0.359) -- (0.890,-0.356) -- (0.896,-0.354) -- (0.903,-0.351) -- (0.909,-0.349) -- (0.916,-0.346) -- (0.922,-0.343) -- (0.928,-0.341) -- (0.935,-0.338) -- (0.941,-0.335) -- (0.947,-0.333) -- (0.954,-0.330) -- (0.960,-0.328) -- (0.966,-0.325) -- (0.973,-0.322) -- (0.979,-0.320) -- (0.985,-0.317) -- (0.992,-0.314) -- (0.998,-0.312) -- (1.004,-0.309) -- (1.011,-0.307) -- (1.017,-0.304) -- (1.023,-0.301) -- (1.030,-0.299) -- (1.036,-0.296) -- (1.042,-0.293) -- (1.049,-0.291) -- (1.055,-0.288) -- (1.061,-0.285) -- (1.068,-0.283) -- (1.074,-0.280) -- (1.081,-0.278) -- (1.087,-0.275) -- (1.093,-0.272) -- (1.100,-0.270) -- (1.106,-0.267) -- (1.112,-0.264) -- (1.119,-0.262) -- (1.125,-0.259) -- (1.131,-0.257) -- (1.138,-0.254) -- (1.144,-0.251) -- (1.150,-0.249) -- (1.157,-0.246) -- (1.163,-0.243) -- (1.169,-0.241) -- (1.176,-0.238) -- (1.182,-0.235) -- (1.188,-0.233) -- (1.195,-0.230) -- (1.201,-0.228) -- (1.207,-0.225) -- (1.214,-0.222) -- (1.220,-0.220) -- (1.226,-0.217) -- (1.233,-0.214) -- (1.239,-0.212) -- (1.245,-0.209) -- (1.252,-0.207) -- (1.258,-0.204) -- (1.265,-0.201) -- (1.271,-0.199) -- (1.277,-0.196) -- (1.284,-0.193) -- (1.290,-0.191) -- (1.296,-0.188) -- (1.303,-0.185) -- (1.309,-0.183) -- (1.315,-0.180) -- (1.322,-0.178) -- (1.328,-0.175) -- (1.334,-0.172) -- (1.341,-0.170) -- (1.347,-0.167) -- (1.353,-0.164) -- (1.360,-0.162) -- (1.366,-0.159) -- (1.372,-0.157) -- (1.379,-0.154) -- (1.385,-0.151) -- (1.391,-0.149) -- (1.398,-0.146) -- (1.404,-0.143) -- (1.410,-0.141) -- (1.417,-0.138) -- (1.423,-0.135) -- (1.430,-0.133) -- (1.436,-0.130) -- (1.442,-0.128) -- (1.449,-0.125) -- (1.455,-0.122) -- (1.461,-0.120) -- (1.468,-0.117) -- (1.474,-0.114) -- (1.480,-0.112) -- (1.487,-0.109) -- (1.493,-0.107) -- (1.499,-0.104) -- (1.506,-0.101) -- (1.512,-0.099) -- (1.518,-0.096) -- (1.525,-0.093) -- (1.531,-0.091) -- (1.537,-0.088) -- (1.544,-0.086) -- (1.550,-0.083) -- (1.556,-0.080) -- (1.563,-0.078) -- (1.569,-0.075) -- (1.575,-0.072) -- (1.582,-0.070) -- (1.588,-0.067) -- (1.595,-0.064) -- (1.601,-0.062) -- (1.607,-0.059) -- (1.614,-0.057) -- (1.620,-0.054) -- (1.626,-0.051) -- (1.633,-0.049) -- (1.639,-0.046) -- (1.645,-0.043) -- (1.652,-0.041) -- (1.658,-0.038) -- (1.664,-0.036) -- (1.671,-0.033) -- (1.677,-0.030) -- (1.683,-0.028) -- (1.690,-0.025) -- (1.696,-0.022) -- (1.702,-0.020) -- (1.709,-0.017) -- (1.715,-0.014) -- (1.721,-0.012) -- (1.728,-0.009) -- (1.734,-0.007) -- (1.740,-0.004) -- (1.744,-0.001) -- (1.744,0.001) -- (1.740,0.004) -- (1.734,0.007) -- (1.728,0.009) -- (1.721,0.012) -- (1.715,0.014) -- (1.709,0.017) -- (1.702,0.020) -- (1.696,0.022) -- (1.690,0.025) -- (1.683,0.028) -- (1.677,0.030) -- (1.671,0.033) -- (1.664,0.036) -- (1.658,0.038) -- (1.652,0.041) -- (1.645,0.043) -- (1.639,0.046) -- (1.633,0.049) -- (1.626,0.051) -- (1.620,0.054) -- (1.614,0.057) -- (1.607,0.059) -- (1.601,0.062) -- (1.595,0.064) -- (1.588,0.067) -- (1.582,0.070) -- (1.575,0.072) -- (1.569,0.075) -- (1.563,0.078) -- (1.556,0.080) -- (1.550,0.083) -- (1.544,0.086) -- (1.537,0.088) -- (1.531,0.091) -- (1.525,0.093) -- (1.518,0.096) -- (1.512,0.099) -- (1.506,0.101) -- (1.499,0.104) -- (1.493,0.107) -- (1.487,0.109) -- (1.480,0.112) -- (1.474,0.114) -- (1.468,0.117) -- (1.461,0.120) -- (1.455,0.122) -- (1.449,0.125) -- (1.442,0.128) -- (1.436,0.130) -- (1.430,0.133) -- (1.423,0.135) -- (1.417,0.138) -- (1.410,0.141) -- (1.404,0.143) -- (1.398,0.146) -- (1.391,0.149) -- (1.385,0.151) -- (1.379,0.154) -- (1.372,0.157) -- (1.366,0.159) -- (1.360,0.162) -- (1.353,0.164) -- (1.347,0.167) -- (1.341,0.170) -- (1.334,0.172) -- (1.328,0.175) -- (1.322,0.178) -- (1.315,0.180) -- (1.309,0.183) -- (1.303,0.185) -- (1.296,0.188) -- (1.290,0.191) -- (1.284,0.193) -- (1.277,0.196) -- (1.271,0.199) -- (1.265,0.201) -- (1.258,0.204) -- (1.252,0.207) -- (1.245,0.209) -- (1.239,0.212) -- (1.233,0.214) -- (1.226,0.217) -- (1.220,0.220) -- (1.214,0.222) -- (1.207,0.225) -- (1.201,0.228) -- (1.195,0.230) -- (1.188,0.233) -- (1.182,0.235) -- (1.176,0.238) -- (1.169,0.241) -- (1.163,0.243) -- (1.157,0.246) -- (1.150,0.249) -- (1.144,0.251) -- (1.138,0.254) -- (1.131,0.257) -- (1.125,0.259) -- (1.119,0.262) -- (1.112,0.264) -- (1.106,0.267) -- (1.100,0.270) -- (1.093,0.272) -- (1.087,0.275) -- (1.081,0.278) -- (1.074,0.280) -- (1.068,0.283) -- (1.061,0.285) -- (1.055,0.288) -- (1.049,0.291) -- (1.042,0.293) -- (1.036,0.296) -- (1.030,0.299) -- (1.023,0.301) -- (1.017,0.304) -- (1.011,0.307) -- (1.004,0.309) -- (0.998,0.312) -- (0.992,0.314) -- (0.985,0.317) -- (0.979,0.320) -- (0.973,0.322) -- (0.966,0.325) -- (0.960,0.328) -- (0.954,0.330) -- (0.947,0.333) -- (0.941,0.335) -- (0.935,0.338) -- (0.928,0.341) -- (0.922,0.343) -- (0.916,0.346) -- (0.909,0.349) -- (0.903,0.351) -- (0.896,0.354) -- (0.890,0.356) -- (0.884,0.359) -- (0.877,0.362) -- (0.871,0.364) -- (0.865,0.367) -- (0.858,0.370) -- (0.852,0.372) -- (0.846,0.375) -- (0.839,0.378) -- (0.833,0.380) -- (0.827,0.383) -- (0.820,0.385) -- (0.814,0.388) -- (0.808,0.391) -- (0.801,0.393) -- (0.795,0.396) -- (0.789,0.399) -- (0.782,0.401) -- (0.776,0.404) -- (0.770,0.406) -- (0.763,0.409) -- (0.757,0.412) -- (0.751,0.414) -- (0.744,0.417) -- (0.738,0.420) -- (0.731,0.422) -- (0.725,0.425) -- (0.719,0.428) -- (0.712,0.430) -- (0.706,0.433) -- (0.700,0.435) -- (0.693,0.438) -- (0.687,0.441) -- (0.681,0.443) -- (0.674,0.446) -- (0.668,0.449) -- (0.662,0.451) -- (0.655,0.454) -- (0.649,0.456) -- (0.643,0.459) -- (0.636,0.462) -- (0.630,0.464) -- (0.624,0.467) -- (0.617,0.470) -- (0.611,0.472) -- (0.605,0.475) -- (0.598,0.478) -- (0.592,0.480) -- (0.586,0.483) -- (0.579,0.485) -- (0.573,0.488) -- (0.566,0.491) -- (0.560,0.493) -- (0.554,0.496) -- (0.547,0.499) -- (0.541,0.501) -- (0.535,0.504) -- (0.528,0.506) -- (0.522,0.509) -- (0.516,0.512) -- (0.509,0.514) -- (0.503,0.517) -- (0.497,0.520) -- (0.490,0.522) -- (0.484,0.525) -- (0.478,0.528) -- (0.471,0.530) -- (0.465,0.533) -- (0.459,0.535) -- (0.452,0.538) -- (0.446,0.541) -- (0.440,0.543) -- (0.433,0.546) -- (0.427,0.549) -- (0.421,0.551) -- (0.414,0.554) -- (0.408,0.556) -- (0.401,0.559) -- (0.395,0.562) -- (0.389,0.564) -- (0.382,0.567) -- (0.376,0.570) -- (0.370,0.572) -- (0.363,0.575) -- (0.357,0.577) -- (0.351,0.580) -- (0.344,0.583) -- (0.338,0.585) -- (0.332,0.588) -- (0.325,0.591) -- (0.319,0.593) -- (0.313,0.596) -- (0.306,0.599) -- (0.300,0.601) -- (0.294,0.604) -- (0.287,0.606) -- (0.281,0.609) -- (0.275,0.612) -- (0.268,0.614) -- (0.262,0.617) -- (0.256,0.620) -- (0.249,0.622) -- (0.243,0.625) -- (0.236,0.627) -- (0.230,0.630) -- (0.224,0.633) -- (0.217,0.635) -- (0.211,0.638) -- (0.205,0.641) -- (0.198,0.643) -- (0.192,0.646) -- (0.186,0.649) -- (0.179,0.651) -- (0.173,0.654) -- (0.167,0.656) -- (0.160,0.659) -- (0.154,0.662) -- (0.148,0.664) -- (0.141,0.667) -- (0.135,0.670) -- (0.129,0.672) -- (0.122,0.675) -- (0.116,0.677) -- (0.110,0.680) -- (0.103,0.683) -- (0.097,0.685) -- (0.091,0.688) -- (0.084,0.691) -- (0.078,0.693) -- (0.071,0.696) -- (0.065,0.699) -- (0.059,0.701) -- (0.052,0.704) -- (0.046,0.706) -- (0.040,0.709) -- (0.033,0.712) -- (0.027,0.714) -- (0.021,0.717) -- (0.014,0.720) -- (0.008,0.722) -- (0.002,0.725) -- (-0.005,0.727) -- (-0.011,0.730) -- (-0.017,0.733) -- (-0.024,0.735) -- (-0.030,0.738) -- (-0.036,0.741) -- (-0.043,0.743) -- (-0.049,0.746) -- (-0.055,0.749) -- (-0.062,0.751) -- (-0.068,0.754) -- (-0.074,0.756) -- (-0.081,0.759) -- (-0.087,0.762) -- (-0.093,0.764) -- (-0.100,0.767) -- (-0.106,0.770) -- (-0.113,0.772) -- (-0.119,0.775) -- (-0.125,0.777) -- (-0.132,0.780) -- (-0.138,0.783) -- (-0.144,0.785) -- (-0.151,0.788) -- (-0.157,0.791) -- (-0.163,0.793) -- (-0.170,0.796) -- (-0.176,0.798) -- (-0.182,0.801) -- (-0.189,0.804) -- (-0.195,0.806) -- (-0.201,0.809) -- (-0.208,0.812) -- (-0.214,0.814) -- (-0.220,0.817) -- (-0.227,0.820) -- (-0.233,0.822) -- (-0.239,0.825) -- (-0.246,0.827) -- (-0.251,0.830) -- (-0.256,0.832) -- cycle}

      \fill[top color=azulvelo!70!white, bottom color=metal] \formagota;
      % Conductor con punta - densidad superficial
% Generado por scripts/generar_datos.py (no editar a mano).
\SG{-0.259}{0.833}{-0.262}{0.834}{0.183}
\SG{-0.262}{0.834}{-0.265}{0.835}{0.181}
\SG{-0.265}{0.835}{-0.268}{0.836}{0.179}
\SG{-0.268}{0.836}{-0.270}{0.837}{0.179}
\SG{-0.270}{0.837}{-0.273}{0.839}{0.181}
\SG{-0.273}{0.839}{-0.276}{0.840}{0.183}
\SG{-0.276}{0.840}{-0.279}{0.841}{0.185}
\SG{-0.279}{0.841}{-0.281}{0.842}{0.187}
\SG{-0.281}{0.842}{-0.284}{0.843}{0.190}
\SG{-0.284}{0.843}{-0.287}{0.844}{0.194}
\SG{-0.287}{0.844}{-0.290}{0.845}{0.197}
\SG{-0.290}{0.845}{-0.293}{0.846}{0.201}
\SG{-0.293}{0.846}{-0.295}{0.847}{0.203}
\SG{-0.295}{0.847}{-0.298}{0.848}{0.206}
\SG{-0.298}{0.848}{-0.301}{0.849}{0.203}
\SG{-0.301}{0.849}{-0.304}{0.850}{0.201}
\SG{-0.304}{0.850}{-0.306}{0.851}{0.200}
\SG{-0.306}{0.851}{-0.309}{0.852}{0.199}
\SG{-0.309}{0.852}{-0.312}{0.853}{0.199}
\SG{-0.312}{0.853}{-0.315}{0.854}{0.201}
\SG{-0.315}{0.854}{-0.318}{0.855}{0.204}
\SG{-0.318}{0.855}{-0.320}{0.855}{0.206}
\SG{-0.320}{0.855}{-0.323}{0.856}{0.209}
\SG{-0.323}{0.856}{-0.326}{0.857}{0.212}
\SG{-0.326}{0.857}{-0.329}{0.858}{0.216}
\SG{-0.329}{0.858}{-0.332}{0.859}{0.221}
\SG{-0.332}{0.859}{-0.335}{0.860}{0.225}
\SG{-0.335}{0.860}{-0.337}{0.861}{0.229}
\SG{-0.337}{0.861}{-0.340}{0.862}{0.233}
\SG{-0.340}{0.862}{-0.343}{0.863}{0.234}
\SG{-0.343}{0.863}{-0.346}{0.863}{0.232}
\SG{-0.346}{0.863}{-0.349}{0.864}{0.231}
\SG{-0.349}{0.864}{-0.352}{0.865}{0.231}
\SG{-0.352}{0.865}{-0.354}{0.866}{0.231}
\SG{-0.354}{0.866}{-0.357}{0.867}{0.232}
\SG{-0.357}{0.867}{-0.360}{0.867}{0.236}
\SG{-0.360}{0.867}{-0.363}{0.868}{0.239}
\SG{-0.363}{0.868}{-0.366}{0.869}{0.242}
\SG{-0.366}{0.869}{-0.369}{0.870}{0.246}
\SG{-0.369}{0.870}{-0.371}{0.870}{0.250}
\SG{-0.371}{0.870}{-0.374}{0.871}{0.257}
\SG{-0.374}{0.871}{-0.377}{0.872}{0.262}
\SG{-0.377}{0.872}{-0.380}{0.873}{0.266}
\SG{-0.380}{0.873}{-0.383}{0.873}{0.271}
\SG{-0.383}{0.873}{-0.386}{0.874}{0.273}
\SG{-0.386}{0.874}{-0.389}{0.875}{0.269}
\SG{-0.389}{0.875}{-0.392}{0.876}{0.266}
\SG{-0.392}{0.876}{-0.394}{0.876}{0.262}
\SG{-0.394}{0.876}{-0.397}{0.877}{0.259}
\SG{-0.397}{0.877}{-0.400}{0.878}{0.257}
\SG{-0.400}{0.878}{-0.403}{0.878}{0.253}
\SG{-0.403}{0.878}{-0.406}{0.879}{0.251}
\SG{-0.406}{0.879}{-0.409}{0.879}{0.248}
\SG{-0.409}{0.879}{-0.412}{0.880}{0.246}
\SG{-0.412}{0.880}{-0.415}{0.881}{0.244}
\SG{-0.415}{0.881}{-0.418}{0.881}{0.244}
\SG{-0.418}{0.881}{-0.420}{0.882}{0.245}
\SG{-0.420}{0.882}{-0.423}{0.882}{0.245}
\SG{-0.423}{0.882}{-0.426}{0.883}{0.246}
\SG{-0.426}{0.883}{-0.429}{0.884}{0.247}
\SG{-0.429}{0.884}{-0.432}{0.884}{0.252}
\SG{-0.432}{0.884}{-0.435}{0.885}{0.257}
\SG{-0.435}{0.885}{-0.438}{0.885}{0.262}
\SG{-0.438}{0.885}{-0.441}{0.886}{0.266}
\SG{-0.441}{0.886}{-0.444}{0.886}{0.271}
\SG{-0.444}{0.886}{-0.447}{0.887}{0.279}
\SG{-0.447}{0.887}{-0.449}{0.887}{0.286}
\SG{-0.449}{0.887}{-0.452}{0.888}{0.292}
\SG{-0.452}{0.888}{-0.455}{0.888}{0.299}
\SG{-0.455}{0.888}{-0.458}{0.889}{0.303}
\SG{-0.458}{0.889}{-0.461}{0.889}{0.302}
\SG{-0.461}{0.889}{-0.464}{0.890}{0.301}
\SG{-0.464}{0.890}{-0.467}{0.890}{0.299}
\SG{-0.467}{0.890}{-0.470}{0.891}{0.298}
\SG{-0.470}{0.891}{-0.473}{0.891}{0.297}
\SG{-0.473}{0.891}{-0.476}{0.891}{0.295}
\SG{-0.476}{0.891}{-0.479}{0.892}{0.293}
\SG{-0.479}{0.892}{-0.482}{0.892}{0.290}
\SG{-0.482}{0.892}{-0.485}{0.893}{0.289}
\SG{-0.485}{0.893}{-0.487}{0.893}{0.287}
\SG{-0.487}{0.893}{-0.490}{0.893}{0.285}
\SG{-0.490}{0.893}{-0.493}{0.894}{0.283}
\SG{-0.493}{0.894}{-0.496}{0.894}{0.281}
\SG{-0.496}{0.894}{-0.499}{0.894}{0.280}
\SG{-0.499}{0.894}{-0.502}{0.895}{0.278}
\SG{-0.502}{0.895}{-0.505}{0.895}{0.277}
\SG{-0.505}{0.895}{-0.508}{0.895}{0.276}
\SG{-0.508}{0.895}{-0.511}{0.896}{0.275}
\SG{-0.511}{0.896}{-0.514}{0.896}{0.273}
\SG{-0.514}{0.896}{-0.517}{0.896}{0.272}
\SG{-0.517}{0.896}{-0.520}{0.896}{0.271}
\SG{-0.520}{0.896}{-0.523}{0.897}{0.271}
\SG{-0.523}{0.897}{-0.526}{0.897}{0.270}
\SG{-0.526}{0.897}{-0.529}{0.897}{0.269}
\SG{-0.529}{0.897}{-0.532}{0.897}{0.268}
\SG{-0.532}{0.897}{-0.535}{0.898}{0.268}
\SG{-0.535}{0.898}{-0.537}{0.898}{0.267}
\SG{-0.537}{0.898}{-0.540}{0.898}{0.267}
\SG{-0.540}{0.898}{-0.543}{0.898}{0.266}
\SG{-0.543}{0.898}{-0.546}{0.898}{0.266}
\SG{-0.546}{0.898}{-0.549}{0.899}{0.266}
\SG{-0.549}{0.899}{-0.552}{0.899}{0.266}
\SG{-0.552}{0.899}{-0.555}{0.899}{0.265}
\SG{-0.555}{0.899}{-0.558}{0.899}{0.265}
\SG{-0.558}{0.899}{-0.561}{0.899}{0.265}
\SG{-0.561}{0.899}{-0.564}{0.899}{0.265}
\SG{-0.564}{0.899}{-0.567}{0.899}{0.265}
\SG{-0.567}{0.899}{-0.570}{0.899}{0.265}
\SG{-0.570}{0.899}{-0.573}{0.900}{0.265}
\SG{-0.573}{0.900}{-0.576}{0.900}{0.265}
\SG{-0.576}{0.900}{-0.579}{0.900}{0.266}
\SG{-0.579}{0.900}{-0.582}{0.900}{0.266}
\SG{-0.582}{0.900}{-0.585}{0.900}{0.266}
\SG{-0.585}{0.900}{-0.588}{0.900}{0.267}
\SG{-0.588}{0.900}{-0.591}{0.900}{0.267}
\SG{-0.591}{0.900}{-0.594}{0.900}{0.268}
\SG{-0.594}{0.900}{-0.597}{0.900}{0.268}
\SG{-0.597}{0.900}{-0.599}{0.900}{0.269}
\SG{-0.599}{0.900}{-0.602}{0.900}{0.269}
\SG{-0.602}{0.900}{-0.605}{0.900}{0.270}
\SG{-0.605}{0.900}{-0.608}{0.900}{0.270}
\SG{-0.608}{0.900}{-0.611}{0.900}{0.271}
\SG{-0.611}{0.900}{-0.614}{0.900}{0.272}
\SG{-0.614}{0.900}{-0.617}{0.900}{0.273}
\SG{-0.617}{0.900}{-0.620}{0.900}{0.273}
\SG{-0.620}{0.900}{-0.623}{0.900}{0.274}
\SG{-0.623}{0.900}{-0.626}{0.900}{0.275}
\SG{-0.626}{0.900}{-0.629}{0.900}{0.276}
\SG{-0.629}{0.900}{-0.632}{0.899}{0.277}
\SG{-0.632}{0.899}{-0.635}{0.899}{0.278}
\SG{-0.635}{0.899}{-0.638}{0.899}{0.280}
\SG{-0.638}{0.899}{-0.641}{0.899}{0.281}
\SG{-0.641}{0.899}{-0.644}{0.899}{0.282}
\SG{-0.644}{0.899}{-0.647}{0.899}{0.283}
\SG{-0.647}{0.899}{-0.650}{0.899}{0.284}
\SG{-0.650}{0.899}{-0.653}{0.898}{0.286}
\SG{-0.653}{0.898}{-0.656}{0.898}{0.287}
\SG{-0.656}{0.898}{-0.659}{0.898}{0.289}
\SG{-0.659}{0.898}{-0.661}{0.898}{0.290}
\SG{-0.661}{0.898}{-0.664}{0.898}{0.292}
\SG{-0.664}{0.898}{-0.667}{0.897}{0.294}
\SG{-0.667}{0.897}{-0.670}{0.897}{0.295}
\SG{-0.670}{0.897}{-0.673}{0.897}{0.297}
\SG{-0.673}{0.897}{-0.676}{0.897}{0.299}
\SG{-0.676}{0.897}{-0.679}{0.896}{0.301}
\SG{-0.679}{0.896}{-0.682}{0.896}{0.303}
\SG{-0.682}{0.896}{-0.685}{0.896}{0.305}
\SG{-0.685}{0.896}{-0.688}{0.896}{0.307}
\SG{-0.688}{0.896}{-0.691}{0.895}{0.310}
\SG{-0.691}{0.895}{-0.694}{0.895}{0.312}
\SG{-0.694}{0.895}{-0.697}{0.895}{0.315}
\SG{-0.697}{0.895}{-0.700}{0.894}{0.317}
\SG{-0.700}{0.894}{-0.703}{0.894}{0.320}
\SG{-0.703}{0.894}{-0.706}{0.894}{0.322}
\SG{-0.706}{0.894}{-0.709}{0.893}{0.325}
\SG{-0.709}{0.893}{-0.711}{0.893}{0.328}
\SG{-0.711}{0.893}{-0.714}{0.893}{0.331}
\SG{-0.714}{0.893}{-0.717}{0.892}{0.334}
\SG{-0.717}{0.892}{-0.720}{0.892}{0.338}
\SG{-0.720}{0.892}{-0.723}{0.892}{0.341}
\SG{-0.723}{0.892}{-0.726}{0.891}{0.344}
\SG{-0.726}{0.891}{-0.729}{0.891}{0.348}
\SG{-0.729}{0.891}{-0.732}{0.890}{0.351}
\SG{-0.732}{0.890}{-0.735}{0.890}{0.354}
\SG{-0.735}{0.890}{-0.738}{0.889}{0.357}
\SG{-0.738}{0.889}{-0.741}{0.889}{0.359}
\SG{-0.741}{0.889}{-0.744}{0.888}{0.362}
\SG{-0.744}{0.888}{-0.747}{0.888}{0.365}
\SG{-0.747}{0.888}{-0.750}{0.887}{0.364}
\SG{-0.750}{0.887}{-0.752}{0.887}{0.359}
\SG{-0.752}{0.887}{-0.755}{0.886}{0.353}
\SG{-0.755}{0.886}{-0.758}{0.886}{0.348}
\SG{-0.758}{0.886}{-0.761}{0.885}{0.342}
\SG{-0.761}{0.885}{-0.764}{0.885}{0.338}
\SG{-0.764}{0.885}{-0.767}{0.884}{0.334}
\SG{-0.767}{0.884}{-0.770}{0.884}{0.331}
\SG{-0.770}{0.884}{-0.773}{0.883}{0.328}
\SG{-0.773}{0.883}{-0.776}{0.883}{0.324}
\SG{-0.776}{0.883}{-0.779}{0.882}{0.323}
\SG{-0.779}{0.882}{-0.781}{0.881}{0.324}
\SG{-0.781}{0.881}{-0.784}{0.881}{0.325}
\SG{-0.784}{0.881}{-0.787}{0.880}{0.326}
\SG{-0.787}{0.880}{-0.790}{0.880}{0.327}
\SG{-0.790}{0.880}{-0.793}{0.879}{0.330}
\SG{-0.793}{0.879}{-0.796}{0.878}{0.334}
\SG{-0.796}{0.878}{-0.799}{0.878}{0.338}
\SG{-0.799}{0.878}{-0.802}{0.877}{0.342}
\SG{-0.802}{0.877}{-0.805}{0.876}{0.346}
\SG{-0.805}{0.876}{-0.807}{0.876}{0.351}
\SG{-0.807}{0.876}{-0.810}{0.875}{0.355}
\SG{-0.810}{0.875}{-0.813}{0.874}{0.360}
\SG{-0.813}{0.874}{-0.816}{0.874}{0.365}
\SG{-0.816}{0.874}{-0.819}{0.873}{0.370}
\SG{-0.819}{0.873}{-0.822}{0.872}{0.373}
\SG{-0.822}{0.872}{-0.825}{0.871}{0.370}
\SG{-0.825}{0.871}{-0.828}{0.871}{0.367}
\SG{-0.828}{0.871}{-0.830}{0.870}{0.363}
\SG{-0.830}{0.870}{-0.833}{0.869}{0.359}
\SG{-0.833}{0.869}{-0.836}{0.868}{0.355}
\SG{-0.836}{0.868}{-0.839}{0.868}{0.353}
\SG{-0.839}{0.868}{-0.842}{0.867}{0.352}
\SG{-0.842}{0.867}{-0.845}{0.866}{0.350}
\SG{-0.845}{0.866}{-0.847}{0.865}{0.349}
\SG{-0.847}{0.865}{-0.850}{0.864}{0.348}
\SG{-0.850}{0.864}{-0.853}{0.864}{0.350}
\SG{-0.853}{0.864}{-0.856}{0.863}{0.352}
\SG{-0.856}{0.863}{-0.859}{0.862}{0.355}
\SG{-0.859}{0.862}{-0.862}{0.861}{0.358}
\SG{-0.862}{0.861}{-0.864}{0.860}{0.360}
\SG{-0.864}{0.860}{-0.867}{0.859}{0.358}
\SG{-0.867}{0.859}{-0.870}{0.858}{0.356}
\SG{-0.870}{0.858}{-0.873}{0.858}{0.354}
\SG{-0.873}{0.858}{-0.876}{0.857}{0.351}
\SG{-0.876}{0.857}{-0.879}{0.856}{0.349}
\SG{-0.879}{0.856}{-0.881}{0.855}{0.348}
\SG{-0.881}{0.855}{-0.884}{0.854}{0.347}
\SG{-0.884}{0.854}{-0.887}{0.853}{0.346}
\SG{-0.887}{0.853}{-0.890}{0.852}{0.345}
\SG{-0.890}{0.852}{-0.893}{0.851}{0.345}
\SG{-0.893}{0.851}{-0.895}{0.850}{0.347}
\SG{-0.895}{0.850}{-0.898}{0.849}{0.350}
\SG{-0.898}{0.849}{-0.901}{0.848}{0.353}
\SG{-0.901}{0.848}{-0.904}{0.847}{0.357}
\SG{-0.904}{0.847}{-0.907}{0.846}{0.362}
\SG{-0.907}{0.846}{-0.909}{0.845}{0.363}
\SG{-0.909}{0.845}{-0.912}{0.844}{0.362}
\SG{-0.912}{0.844}{-0.915}{0.843}{0.361}
\SG{-0.915}{0.843}{-0.918}{0.842}{0.360}
\SG{-0.918}{0.842}{-0.920}{0.841}{0.358}
\SG{-0.920}{0.841}{-0.923}{0.840}{0.356}
\SG{-0.923}{0.840}{-0.926}{0.839}{0.356}
\SG{-0.926}{0.839}{-0.929}{0.838}{0.355}
\SG{-0.929}{0.838}{-0.931}{0.837}{0.355}
\SG{-0.931}{0.837}{-0.934}{0.836}{0.355}
\SG{-0.934}{0.836}{-0.937}{0.835}{0.354}
\SG{-0.937}{0.835}{-0.940}{0.833}{0.353}
\SG{-0.940}{0.833}{-0.942}{0.832}{0.351}
\SG{-0.942}{0.832}{-0.945}{0.831}{0.350}
\SG{-0.945}{0.831}{-0.948}{0.830}{0.348}
\SG{-0.948}{0.830}{-0.951}{0.829}{0.345}
\SG{-0.951}{0.829}{-0.953}{0.828}{0.345}
\SG{-0.953}{0.828}{-0.956}{0.827}{0.346}
\SG{-0.956}{0.827}{-0.959}{0.825}{0.346}
\SG{-0.959}{0.825}{-0.961}{0.824}{0.346}
\SG{-0.961}{0.824}{-0.964}{0.823}{0.346}
\SG{-0.964}{0.823}{-0.967}{0.822}{0.348}
\SG{-0.967}{0.822}{-0.969}{0.821}{0.352}
\SG{-0.969}{0.821}{-0.972}{0.819}{0.357}
\SG{-0.972}{0.819}{-0.975}{0.818}{0.362}
\SG{-0.975}{0.818}{-0.978}{0.817}{0.369}
\SG{-0.978}{0.817}{-0.980}{0.816}{0.375}
\SG{-0.980}{0.816}{-0.983}{0.814}{0.377}
\SG{-0.983}{0.814}{-0.986}{0.813}{0.378}
\SG{-0.986}{0.813}{-0.988}{0.812}{0.379}
\SG{-0.988}{0.812}{-0.991}{0.811}{0.380}
\SG{-0.991}{0.811}{-0.994}{0.809}{0.380}
\SG{-0.994}{0.809}{-0.996}{0.808}{0.381}
\SG{-0.996}{0.808}{-0.999}{0.807}{0.383}
\SG{-0.999}{0.807}{-1.002}{0.805}{0.385}
\SG{-1.002}{0.805}{-1.004}{0.804}{0.388}
\SG{-1.004}{0.804}{-1.007}{0.803}{0.391}
\SG{-1.007}{0.803}{-1.009}{0.801}{0.394}
\SG{-1.009}{0.801}{-1.012}{0.800}{0.393}
\SG{-1.012}{0.800}{-1.015}{0.799}{0.392}
\SG{-1.015}{0.799}{-1.017}{0.797}{0.391}
\SG{-1.017}{0.797}{-1.020}{0.796}{0.390}
\SG{-1.020}{0.796}{-1.023}{0.795}{0.389}
\SG{-1.023}{0.795}{-1.025}{0.793}{0.386}
\SG{-1.025}{0.793}{-1.028}{0.792}{0.383}
\SG{-1.028}{0.792}{-1.030}{0.790}{0.380}
\SG{-1.030}{0.790}{-1.033}{0.789}{0.377}
\SG{-1.033}{0.789}{-1.036}{0.788}{0.373}
\SG{-1.036}{0.788}{-1.038}{0.786}{0.368}
\SG{-1.038}{0.786}{-1.041}{0.785}{0.367}
\SG{-1.041}{0.785}{-1.043}{0.783}{0.366}
\SG{-1.043}{0.783}{-1.046}{0.782}{0.365}
\SG{-1.046}{0.782}{-1.048}{0.780}{0.363}
\SG{-1.048}{0.780}{-1.051}{0.779}{0.363}
\SG{-1.051}{0.779}{-1.054}{0.777}{0.363}
\SG{-1.054}{0.777}{-1.056}{0.776}{0.363}
\SG{-1.056}{0.776}{-1.059}{0.774}{0.363}
\SG{-1.059}{0.774}{-1.061}{0.773}{0.363}
\SG{-1.061}{0.773}{-1.064}{0.771}{0.363}
\SG{-1.064}{0.771}{-1.066}{0.770}{0.363}
\SG{-1.066}{0.770}{-1.069}{0.768}{0.365}
\SG{-1.069}{0.768}{-1.071}{0.767}{0.367}
\SG{-1.071}{0.767}{-1.074}{0.765}{0.369}
\SG{-1.074}{0.765}{-1.076}{0.764}{0.372}
\SG{-1.076}{0.764}{-1.079}{0.762}{0.377}
\SG{-1.079}{0.762}{-1.081}{0.760}{0.382}
\SG{-1.081}{0.760}{-1.084}{0.759}{0.383}
\SG{-1.084}{0.759}{-1.086}{0.757}{0.385}
\SG{-1.086}{0.757}{-1.089}{0.756}{0.386}
\SG{-1.089}{0.756}{-1.091}{0.754}{0.388}
\SG{-1.091}{0.754}{-1.094}{0.752}{0.390}
\SG{-1.094}{0.752}{-1.096}{0.751}{0.391}
\SG{-1.096}{0.751}{-1.099}{0.749}{0.391}
\SG{-1.099}{0.749}{-1.101}{0.748}{0.391}
\SG{-1.101}{0.748}{-1.104}{0.746}{0.390}
\SG{-1.104}{0.746}{-1.106}{0.744}{0.388}
\SG{-1.106}{0.744}{-1.108}{0.743}{0.386}
\SG{-1.108}{0.743}{-1.111}{0.741}{0.384}
\SG{-1.111}{0.741}{-1.113}{0.739}{0.382}
\SG{-1.113}{0.739}{-1.116}{0.738}{0.380}
\SG{-1.116}{0.738}{-1.118}{0.736}{0.377}
\SG{-1.118}{0.736}{-1.121}{0.734}{0.374}
\SG{-1.121}{0.734}{-1.123}{0.732}{0.371}
\SG{-1.123}{0.732}{-1.125}{0.731}{0.368}
\SG{-1.125}{0.731}{-1.128}{0.729}{0.367}
\SG{-1.128}{0.729}{-1.130}{0.727}{0.367}
\SG{-1.130}{0.727}{-1.133}{0.725}{0.368}
\SG{-1.133}{0.725}{-1.135}{0.724}{0.368}
\SG{-1.135}{0.724}{-1.137}{0.722}{0.369}
\SG{-1.137}{0.722}{-1.140}{0.720}{0.372}
\SG{-1.140}{0.720}{-1.142}{0.718}{0.373}
\SG{-1.142}{0.718}{-1.144}{0.717}{0.375}
\SG{-1.144}{0.717}{-1.147}{0.715}{0.377}
\SG{-1.147}{0.715}{-1.149}{0.713}{0.379}
\SG{-1.149}{0.713}{-1.151}{0.711}{0.381}
\SG{-1.151}{0.711}{-1.154}{0.709}{0.384}
\SG{-1.154}{0.709}{-1.156}{0.708}{0.386}
\SG{-1.156}{0.708}{-1.158}{0.706}{0.388}
\SG{-1.158}{0.706}{-1.161}{0.704}{0.390}
\SG{-1.161}{0.704}{-1.163}{0.702}{0.391}
\SG{-1.163}{0.702}{-1.165}{0.700}{0.391}
\SG{-1.165}{0.700}{-1.168}{0.698}{0.392}
\SG{-1.168}{0.698}{-1.170}{0.697}{0.392}
\SG{-1.170}{0.697}{-1.172}{0.695}{0.391}
\SG{-1.172}{0.695}{-1.174}{0.693}{0.391}
\SG{-1.174}{0.693}{-1.177}{0.691}{0.390}
\SG{-1.177}{0.691}{-1.179}{0.689}{0.388}
\SG{-1.179}{0.689}{-1.181}{0.687}{0.387}
\SG{-1.181}{0.687}{-1.184}{0.685}{0.386}
\SG{-1.184}{0.685}{-1.186}{0.683}{0.385}
\SG{-1.186}{0.683}{-1.188}{0.681}{0.384}
\SG{-1.188}{0.681}{-1.190}{0.679}{0.382}
\SG{-1.190}{0.679}{-1.192}{0.677}{0.380}
\SG{-1.192}{0.677}{-1.195}{0.675}{0.379}
\SG{-1.195}{0.675}{-1.197}{0.674}{0.377}
\SG{-1.197}{0.674}{-1.199}{0.672}{0.376}
\SG{-1.199}{0.672}{-1.201}{0.670}{0.375}
\SG{-1.201}{0.670}{-1.204}{0.668}{0.374}
\SG{-1.204}{0.668}{-1.206}{0.666}{0.373}
\SG{-1.206}{0.666}{-1.208}{0.664}{0.372}
\SG{-1.208}{0.664}{-1.210}{0.662}{0.371}
\SG{-1.210}{0.662}{-1.212}{0.660}{0.370}
\SG{-1.212}{0.660}{-1.214}{0.658}{0.370}
\SG{-1.214}{0.658}{-1.217}{0.656}{0.369}
\SG{-1.217}{0.656}{-1.219}{0.654}{0.369}
\SG{-1.219}{0.654}{-1.221}{0.652}{0.368}
\SG{-1.221}{0.652}{-1.223}{0.649}{0.368}
\SG{-1.223}{0.649}{-1.225}{0.647}{0.367}
\SG{-1.225}{0.647}{-1.227}{0.645}{0.367}
\SG{-1.227}{0.645}{-1.229}{0.643}{0.367}
\SG{-1.229}{0.643}{-1.231}{0.641}{0.367}
\SG{-1.231}{0.641}{-1.234}{0.639}{0.367}
\SG{-1.234}{0.639}{-1.236}{0.637}{0.367}
\SG{-1.236}{0.637}{-1.238}{0.635}{0.367}
\SG{-1.238}{0.635}{-1.240}{0.633}{0.367}
\SG{-1.240}{0.633}{-1.242}{0.631}{0.367}
\SG{-1.242}{0.631}{-1.244}{0.629}{0.368}
\SG{-1.244}{0.629}{-1.246}{0.627}{0.368}
\SG{-1.246}{0.627}{-1.248}{0.624}{0.368}
\SG{-1.248}{0.624}{-1.250}{0.622}{0.369}
\SG{-1.250}{0.622}{-1.252}{0.620}{0.370}
\SG{-1.252}{0.620}{-1.254}{0.618}{0.370}
\SG{-1.254}{0.618}{-1.256}{0.616}{0.371}
\SG{-1.256}{0.616}{-1.258}{0.614}{0.372}
\SG{-1.258}{0.614}{-1.260}{0.612}{0.373}
\SG{-1.260}{0.612}{-1.262}{0.609}{0.374}
\SG{-1.262}{0.609}{-1.264}{0.607}{0.375}
\SG{-1.264}{0.607}{-1.266}{0.605}{0.376}
\SG{-1.266}{0.605}{-1.268}{0.603}{0.377}
\SG{-1.268}{0.603}{-1.270}{0.601}{0.379}
\SG{-1.270}{0.601}{-1.272}{0.598}{0.380}
\SG{-1.272}{0.598}{-1.274}{0.596}{0.381}
\SG{-1.274}{0.596}{-1.276}{0.594}{0.382}
\SG{-1.276}{0.594}{-1.278}{0.592}{0.384}
\SG{-1.278}{0.592}{-1.280}{0.590}{0.386}
\SG{-1.280}{0.590}{-1.282}{0.587}{0.388}
\SG{-1.282}{0.587}{-1.284}{0.585}{0.390}
\SG{-1.284}{0.585}{-1.286}{0.583}{0.392}
\SG{-1.286}{0.583}{-1.288}{0.581}{0.393}
\SG{-1.288}{0.581}{-1.290}{0.578}{0.394}
\SG{-1.290}{0.578}{-1.291}{0.576}{0.396}
\SG{-1.291}{0.576}{-1.293}{0.574}{0.398}
\SG{-1.293}{0.574}{-1.295}{0.571}{0.399}
\SG{-1.295}{0.571}{-1.297}{0.569}{0.400}
\SG{-1.297}{0.569}{-1.299}{0.567}{0.401}
\SG{-1.299}{0.567}{-1.301}{0.565}{0.402}
\SG{-1.301}{0.565}{-1.303}{0.562}{0.402}
\SG{-1.303}{0.562}{-1.305}{0.560}{0.402}
\SG{-1.305}{0.560}{-1.306}{0.558}{0.402}
\SG{-1.306}{0.558}{-1.308}{0.555}{0.401}
\SG{-1.308}{0.555}{-1.310}{0.553}{0.400}
\SG{-1.310}{0.553}{-1.312}{0.551}{0.398}
\SG{-1.312}{0.551}{-1.314}{0.548}{0.396}
\SG{-1.314}{0.548}{-1.315}{0.546}{0.393}
\SG{-1.315}{0.546}{-1.317}{0.544}{0.391}
\SG{-1.317}{0.544}{-1.319}{0.541}{0.389}
\SG{-1.319}{0.541}{-1.321}{0.539}{0.388}
\SG{-1.321}{0.539}{-1.323}{0.537}{0.387}
\SG{-1.323}{0.537}{-1.324}{0.534}{0.385}
\SG{-1.324}{0.534}{-1.326}{0.532}{0.384}
\SG{-1.326}{0.532}{-1.328}{0.529}{0.383}
\SG{-1.328}{0.529}{-1.330}{0.527}{0.383}
\SG{-1.330}{0.527}{-1.331}{0.525}{0.383}
\SG{-1.331}{0.525}{-1.333}{0.522}{0.384}
\SG{-1.333}{0.522}{-1.335}{0.520}{0.385}
\SG{-1.335}{0.520}{-1.336}{0.517}{0.389}
\SG{-1.336}{0.517}{-1.338}{0.515}{0.392}
\SG{-1.338}{0.515}{-1.340}{0.513}{0.396}
\SG{-1.340}{0.513}{-1.341}{0.510}{0.399}
\SG{-1.341}{0.510}{-1.343}{0.508}{0.401}
\SG{-1.343}{0.508}{-1.345}{0.505}{0.403}
\SG{-1.345}{0.505}{-1.346}{0.503}{0.406}
\SG{-1.346}{0.503}{-1.348}{0.500}{0.408}
\SG{-1.348}{0.500}{-1.350}{0.498}{0.410}
\SG{-1.350}{0.498}{-1.351}{0.495}{0.411}
\SG{-1.351}{0.495}{-1.353}{0.493}{0.412}
\SG{-1.353}{0.493}{-1.355}{0.490}{0.411}
\SG{-1.355}{0.490}{-1.356}{0.488}{0.410}
\SG{-1.356}{0.488}{-1.358}{0.485}{0.408}
\SG{-1.358}{0.485}{-1.359}{0.483}{0.407}
\SG{-1.359}{0.483}{-1.361}{0.480}{0.406}
\SG{-1.361}{0.480}{-1.363}{0.478}{0.405}
\SG{-1.363}{0.478}{-1.364}{0.475}{0.402}
\SG{-1.364}{0.475}{-1.366}{0.473}{0.398}
\SG{-1.366}{0.473}{-1.367}{0.470}{0.394}
\SG{-1.367}{0.470}{-1.369}{0.468}{0.392}
\SG{-1.369}{0.468}{-1.370}{0.465}{0.390}
\SG{-1.370}{0.465}{-1.372}{0.463}{0.388}
\SG{-1.372}{0.463}{-1.373}{0.460}{0.388}
\SG{-1.373}{0.460}{-1.375}{0.458}{0.388}
\SG{-1.375}{0.458}{-1.376}{0.455}{0.388}
\SG{-1.376}{0.455}{-1.378}{0.453}{0.388}
\SG{-1.378}{0.453}{-1.379}{0.450}{0.389}
\SG{-1.379}{0.450}{-1.381}{0.448}{0.389}
\SG{-1.381}{0.448}{-1.382}{0.445}{0.389}
\SG{-1.382}{0.445}{-1.384}{0.442}{0.391}
\SG{-1.384}{0.442}{-1.385}{0.440}{0.393}
\SG{-1.385}{0.440}{-1.387}{0.437}{0.394}
\SG{-1.387}{0.437}{-1.388}{0.435}{0.395}
\SG{-1.388}{0.435}{-1.389}{0.432}{0.398}
\SG{-1.389}{0.432}{-1.391}{0.430}{0.403}
\SG{-1.391}{0.430}{-1.392}{0.427}{0.407}
\SG{-1.392}{0.427}{-1.394}{0.424}{0.411}
\SG{-1.394}{0.424}{-1.395}{0.422}{0.414}
\SG{-1.395}{0.422}{-1.396}{0.419}{0.417}
\SG{-1.396}{0.419}{-1.398}{0.416}{0.420}
\SG{-1.398}{0.416}{-1.399}{0.414}{0.422}
\SG{-1.399}{0.414}{-1.401}{0.411}{0.423}
\SG{-1.401}{0.411}{-1.402}{0.409}{0.424}
\SG{-1.402}{0.409}{-1.403}{0.406}{0.426}
\SG{-1.403}{0.406}{-1.405}{0.403}{0.426}
\SG{-1.405}{0.403}{-1.406}{0.401}{0.423}
\SG{-1.406}{0.401}{-1.407}{0.398}{0.420}
\SG{-1.407}{0.398}{-1.408}{0.395}{0.418}
\SG{-1.408}{0.395}{-1.410}{0.393}{0.416}
\SG{-1.410}{0.393}{-1.411}{0.390}{0.414}
\SG{-1.411}{0.390}{-1.412}{0.387}{0.415}
\SG{-1.412}{0.387}{-1.414}{0.385}{0.415}
\SG{-1.414}{0.385}{-1.415}{0.382}{0.414}
\SG{-1.415}{0.382}{-1.416}{0.379}{0.413}
\SG{-1.416}{0.379}{-1.417}{0.377}{0.412}
\SG{-1.417}{0.377}{-1.419}{0.374}{0.409}
\SG{-1.419}{0.374}{-1.420}{0.371}{0.403}
\SG{-1.420}{0.371}{-1.421}{0.369}{0.397}
\SG{-1.421}{0.369}{-1.422}{0.366}{0.392}
\SG{-1.422}{0.366}{-1.423}{0.363}{0.388}
\SG{-1.423}{0.363}{-1.425}{0.361}{0.385}
\SG{-1.425}{0.361}{-1.426}{0.358}{0.384}
\SG{-1.426}{0.358}{-1.427}{0.355}{0.385}
\SG{-1.427}{0.355}{-1.428}{0.352}{0.385}
\SG{-1.428}{0.352}{-1.429}{0.350}{0.385}
\SG{-1.429}{0.350}{-1.430}{0.347}{0.385}
\SG{-1.430}{0.347}{-1.432}{0.344}{0.387}
\SG{-1.432}{0.344}{-1.433}{0.341}{0.389}
\SG{-1.433}{0.341}{-1.434}{0.339}{0.392}
\SG{-1.434}{0.339}{-1.435}{0.336}{0.394}
\SG{-1.435}{0.336}{-1.436}{0.333}{0.396}
\SG{-1.436}{0.333}{-1.437}{0.330}{0.397}
\SG{-1.437}{0.330}{-1.438}{0.328}{0.397}
\SG{-1.438}{0.328}{-1.439}{0.325}{0.398}
\SG{-1.439}{0.325}{-1.440}{0.322}{0.399}
\SG{-1.440}{0.322}{-1.441}{0.319}{0.400}
\SG{-1.441}{0.319}{-1.442}{0.317}{0.401}
\SG{-1.442}{0.317}{-1.443}{0.314}{0.403}
\SG{-1.443}{0.314}{-1.444}{0.311}{0.405}
\SG{-1.444}{0.311}{-1.445}{0.308}{0.407}
\SG{-1.445}{0.308}{-1.447}{0.306}{0.409}
\SG{-1.447}{0.306}{-1.448}{0.303}{0.410}
\SG{-1.448}{0.303}{-1.448}{0.300}{0.407}
\SG{-1.448}{0.300}{-1.449}{0.297}{0.403}
\SG{-1.449}{0.297}{-1.450}{0.294}{0.399}
\SG{-1.450}{0.294}{-1.451}{0.292}{0.397}
\SG{-1.451}{0.292}{-1.452}{0.289}{0.394}
\SG{-1.452}{0.289}{-1.453}{0.286}{0.394}
\SG{-1.453}{0.286}{-1.454}{0.283}{0.395}
\SG{-1.454}{0.283}{-1.455}{0.280}{0.396}
\SG{-1.455}{0.280}{-1.456}{0.278}{0.397}
\SG{-1.456}{0.278}{-1.457}{0.275}{0.398}
\SG{-1.457}{0.275}{-1.458}{0.272}{0.400}
\SG{-1.458}{0.272}{-1.459}{0.269}{0.403}
\SG{-1.459}{0.269}{-1.460}{0.266}{0.406}
\SG{-1.460}{0.266}{-1.461}{0.264}{0.409}
\SG{-1.461}{0.264}{-1.461}{0.261}{0.411}
\SG{-1.461}{0.261}{-1.462}{0.258}{0.412}
\SG{-1.462}{0.258}{-1.463}{0.255}{0.409}
\SG{-1.463}{0.255}{-1.464}{0.252}{0.407}
\SG{-1.464}{0.252}{-1.465}{0.249}{0.405}
\SG{-1.465}{0.249}{-1.466}{0.247}{0.403}
\SG{-1.466}{0.247}{-1.466}{0.244}{0.402}
\SG{-1.466}{0.244}{-1.467}{0.241}{0.404}
\SG{-1.467}{0.241}{-1.468}{0.238}{0.406}
\SG{-1.468}{0.238}{-1.469}{0.235}{0.408}
\SG{-1.469}{0.235}{-1.469}{0.232}{0.410}
\SG{-1.469}{0.232}{-1.470}{0.229}{0.412}
\SG{-1.470}{0.229}{-1.471}{0.227}{0.418}
\SG{-1.471}{0.227}{-1.472}{0.224}{0.422}
\SG{-1.472}{0.224}{-1.472}{0.221}{0.426}
\SG{-1.472}{0.221}{-1.473}{0.218}{0.429}
\SG{-1.473}{0.218}{-1.474}{0.215}{0.432}
\SG{-1.474}{0.215}{-1.475}{0.212}{0.428}
\SG{-1.475}{0.212}{-1.475}{0.209}{0.423}
\SG{-1.475}{0.209}{-1.476}{0.207}{0.419}
\SG{-1.476}{0.207}{-1.477}{0.204}{0.414}
\SG{-1.477}{0.204}{-1.477}{0.201}{0.411}
\SG{-1.477}{0.201}{-1.478}{0.198}{0.406}
\SG{-1.478}{0.198}{-1.479}{0.195}{0.402}
\SG{-1.479}{0.195}{-1.479}{0.192}{0.399}
\SG{-1.479}{0.192}{-1.480}{0.189}{0.395}
\SG{-1.480}{0.189}{-1.480}{0.186}{0.392}
\SG{-1.480}{0.186}{-1.481}{0.183}{0.391}
\SG{-1.481}{0.183}{-1.482}{0.181}{0.390}
\SG{-1.482}{0.181}{-1.482}{0.178}{0.389}
\SG{-1.482}{0.178}{-1.483}{0.175}{0.389}
\SG{-1.483}{0.175}{-1.483}{0.172}{0.389}
\SG{-1.483}{0.172}{-1.484}{0.169}{0.392}
\SG{-1.484}{0.169}{-1.485}{0.166}{0.396}
\SG{-1.485}{0.166}{-1.485}{0.163}{0.400}
\SG{-1.485}{0.163}{-1.486}{0.160}{0.403}
\SG{-1.486}{0.160}{-1.486}{0.157}{0.407}
\SG{-1.486}{0.157}{-1.487}{0.154}{0.413}
\SG{-1.487}{0.154}{-1.487}{0.151}{0.419}
\SG{-1.487}{0.151}{-1.488}{0.149}{0.425}
\SG{-1.488}{0.149}{-1.488}{0.146}{0.430}
\SG{-1.488}{0.146}{-1.489}{0.143}{0.436}
\SG{-1.489}{0.143}{-1.489}{0.140}{0.434}
\SG{-1.489}{0.140}{-1.490}{0.137}{0.432}
\SG{-1.490}{0.137}{-1.490}{0.134}{0.430}
\SG{-1.490}{0.134}{-1.490}{0.131}{0.428}
\SG{-1.490}{0.131}{-1.491}{0.128}{0.426}
\SG{-1.491}{0.128}{-1.491}{0.125}{0.423}
\SG{-1.491}{0.125}{-1.492}{0.122}{0.420}
\SG{-1.492}{0.122}{-1.492}{0.119}{0.417}
\SG{-1.492}{0.119}{-1.492}{0.116}{0.414}
\SG{-1.492}{0.116}{-1.493}{0.113}{0.412}
\SG{-1.493}{0.113}{-1.493}{0.111}{0.409}
\SG{-1.493}{0.111}{-1.494}{0.108}{0.407}
\SG{-1.494}{0.108}{-1.494}{0.105}{0.404}
\SG{-1.494}{0.105}{-1.494}{0.102}{0.402}
\SG{-1.494}{0.102}{-1.495}{0.099}{0.399}
\SG{-1.495}{0.099}{-1.495}{0.096}{0.397}
\SG{-1.495}{0.096}{-1.495}{0.093}{0.395}
\SG{-1.495}{0.093}{-1.495}{0.090}{0.393}
\SG{-1.495}{0.090}{-1.496}{0.087}{0.391}
\SG{-1.496}{0.087}{-1.496}{0.084}{0.389}
\SG{-1.496}{0.084}{-1.496}{0.081}{0.388}
\SG{-1.496}{0.081}{-1.497}{0.078}{0.386}
\SG{-1.497}{0.078}{-1.497}{0.075}{0.384}
\SG{-1.497}{0.075}{-1.497}{0.072}{0.383}
\SG{-1.497}{0.072}{-1.497}{0.069}{0.381}
\SG{-1.497}{0.069}{-1.498}{0.066}{0.380}
\SG{-1.498}{0.066}{-1.498}{0.063}{0.379}
\SG{-1.498}{0.063}{-1.498}{0.061}{0.378}
\SG{-1.498}{0.061}{-1.498}{0.058}{0.377}
\SG{-1.498}{0.058}{-1.498}{0.055}{0.375}
\SG{-1.498}{0.055}{-1.498}{0.052}{0.375}
\SG{-1.498}{0.052}{-1.499}{0.049}{0.374}
\SG{-1.499}{0.049}{-1.499}{0.046}{0.373}
\SG{-1.499}{0.046}{-1.499}{0.043}{0.372}
\SG{-1.499}{0.043}{-1.499}{0.040}{0.371}
\SG{-1.499}{0.040}{-1.499}{0.037}{0.370}
\SG{-1.499}{0.037}{-1.499}{0.034}{0.370}
\SG{-1.499}{0.034}{-1.499}{0.031}{0.369}
\SG{-1.499}{0.031}{-1.500}{0.028}{0.369}
\SG{-1.500}{0.028}{-1.500}{0.025}{0.368}
\SG{-1.500}{0.025}{-1.500}{0.022}{0.368}
\SG{-1.500}{0.022}{-1.500}{0.019}{0.367}
\SG{-1.500}{0.019}{-1.500}{0.016}{0.367}
\SG{-1.500}{0.016}{-1.500}{0.013}{0.367}
\SG{-1.500}{0.013}{-1.500}{0.010}{0.366}
\SG{-1.500}{0.010}{-1.500}{0.007}{0.366}
\SG{-1.500}{0.007}{-1.500}{0.004}{0.366}
\SG{-1.500}{0.004}{-1.500}{0.001}{0.366}
\SG{-1.500}{0.001}{-1.500}{-0.001}{0.366}
\SG{-1.500}{-0.001}{-1.500}{-0.004}{0.366}
\SG{-1.500}{-0.004}{-1.500}{-0.007}{0.366}
\SG{-1.500}{-0.007}{-1.500}{-0.010}{0.366}
\SG{-1.500}{-0.010}{-1.500}{-0.013}{0.366}
\SG{-1.500}{-0.013}{-1.500}{-0.016}{0.366}
\SG{-1.500}{-0.016}{-1.500}{-0.019}{0.367}
\SG{-1.500}{-0.019}{-1.500}{-0.022}{0.367}
\SG{-1.500}{-0.022}{-1.500}{-0.025}{0.367}
\SG{-1.500}{-0.025}{-1.500}{-0.028}{0.368}
\SG{-1.500}{-0.028}{-1.499}{-0.031}{0.368}
\SG{-1.499}{-0.031}{-1.499}{-0.034}{0.369}
\SG{-1.499}{-0.034}{-1.499}{-0.037}{0.369}
\SG{-1.499}{-0.037}{-1.499}{-0.040}{0.370}
\SG{-1.499}{-0.040}{-1.499}{-0.043}{0.370}
\SG{-1.499}{-0.043}{-1.499}{-0.046}{0.371}
\SG{-1.499}{-0.046}{-1.499}{-0.049}{0.372}
\SG{-1.499}{-0.049}{-1.498}{-0.052}{0.373}
\SG{-1.498}{-0.052}{-1.498}{-0.055}{0.374}
\SG{-1.498}{-0.055}{-1.498}{-0.058}{0.374}
\SG{-1.498}{-0.058}{-1.498}{-0.061}{0.375}
\SG{-1.498}{-0.061}{-1.498}{-0.063}{0.376}
\SG{-1.498}{-0.063}{-1.498}{-0.066}{0.378}
\SG{-1.498}{-0.066}{-1.497}{-0.069}{0.379}
\SG{-1.497}{-0.069}{-1.497}{-0.072}{0.380}
\SG{-1.497}{-0.072}{-1.497}{-0.075}{0.381}
\SG{-1.497}{-0.075}{-1.497}{-0.078}{0.383}
\SG{-1.497}{-0.078}{-1.496}{-0.081}{0.384}
\SG{-1.496}{-0.081}{-1.496}{-0.084}{0.386}
\SG{-1.496}{-0.084}{-1.496}{-0.087}{0.388}
\SG{-1.496}{-0.087}{-1.495}{-0.090}{0.389}
\SG{-1.495}{-0.090}{-1.495}{-0.093}{0.391}
\SG{-1.495}{-0.093}{-1.495}{-0.096}{0.393}
\SG{-1.495}{-0.096}{-1.495}{-0.099}{0.395}
\SG{-1.495}{-0.099}{-1.494}{-0.102}{0.397}
\SG{-1.494}{-0.102}{-1.494}{-0.105}{0.399}
\SG{-1.494}{-0.105}{-1.494}{-0.108}{0.402}
\SG{-1.494}{-0.108}{-1.493}{-0.111}{0.404}
\SG{-1.493}{-0.111}{-1.493}{-0.113}{0.407}
\SG{-1.493}{-0.113}{-1.492}{-0.116}{0.409}
\SG{-1.492}{-0.116}{-1.492}{-0.119}{0.412}
\SG{-1.492}{-0.119}{-1.492}{-0.122}{0.414}
\SG{-1.492}{-0.122}{-1.491}{-0.125}{0.417}
\SG{-1.491}{-0.125}{-1.491}{-0.128}{0.420}
\SG{-1.491}{-0.128}{-1.490}{-0.131}{0.423}
\SG{-1.490}{-0.131}{-1.490}{-0.134}{0.426}
\SG{-1.490}{-0.134}{-1.490}{-0.137}{0.428}
\SG{-1.490}{-0.137}{-1.489}{-0.140}{0.430}
\SG{-1.489}{-0.140}{-1.489}{-0.143}{0.432}
\SG{-1.489}{-0.143}{-1.488}{-0.146}{0.435}
\SG{-1.488}{-0.146}{-1.488}{-0.149}{0.436}
\SG{-1.488}{-0.149}{-1.487}{-0.151}{0.430}
\SG{-1.487}{-0.151}{-1.487}{-0.154}{0.425}
\SG{-1.487}{-0.154}{-1.486}{-0.157}{0.419}
\SG{-1.486}{-0.157}{-1.486}{-0.160}{0.412}
\SG{-1.486}{-0.160}{-1.485}{-0.163}{0.407}
\SG{-1.485}{-0.163}{-1.485}{-0.166}{0.403}
\SG{-1.485}{-0.166}{-1.484}{-0.169}{0.399}
\SG{-1.484}{-0.169}{-1.483}{-0.172}{0.396}
\SG{-1.483}{-0.172}{-1.483}{-0.175}{0.392}
\SG{-1.483}{-0.175}{-1.482}{-0.178}{0.389}
\SG{-1.482}{-0.178}{-1.482}{-0.181}{0.389}
\SG{-1.482}{-0.181}{-1.481}{-0.183}{0.389}
\SG{-1.481}{-0.183}{-1.480}{-0.186}{0.390}
\SG{-1.480}{-0.186}{-1.480}{-0.189}{0.391}
\SG{-1.480}{-0.189}{-1.479}{-0.192}{0.392}
\SG{-1.479}{-0.192}{-1.479}{-0.195}{0.395}
\SG{-1.479}{-0.195}{-1.478}{-0.198}{0.399}
\SG{-1.478}{-0.198}{-1.477}{-0.201}{0.402}
\SG{-1.477}{-0.201}{-1.477}{-0.204}{0.406}
\SG{-1.477}{-0.204}{-1.476}{-0.207}{0.411}
\SG{-1.476}{-0.207}{-1.475}{-0.209}{0.414}
\SG{-1.475}{-0.209}{-1.475}{-0.212}{0.419}
\SG{-1.475}{-0.212}{-1.474}{-0.215}{0.423}
\SG{-1.474}{-0.215}{-1.473}{-0.218}{0.428}
\SG{-1.473}{-0.218}{-1.472}{-0.221}{0.432}
\SG{-1.472}{-0.221}{-1.472}{-0.224}{0.429}
\SG{-1.472}{-0.224}{-1.471}{-0.227}{0.426}
\SG{-1.471}{-0.227}{-1.470}{-0.229}{0.422}
\SG{-1.470}{-0.229}{-1.469}{-0.232}{0.418}
\SG{-1.469}{-0.232}{-1.469}{-0.235}{0.412}
\SG{-1.469}{-0.235}{-1.468}{-0.238}{0.410}
\SG{-1.468}{-0.238}{-1.467}{-0.241}{0.408}
\SG{-1.467}{-0.241}{-1.466}{-0.244}{0.406}
\SG{-1.466}{-0.244}{-1.466}{-0.247}{0.404}
\SG{-1.466}{-0.247}{-1.465}{-0.249}{0.402}
\SG{-1.465}{-0.249}{-1.464}{-0.252}{0.403}
\SG{-1.464}{-0.252}{-1.463}{-0.255}{0.405}
\SG{-1.463}{-0.255}{-1.462}{-0.258}{0.407}
\SG{-1.462}{-0.258}{-1.461}{-0.261}{0.410}
\SG{-1.461}{-0.261}{-1.461}{-0.264}{0.412}
\SG{-1.461}{-0.264}{-1.460}{-0.266}{0.411}
\SG{-1.460}{-0.266}{-1.459}{-0.269}{0.409}
\SG{-1.459}{-0.269}{-1.458}{-0.272}{0.406}
\SG{-1.458}{-0.272}{-1.457}{-0.275}{0.403}
\SG{-1.457}{-0.275}{-1.456}{-0.278}{0.400}
\SG{-1.456}{-0.278}{-1.455}{-0.280}{0.398}
\SG{-1.455}{-0.280}{-1.454}{-0.283}{0.397}
\SG{-1.454}{-0.283}{-1.453}{-0.286}{0.396}
\SG{-1.453}{-0.286}{-1.452}{-0.289}{0.395}
\SG{-1.452}{-0.289}{-1.451}{-0.292}{0.394}
\SG{-1.451}{-0.292}{-1.450}{-0.294}{0.394}
\SG{-1.450}{-0.294}{-1.449}{-0.297}{0.397}
\SG{-1.449}{-0.297}{-1.448}{-0.300}{0.399}
\SG{-1.448}{-0.300}{-1.448}{-0.303}{0.403}
\SG{-1.448}{-0.303}{-1.447}{-0.306}{0.407}
\SG{-1.447}{-0.306}{-1.445}{-0.308}{0.410}
\SG{-1.445}{-0.308}{-1.444}{-0.311}{0.409}
\SG{-1.444}{-0.311}{-1.443}{-0.314}{0.407}
\SG{-1.443}{-0.314}{-1.442}{-0.317}{0.405}
\SG{-1.442}{-0.317}{-1.441}{-0.319}{0.403}
\SG{-1.441}{-0.319}{-1.440}{-0.322}{0.401}
\SG{-1.440}{-0.322}{-1.439}{-0.325}{0.400}
\SG{-1.439}{-0.325}{-1.438}{-0.328}{0.399}
\SG{-1.438}{-0.328}{-1.437}{-0.330}{0.398}
\SG{-1.437}{-0.330}{-1.436}{-0.333}{0.397}
\SG{-1.436}{-0.333}{-1.435}{-0.336}{0.397}
\SG{-1.435}{-0.336}{-1.434}{-0.339}{0.396}
\SG{-1.434}{-0.339}{-1.433}{-0.341}{0.394}
\SG{-1.433}{-0.341}{-1.432}{-0.344}{0.392}
\SG{-1.432}{-0.344}{-1.430}{-0.347}{0.389}
\SG{-1.430}{-0.347}{-1.429}{-0.350}{0.386}
\SG{-1.429}{-0.350}{-1.428}{-0.352}{0.385}
\SG{-1.428}{-0.352}{-1.427}{-0.355}{0.385}
\SG{-1.427}{-0.355}{-1.426}{-0.358}{0.385}
\SG{-1.426}{-0.358}{-1.425}{-0.361}{0.385}
\SG{-1.425}{-0.361}{-1.423}{-0.363}{0.384}
\SG{-1.423}{-0.363}{-1.422}{-0.366}{0.385}
\SG{-1.422}{-0.366}{-1.421}{-0.369}{0.388}
\SG{-1.421}{-0.369}{-1.420}{-0.371}{0.392}
\SG{-1.420}{-0.371}{-1.419}{-0.374}{0.397}
\SG{-1.419}{-0.374}{-1.417}{-0.377}{0.403}
\SG{-1.417}{-0.377}{-1.416}{-0.379}{0.410}
\SG{-1.416}{-0.379}{-1.415}{-0.382}{0.412}
\SG{-1.415}{-0.382}{-1.414}{-0.385}{0.413}
\SG{-1.414}{-0.385}{-1.412}{-0.387}{0.414}
\SG{-1.412}{-0.387}{-1.411}{-0.390}{0.415}
\SG{-1.411}{-0.390}{-1.410}{-0.393}{0.415}
\SG{-1.410}{-0.393}{-1.408}{-0.395}{0.414}
\SG{-1.408}{-0.395}{-1.407}{-0.398}{0.416}
\SG{-1.407}{-0.398}{-1.406}{-0.401}{0.418}
\SG{-1.406}{-0.401}{-1.405}{-0.403}{0.420}
\SG{-1.405}{-0.403}{-1.403}{-0.406}{0.423}
\SG{-1.403}{-0.406}{-1.402}{-0.409}{0.426}
\SG{-1.402}{-0.409}{-1.401}{-0.411}{0.425}
\SG{-1.401}{-0.411}{-1.399}{-0.414}{0.424}
\SG{-1.399}{-0.414}{-1.398}{-0.416}{0.423}
\SG{-1.398}{-0.416}{-1.396}{-0.419}{0.422}
\SG{-1.396}{-0.419}{-1.395}{-0.422}{0.420}
\SG{-1.395}{-0.422}{-1.394}{-0.424}{0.417}
\SG{-1.394}{-0.424}{-1.392}{-0.427}{0.414}
\SG{-1.392}{-0.427}{-1.391}{-0.430}{0.411}
\SG{-1.391}{-0.430}{-1.389}{-0.432}{0.407}
\SG{-1.389}{-0.432}{-1.388}{-0.435}{0.403}
\SG{-1.388}{-0.435}{-1.387}{-0.437}{0.398}
\SG{-1.387}{-0.437}{-1.385}{-0.440}{0.395}
\SG{-1.385}{-0.440}{-1.384}{-0.442}{0.394}
\SG{-1.384}{-0.442}{-1.382}{-0.445}{0.393}
\SG{-1.382}{-0.445}{-1.381}{-0.448}{0.391}
\SG{-1.381}{-0.448}{-1.379}{-0.450}{0.389}
\SG{-1.379}{-0.450}{-1.378}{-0.453}{0.389}
\SG{-1.378}{-0.453}{-1.376}{-0.455}{0.389}
\SG{-1.376}{-0.455}{-1.375}{-0.458}{0.388}
\SG{-1.375}{-0.458}{-1.373}{-0.460}{0.388}
\SG{-1.373}{-0.460}{-1.372}{-0.463}{0.388}
\SG{-1.372}{-0.463}{-1.370}{-0.465}{0.388}
\SG{-1.370}{-0.465}{-1.369}{-0.468}{0.388}
\SG{-1.369}{-0.468}{-1.367}{-0.470}{0.390}
\SG{-1.367}{-0.470}{-1.366}{-0.473}{0.392}
\SG{-1.366}{-0.473}{-1.364}{-0.475}{0.394}
\SG{-1.364}{-0.475}{-1.363}{-0.478}{0.398}
\SG{-1.363}{-0.478}{-1.361}{-0.480}{0.402}
\SG{-1.361}{-0.480}{-1.359}{-0.483}{0.405}
\SG{-1.359}{-0.483}{-1.358}{-0.485}{0.406}
\SG{-1.358}{-0.485}{-1.356}{-0.488}{0.407}
\SG{-1.356}{-0.488}{-1.355}{-0.490}{0.408}
\SG{-1.355}{-0.490}{-1.353}{-0.493}{0.410}
\SG{-1.353}{-0.493}{-1.351}{-0.495}{0.412}
\SG{-1.351}{-0.495}{-1.350}{-0.498}{0.412}
\SG{-1.350}{-0.498}{-1.348}{-0.500}{0.411}
\SG{-1.348}{-0.500}{-1.346}{-0.503}{0.410}
\SG{-1.346}{-0.503}{-1.345}{-0.505}{0.408}
\SG{-1.345}{-0.505}{-1.343}{-0.508}{0.406}
\SG{-1.343}{-0.508}{-1.341}{-0.510}{0.403}
\SG{-1.341}{-0.510}{-1.340}{-0.513}{0.401}
\SG{-1.340}{-0.513}{-1.338}{-0.515}{0.399}
\SG{-1.338}{-0.515}{-1.336}{-0.517}{0.396}
\SG{-1.336}{-0.517}{-1.335}{-0.520}{0.392}
\SG{-1.335}{-0.520}{-1.333}{-0.522}{0.389}
\SG{-1.333}{-0.522}{-1.331}{-0.525}{0.385}
\SG{-1.331}{-0.525}{-1.330}{-0.527}{0.384}
\SG{-1.330}{-0.527}{-1.328}{-0.529}{0.383}
\SG{-1.328}{-0.529}{-1.326}{-0.532}{0.383}
\SG{-1.326}{-0.532}{-1.324}{-0.534}{0.383}
\SG{-1.324}{-0.534}{-1.323}{-0.537}{0.384}
\SG{-1.323}{-0.537}{-1.321}{-0.539}{0.385}
\SG{-1.321}{-0.539}{-1.319}{-0.541}{0.387}
\SG{-1.319}{-0.541}{-1.317}{-0.544}{0.388}
\SG{-1.317}{-0.544}{-1.315}{-0.546}{0.389}
\SG{-1.315}{-0.546}{-1.314}{-0.548}{0.391}
\SG{-1.314}{-0.548}{-1.312}{-0.551}{0.393}
\SG{-1.312}{-0.551}{-1.310}{-0.553}{0.396}
\SG{-1.310}{-0.553}{-1.308}{-0.555}{0.398}
\SG{-1.308}{-0.555}{-1.306}{-0.558}{0.400}
\SG{-1.306}{-0.558}{-1.305}{-0.560}{0.401}
\SG{-1.305}{-0.560}{-1.303}{-0.562}{0.402}
\SG{-1.303}{-0.562}{-1.301}{-0.565}{0.402}
\SG{-1.301}{-0.565}{-1.299}{-0.567}{0.402}
\SG{-1.299}{-0.567}{-1.297}{-0.569}{0.402}
\SG{-1.297}{-0.569}{-1.295}{-0.571}{0.401}
\SG{-1.295}{-0.571}{-1.293}{-0.574}{0.400}
\SG{-1.293}{-0.574}{-1.291}{-0.576}{0.399}
\SG{-1.291}{-0.576}{-1.290}{-0.578}{0.398}
\SG{-1.290}{-0.578}{-1.288}{-0.581}{0.396}
\SG{-1.288}{-0.581}{-1.286}{-0.583}{0.394}
\SG{-1.286}{-0.583}{-1.284}{-0.585}{0.393}
\SG{-1.284}{-0.585}{-1.282}{-0.587}{0.392}
\SG{-1.282}{-0.587}{-1.280}{-0.590}{0.390}
\SG{-1.280}{-0.590}{-1.278}{-0.592}{0.388}
\SG{-1.278}{-0.592}{-1.276}{-0.594}{0.386}
\SG{-1.276}{-0.594}{-1.274}{-0.596}{0.384}
\SG{-1.274}{-0.596}{-1.272}{-0.598}{0.382}
\SG{-1.272}{-0.598}{-1.270}{-0.601}{0.381}
\SG{-1.270}{-0.601}{-1.268}{-0.603}{0.380}
\SG{-1.268}{-0.603}{-1.266}{-0.605}{0.379}
\SG{-1.266}{-0.605}{-1.264}{-0.607}{0.377}
\SG{-1.264}{-0.607}{-1.262}{-0.609}{0.376}
\SG{-1.262}{-0.609}{-1.260}{-0.612}{0.375}
\SG{-1.260}{-0.612}{-1.258}{-0.614}{0.374}
\SG{-1.258}{-0.614}{-1.256}{-0.616}{0.373}
\SG{-1.256}{-0.616}{-1.254}{-0.618}{0.372}
\SG{-1.254}{-0.618}{-1.252}{-0.620}{0.371}
\SG{-1.252}{-0.620}{-1.250}{-0.622}{0.370}
\SG{-1.250}{-0.622}{-1.248}{-0.624}{0.370}
\SG{-1.248}{-0.624}{-1.246}{-0.627}{0.369}
\SG{-1.246}{-0.627}{-1.244}{-0.629}{0.368}
\SG{-1.244}{-0.629}{-1.242}{-0.631}{0.368}
\SG{-1.242}{-0.631}{-1.240}{-0.633}{0.368}
\SG{-1.240}{-0.633}{-1.238}{-0.635}{0.367}
\SG{-1.238}{-0.635}{-1.236}{-0.637}{0.367}
\SG{-1.236}{-0.637}{-1.234}{-0.639}{0.367}
\SG{-1.234}{-0.639}{-1.231}{-0.641}{0.367}
\SG{-1.231}{-0.641}{-1.229}{-0.643}{0.367}
\SG{-1.229}{-0.643}{-1.227}{-0.645}{0.367}
\SG{-1.227}{-0.645}{-1.225}{-0.647}{0.367}
\SG{-1.225}{-0.647}{-1.223}{-0.649}{0.367}
\SG{-1.223}{-0.649}{-1.221}{-0.652}{0.367}
\SG{-1.221}{-0.652}{-1.219}{-0.654}{0.368}
\SG{-1.219}{-0.654}{-1.217}{-0.656}{0.368}
\SG{-1.217}{-0.656}{-1.214}{-0.658}{0.369}
\SG{-1.214}{-0.658}{-1.212}{-0.660}{0.369}
\SG{-1.212}{-0.660}{-1.210}{-0.662}{0.370}
\SG{-1.210}{-0.662}{-1.208}{-0.664}{0.370}
\SG{-1.208}{-0.664}{-1.206}{-0.666}{0.371}
\SG{-1.206}{-0.666}{-1.204}{-0.668}{0.372}
\SG{-1.204}{-0.668}{-1.201}{-0.670}{0.373}
\SG{-1.201}{-0.670}{-1.199}{-0.672}{0.374}
\SG{-1.199}{-0.672}{-1.197}{-0.674}{0.375}
\SG{-1.197}{-0.674}{-1.195}{-0.675}{0.376}
\SG{-1.195}{-0.675}{-1.192}{-0.677}{0.377}
\SG{-1.192}{-0.677}{-1.190}{-0.679}{0.379}
\SG{-1.190}{-0.679}{-1.188}{-0.681}{0.380}
\SG{-1.188}{-0.681}{-1.186}{-0.683}{0.382}
\SG{-1.186}{-0.683}{-1.184}{-0.685}{0.384}
\SG{-1.184}{-0.685}{-1.181}{-0.687}{0.385}
\SG{-1.181}{-0.687}{-1.179}{-0.689}{0.385}
\SG{-1.179}{-0.689}{-1.177}{-0.691}{0.387}
\SG{-1.177}{-0.691}{-1.174}{-0.693}{0.388}
\SG{-1.174}{-0.693}{-1.172}{-0.695}{0.390}
\SG{-1.172}{-0.695}{-1.170}{-0.697}{0.391}
\SG{-1.170}{-0.697}{-1.168}{-0.698}{0.391}
\SG{-1.168}{-0.698}{-1.165}{-0.700}{0.391}
\SG{-1.165}{-0.700}{-1.163}{-0.702}{0.392}
\SG{-1.163}{-0.702}{-1.161}{-0.704}{0.391}
\SG{-1.161}{-0.704}{-1.158}{-0.706}{0.391}
\SG{-1.158}{-0.706}{-1.156}{-0.708}{0.390}
\SG{-1.156}{-0.708}{-1.154}{-0.709}{0.388}
\SG{-1.154}{-0.709}{-1.151}{-0.711}{0.386}
\SG{-1.151}{-0.711}{-1.149}{-0.713}{0.384}
\SG{-1.149}{-0.713}{-1.147}{-0.715}{0.381}
\SG{-1.147}{-0.715}{-1.144}{-0.717}{0.379}
\SG{-1.144}{-0.717}{-1.142}{-0.718}{0.377}
\SG{-1.142}{-0.718}{-1.140}{-0.720}{0.375}
\SG{-1.140}{-0.720}{-1.137}{-0.722}{0.373}
\SG{-1.137}{-0.722}{-1.135}{-0.724}{0.372}
\SG{-1.135}{-0.724}{-1.133}{-0.725}{0.369}
\SG{-1.133}{-0.725}{-1.130}{-0.727}{0.368}
\SG{-1.130}{-0.727}{-1.128}{-0.729}{0.368}
\SG{-1.128}{-0.729}{-1.125}{-0.731}{0.367}
\SG{-1.125}{-0.731}{-1.123}{-0.732}{0.367}
\SG{-1.123}{-0.732}{-1.121}{-0.734}{0.368}
\SG{-1.121}{-0.734}{-1.118}{-0.736}{0.371}
\SG{-1.118}{-0.736}{-1.116}{-0.738}{0.374}
\SG{-1.116}{-0.738}{-1.113}{-0.739}{0.377}
\SG{-1.113}{-0.739}{-1.111}{-0.741}{0.380}
\SG{-1.111}{-0.741}{-1.108}{-0.743}{0.382}
\SG{-1.108}{-0.743}{-1.106}{-0.744}{0.384}
\SG{-1.106}{-0.744}{-1.104}{-0.746}{0.386}
\SG{-1.104}{-0.746}{-1.101}{-0.748}{0.388}
\SG{-1.101}{-0.748}{-1.099}{-0.749}{0.390}
\SG{-1.099}{-0.749}{-1.096}{-0.751}{0.391}
\SG{-1.096}{-0.751}{-1.094}{-0.752}{0.391}
\SG{-1.094}{-0.752}{-1.091}{-0.754}{0.391}
\SG{-1.091}{-0.754}{-1.089}{-0.756}{0.390}
\SG{-1.089}{-0.756}{-1.086}{-0.757}{0.388}
\SG{-1.086}{-0.757}{-1.084}{-0.759}{0.386}
\SG{-1.084}{-0.759}{-1.081}{-0.760}{0.385}
\SG{-1.081}{-0.760}{-1.079}{-0.762}{0.383}
\SG{-1.079}{-0.762}{-1.076}{-0.764}{0.381}
\SG{-1.076}{-0.764}{-1.074}{-0.765}{0.376}
\SG{-1.074}{-0.765}{-1.071}{-0.767}{0.372}
\SG{-1.071}{-0.767}{-1.069}{-0.768}{0.369}
\SG{-1.069}{-0.768}{-1.066}{-0.770}{0.367}
\SG{-1.066}{-0.770}{-1.064}{-0.771}{0.365}
\SG{-1.064}{-0.771}{-1.061}{-0.773}{0.363}
\SG{-1.061}{-0.773}{-1.059}{-0.774}{0.363}
\SG{-1.059}{-0.774}{-1.056}{-0.776}{0.363}
\SG{-1.056}{-0.776}{-1.054}{-0.777}{0.363}
\SG{-1.054}{-0.777}{-1.051}{-0.779}{0.362}
\SG{-1.051}{-0.779}{-1.048}{-0.780}{0.363}
\SG{-1.048}{-0.780}{-1.046}{-0.782}{0.363}
\SG{-1.046}{-0.782}{-1.043}{-0.783}{0.363}
\SG{-1.043}{-0.783}{-1.041}{-0.785}{0.365}
\SG{-1.041}{-0.785}{-1.038}{-0.786}{0.366}
\SG{-1.038}{-0.786}{-1.036}{-0.788}{0.367}
\SG{-1.036}{-0.788}{-1.033}{-0.789}{0.368}
\SG{-1.033}{-0.789}{-1.030}{-0.790}{0.373}
\SG{-1.030}{-0.790}{-1.028}{-0.792}{0.377}
\SG{-1.028}{-0.792}{-1.025}{-0.793}{0.380}
\SG{-1.025}{-0.793}{-1.023}{-0.795}{0.383}
\SG{-1.023}{-0.795}{-1.020}{-0.796}{0.386}
\SG{-1.020}{-0.796}{-1.017}{-0.797}{0.389}
\SG{-1.017}{-0.797}{-1.015}{-0.799}{0.390}
\SG{-1.015}{-0.799}{-1.012}{-0.800}{0.391}
\SG{-1.012}{-0.800}{-1.009}{-0.801}{0.392}
\SG{-1.009}{-0.801}{-1.007}{-0.803}{0.393}
\SG{-1.007}{-0.803}{-1.004}{-0.804}{0.394}
\SG{-1.004}{-0.804}{-1.002}{-0.805}{0.391}
\SG{-1.002}{-0.805}{-0.999}{-0.807}{0.388}
\SG{-0.999}{-0.807}{-0.996}{-0.808}{0.385}
\SG{-0.996}{-0.808}{-0.994}{-0.809}{0.383}
\SG{-0.994}{-0.809}{-0.991}{-0.811}{0.381}
\SG{-0.991}{-0.811}{-0.988}{-0.812}{0.380}
\SG{-0.988}{-0.812}{-0.986}{-0.813}{0.380}
\SG{-0.986}{-0.813}{-0.983}{-0.814}{0.379}
\SG{-0.983}{-0.814}{-0.980}{-0.816}{0.378}
\SG{-0.980}{-0.816}{-0.978}{-0.817}{0.377}
\SG{-0.978}{-0.817}{-0.975}{-0.818}{0.375}
\SG{-0.975}{-0.818}{-0.972}{-0.819}{0.368}
\SG{-0.972}{-0.819}{-0.969}{-0.821}{0.362}
\SG{-0.969}{-0.821}{-0.967}{-0.822}{0.356}
\SG{-0.967}{-0.822}{-0.964}{-0.823}{0.352}
\SG{-0.964}{-0.823}{-0.961}{-0.824}{0.348}
\SG{-0.961}{-0.824}{-0.959}{-0.825}{0.346}
\SG{-0.959}{-0.825}{-0.956}{-0.827}{0.346}
\SG{-0.956}{-0.827}{-0.953}{-0.828}{0.346}
\SG{-0.953}{-0.828}{-0.951}{-0.829}{0.346}
\SG{-0.951}{-0.829}{-0.948}{-0.830}{0.345}
\SG{-0.948}{-0.830}{-0.945}{-0.831}{0.345}
\SG{-0.945}{-0.831}{-0.942}{-0.832}{0.348}
\SG{-0.942}{-0.832}{-0.940}{-0.833}{0.350}
\SG{-0.940}{-0.833}{-0.937}{-0.835}{0.352}
\SG{-0.937}{-0.835}{-0.934}{-0.836}{0.353}
\SG{-0.934}{-0.836}{-0.931}{-0.837}{0.354}
\SG{-0.931}{-0.837}{-0.929}{-0.838}{0.355}
\SG{-0.929}{-0.838}{-0.926}{-0.839}{0.355}
\SG{-0.926}{-0.839}{-0.923}{-0.840}{0.355}
\SG{-0.923}{-0.840}{-0.920}{-0.841}{0.356}
\SG{-0.920}{-0.841}{-0.918}{-0.842}{0.356}
\SG{-0.918}{-0.842}{-0.915}{-0.843}{0.358}
\SG{-0.915}{-0.843}{-0.912}{-0.844}{0.360}
\SG{-0.912}{-0.844}{-0.909}{-0.845}{0.361}
\SG{-0.909}{-0.845}{-0.907}{-0.846}{0.363}
\SG{-0.907}{-0.846}{-0.904}{-0.847}{0.363}
\SG{-0.904}{-0.847}{-0.901}{-0.848}{0.362}
\SG{-0.901}{-0.848}{-0.898}{-0.849}{0.357}
\SG{-0.898}{-0.849}{-0.895}{-0.850}{0.353}
\SG{-0.895}{-0.850}{-0.893}{-0.851}{0.350}
\SG{-0.893}{-0.851}{-0.890}{-0.852}{0.347}
\SG{-0.890}{-0.852}{-0.887}{-0.853}{0.345}
\SG{-0.887}{-0.853}{-0.884}{-0.854}{0.345}
\SG{-0.884}{-0.854}{-0.881}{-0.855}{0.346}
\SG{-0.881}{-0.855}{-0.879}{-0.856}{0.347}
\SG{-0.879}{-0.856}{-0.876}{-0.857}{0.348}
\SG{-0.876}{-0.857}{-0.873}{-0.858}{0.349}
\SG{-0.873}{-0.858}{-0.870}{-0.858}{0.352}
\SG{-0.870}{-0.858}{-0.867}{-0.859}{0.354}
\SG{-0.867}{-0.859}{-0.864}{-0.860}{0.356}
\SG{-0.864}{-0.860}{-0.862}{-0.861}{0.358}
\SG{-0.862}{-0.861}{-0.859}{-0.862}{0.360}
\SG{-0.859}{-0.862}{-0.856}{-0.863}{0.357}
\SG{-0.856}{-0.863}{-0.853}{-0.864}{0.354}
\SG{-0.853}{-0.864}{-0.850}{-0.864}{0.352}
\SG{-0.850}{-0.864}{-0.847}{-0.865}{0.350}
\SG{-0.847}{-0.865}{-0.845}{-0.866}{0.348}
\SG{-0.845}{-0.866}{-0.842}{-0.867}{0.349}
\SG{-0.842}{-0.867}{-0.839}{-0.868}{0.350}
\SG{-0.839}{-0.868}{-0.836}{-0.868}{0.352}
\SG{-0.836}{-0.868}{-0.833}{-0.869}{0.353}
\SG{-0.833}{-0.869}{-0.830}{-0.870}{0.355}
\SG{-0.830}{-0.870}{-0.828}{-0.871}{0.359}
\SG{-0.828}{-0.871}{-0.825}{-0.871}{0.364}
\SG{-0.825}{-0.871}{-0.822}{-0.872}{0.367}
\SG{-0.822}{-0.872}{-0.819}{-0.873}{0.370}
\SG{-0.819}{-0.873}{-0.816}{-0.874}{0.373}
\SG{-0.816}{-0.874}{-0.813}{-0.874}{0.370}
\SG{-0.813}{-0.874}{-0.810}{-0.875}{0.365}
\SG{-0.810}{-0.875}{-0.807}{-0.876}{0.360}
\SG{-0.807}{-0.876}{-0.805}{-0.876}{0.355}
\SG{-0.805}{-0.876}{-0.802}{-0.877}{0.351}
\SG{-0.802}{-0.877}{-0.799}{-0.878}{0.346}
\SG{-0.799}{-0.878}{-0.796}{-0.878}{0.342}
\SG{-0.796}{-0.878}{-0.793}{-0.879}{0.338}
\SG{-0.793}{-0.879}{-0.790}{-0.880}{0.334}
\SG{-0.790}{-0.880}{-0.787}{-0.880}{0.330}
\SG{-0.787}{-0.880}{-0.784}{-0.881}{0.327}
\SG{-0.784}{-0.881}{-0.781}{-0.881}{0.326}
\SG{-0.781}{-0.881}{-0.779}{-0.882}{0.325}
\SG{-0.779}{-0.882}{-0.776}{-0.883}{0.324}
\SG{-0.776}{-0.883}{-0.773}{-0.883}{0.323}
\SG{-0.773}{-0.883}{-0.770}{-0.884}{0.325}
\SG{-0.770}{-0.884}{-0.767}{-0.884}{0.328}
\SG{-0.767}{-0.884}{-0.764}{-0.885}{0.331}
\SG{-0.764}{-0.885}{-0.761}{-0.885}{0.335}
\SG{-0.761}{-0.885}{-0.758}{-0.886}{0.338}
\SG{-0.758}{-0.886}{-0.755}{-0.886}{0.342}
\SG{-0.755}{-0.886}{-0.752}{-0.887}{0.348}
\SG{-0.752}{-0.887}{-0.750}{-0.887}{0.354}
\SG{-0.750}{-0.887}{-0.747}{-0.888}{0.359}
\SG{-0.747}{-0.888}{-0.744}{-0.888}{0.364}
\SG{-0.744}{-0.888}{-0.741}{-0.889}{0.365}
\SG{-0.741}{-0.889}{-0.738}{-0.889}{0.362}
\SG{-0.738}{-0.889}{-0.735}{-0.890}{0.359}
\SG{-0.735}{-0.890}{-0.732}{-0.890}{0.357}
\SG{-0.732}{-0.890}{-0.729}{-0.891}{0.354}
\SG{-0.729}{-0.891}{-0.726}{-0.891}{0.351}
\SG{-0.726}{-0.891}{-0.723}{-0.892}{0.348}
\SG{-0.723}{-0.892}{-0.720}{-0.892}{0.344}
\SG{-0.720}{-0.892}{-0.717}{-0.892}{0.341}
\SG{-0.717}{-0.892}{-0.714}{-0.893}{0.338}
\SG{-0.714}{-0.893}{-0.711}{-0.893}{0.334}
\SG{-0.711}{-0.893}{-0.709}{-0.893}{0.331}
\SG{-0.709}{-0.893}{-0.706}{-0.894}{0.328}
\SG{-0.706}{-0.894}{-0.703}{-0.894}{0.325}
\SG{-0.703}{-0.894}{-0.700}{-0.894}{0.322}
\SG{-0.700}{-0.894}{-0.697}{-0.895}{0.320}
\SG{-0.697}{-0.895}{-0.694}{-0.895}{0.317}
\SG{-0.694}{-0.895}{-0.691}{-0.895}{0.315}
\SG{-0.691}{-0.895}{-0.688}{-0.896}{0.312}
\SG{-0.688}{-0.896}{-0.685}{-0.896}{0.310}
\SG{-0.685}{-0.896}{-0.682}{-0.896}{0.307}
\SG{-0.682}{-0.896}{-0.679}{-0.896}{0.305}
\SG{-0.679}{-0.896}{-0.676}{-0.897}{0.303}
\SG{-0.676}{-0.897}{-0.673}{-0.897}{0.301}
\SG{-0.673}{-0.897}{-0.670}{-0.897}{0.299}
\SG{-0.670}{-0.897}{-0.667}{-0.897}{0.297}
\SG{-0.667}{-0.897}{-0.664}{-0.898}{0.295}
\SG{-0.664}{-0.898}{-0.661}{-0.898}{0.294}
\SG{-0.661}{-0.898}{-0.659}{-0.898}{0.292}
\SG{-0.659}{-0.898}{-0.656}{-0.898}{0.290}
\SG{-0.656}{-0.898}{-0.653}{-0.898}{0.289}
\SG{-0.653}{-0.898}{-0.650}{-0.899}{0.287}
\SG{-0.650}{-0.899}{-0.647}{-0.899}{0.286}
\SG{-0.647}{-0.899}{-0.644}{-0.899}{0.284}
\SG{-0.644}{-0.899}{-0.641}{-0.899}{0.283}
\SG{-0.641}{-0.899}{-0.638}{-0.899}{0.282}
\SG{-0.638}{-0.899}{-0.635}{-0.899}{0.281}
\SG{-0.635}{-0.899}{-0.632}{-0.899}{0.280}
\SG{-0.632}{-0.899}{-0.629}{-0.900}{0.278}
\SG{-0.629}{-0.900}{-0.626}{-0.900}{0.277}
\SG{-0.626}{-0.900}{-0.623}{-0.900}{0.276}
\SG{-0.623}{-0.900}{-0.620}{-0.900}{0.275}
\SG{-0.620}{-0.900}{-0.617}{-0.900}{0.274}
\SG{-0.617}{-0.900}{-0.614}{-0.900}{0.273}
\SG{-0.614}{-0.900}{-0.611}{-0.900}{0.273}
\SG{-0.611}{-0.900}{-0.608}{-0.900}{0.272}
\SG{-0.608}{-0.900}{-0.605}{-0.900}{0.271}
\SG{-0.605}{-0.900}{-0.602}{-0.900}{0.270}
\SG{-0.602}{-0.900}{-0.599}{-0.900}{0.270}
\SG{-0.599}{-0.900}{-0.597}{-0.900}{0.269}
\SG{-0.597}{-0.900}{-0.594}{-0.900}{0.268}
\SG{-0.594}{-0.900}{-0.591}{-0.900}{0.268}
\SG{-0.591}{-0.900}{-0.588}{-0.900}{0.267}
\SG{-0.588}{-0.900}{-0.585}{-0.900}{0.267}
\SG{-0.585}{-0.900}{-0.582}{-0.900}{0.267}
\SG{-0.582}{-0.900}{-0.579}{-0.900}{0.266}
\SG{-0.579}{-0.900}{-0.576}{-0.900}{0.266}
\SG{-0.576}{-0.900}{-0.573}{-0.900}{0.266}
\SG{-0.573}{-0.900}{-0.570}{-0.899}{0.265}
\SG{-0.570}{-0.899}{-0.567}{-0.899}{0.265}
\SG{-0.567}{-0.899}{-0.564}{-0.899}{0.265}
\SG{-0.564}{-0.899}{-0.561}{-0.899}{0.265}
\SG{-0.561}{-0.899}{-0.558}{-0.899}{0.265}
\SG{-0.558}{-0.899}{-0.555}{-0.899}{0.265}
\SG{-0.555}{-0.899}{-0.552}{-0.899}{0.265}
\SG{-0.552}{-0.899}{-0.549}{-0.899}{0.265}
\SG{-0.549}{-0.899}{-0.546}{-0.898}{0.265}
\SG{-0.546}{-0.898}{-0.543}{-0.898}{0.266}
\SG{-0.543}{-0.898}{-0.540}{-0.898}{0.266}
\SG{-0.540}{-0.898}{-0.537}{-0.898}{0.266}
\SG{-0.537}{-0.898}{-0.535}{-0.898}{0.267}
\SG{-0.535}{-0.898}{-0.532}{-0.897}{0.267}
\SG{-0.532}{-0.897}{-0.529}{-0.897}{0.268}
\SG{-0.529}{-0.897}{-0.526}{-0.897}{0.268}
\SG{-0.526}{-0.897}{-0.523}{-0.897}{0.269}
\SG{-0.523}{-0.897}{-0.520}{-0.896}{0.270}
\SG{-0.520}{-0.896}{-0.517}{-0.896}{0.270}
\SG{-0.517}{-0.896}{-0.514}{-0.896}{0.271}
\SG{-0.514}{-0.896}{-0.511}{-0.896}{0.272}
\SG{-0.511}{-0.896}{-0.508}{-0.895}{0.273}
\SG{-0.508}{-0.895}{-0.505}{-0.895}{0.274}
\SG{-0.505}{-0.895}{-0.502}{-0.895}{0.276}
\SG{-0.502}{-0.895}{-0.499}{-0.894}{0.277}
\SG{-0.499}{-0.894}{-0.496}{-0.894}{0.278}
\SG{-0.496}{-0.894}{-0.493}{-0.894}{0.280}
\SG{-0.493}{-0.894}{-0.490}{-0.893}{0.281}
\SG{-0.490}{-0.893}{-0.487}{-0.893}{0.283}
\SG{-0.487}{-0.893}{-0.485}{-0.893}{0.285}
\SG{-0.485}{-0.893}{-0.482}{-0.892}{0.287}
\SG{-0.482}{-0.892}{-0.479}{-0.892}{0.288}
\SG{-0.479}{-0.892}{-0.476}{-0.891}{0.290}
\SG{-0.476}{-0.891}{-0.473}{-0.891}{0.292}
\SG{-0.473}{-0.891}{-0.470}{-0.891}{0.295}
\SG{-0.470}{-0.891}{-0.467}{-0.890}{0.297}
\SG{-0.467}{-0.890}{-0.464}{-0.890}{0.298}
\SG{-0.464}{-0.890}{-0.461}{-0.889}{0.299}
\SG{-0.461}{-0.889}{-0.458}{-0.889}{0.300}
\SG{-0.458}{-0.889}{-0.455}{-0.888}{0.302}
\SG{-0.455}{-0.888}{-0.452}{-0.888}{0.303}
\SG{-0.452}{-0.888}{-0.449}{-0.887}{0.298}
\SG{-0.449}{-0.887}{-0.447}{-0.887}{0.292}
\SG{-0.447}{-0.887}{-0.444}{-0.886}{0.285}
\SG{-0.444}{-0.886}{-0.441}{-0.886}{0.278}
\SG{-0.441}{-0.886}{-0.438}{-0.885}{0.271}
\SG{-0.438}{-0.885}{-0.435}{-0.885}{0.266}
\SG{-0.435}{-0.885}{-0.432}{-0.884}{0.262}
\SG{-0.432}{-0.884}{-0.429}{-0.884}{0.257}
\SG{-0.429}{-0.884}{-0.426}{-0.883}{0.252}
\SG{-0.426}{-0.883}{-0.423}{-0.882}{0.247}
\SG{-0.423}{-0.882}{-0.420}{-0.882}{0.246}
\SG{-0.420}{-0.882}{-0.418}{-0.881}{0.245}
\SG{-0.418}{-0.881}{-0.415}{-0.881}{0.245}
\SG{-0.415}{-0.881}{-0.412}{-0.880}{0.244}
\SG{-0.412}{-0.880}{-0.409}{-0.879}{0.244}
\SG{-0.409}{-0.879}{-0.406}{-0.879}{0.246}
\SG{-0.406}{-0.879}{-0.403}{-0.878}{0.248}
\SG{-0.403}{-0.878}{-0.400}{-0.878}{0.251}
\SG{-0.400}{-0.878}{-0.397}{-0.877}{0.253}
\SG{-0.397}{-0.877}{-0.394}{-0.876}{0.256}
\SG{-0.394}{-0.876}{-0.392}{-0.876}{0.259}
\SG{-0.392}{-0.876}{-0.389}{-0.875}{0.262}
\SG{-0.389}{-0.875}{-0.386}{-0.874}{0.266}
\SG{-0.386}{-0.874}{-0.383}{-0.873}{0.269}
\SG{-0.383}{-0.873}{-0.380}{-0.873}{0.273}
\SG{-0.380}{-0.873}{-0.377}{-0.872}{0.271}
\SG{-0.377}{-0.872}{-0.374}{-0.871}{0.266}
\SG{-0.374}{-0.871}{-0.371}{-0.870}{0.262}
\SG{-0.371}{-0.870}{-0.369}{-0.870}{0.256}
\SG{-0.369}{-0.870}{-0.366}{-0.869}{0.250}
\SG{-0.366}{-0.869}{-0.363}{-0.868}{0.246}
\SG{-0.363}{-0.868}{-0.360}{-0.867}{0.242}
\SG{-0.360}{-0.867}{-0.357}{-0.867}{0.239}
\SG{-0.357}{-0.867}{-0.354}{-0.866}{0.236}
\SG{-0.354}{-0.866}{-0.352}{-0.865}{0.232}
\SG{-0.352}{-0.865}{-0.349}{-0.864}{0.231}
\SG{-0.349}{-0.864}{-0.346}{-0.863}{0.231}
\SG{-0.346}{-0.863}{-0.343}{-0.863}{0.231}
\SG{-0.343}{-0.863}{-0.340}{-0.862}{0.232}
\SG{-0.340}{-0.862}{-0.337}{-0.861}{0.234}
\SG{-0.337}{-0.861}{-0.335}{-0.860}{0.233}
\SG{-0.335}{-0.860}{-0.332}{-0.859}{0.229}
\SG{-0.332}{-0.859}{-0.329}{-0.858}{0.225}
\SG{-0.329}{-0.858}{-0.326}{-0.857}{0.221}
\SG{-0.326}{-0.857}{-0.323}{-0.856}{0.216}
\SG{-0.323}{-0.856}{-0.320}{-0.855}{0.212}
\SG{-0.320}{-0.855}{-0.318}{-0.855}{0.209}
\SG{-0.318}{-0.855}{-0.315}{-0.854}{0.206}
\SG{-0.315}{-0.854}{-0.312}{-0.853}{0.203}
\SG{-0.312}{-0.853}{-0.309}{-0.852}{0.201}
\SG{-0.309}{-0.852}{-0.306}{-0.851}{0.199}
\SG{-0.306}{-0.851}{-0.304}{-0.850}{0.199}
\SG{-0.304}{-0.850}{-0.301}{-0.849}{0.200}
\SG{-0.301}{-0.849}{-0.298}{-0.848}{0.201}
\SG{-0.298}{-0.848}{-0.295}{-0.847}{0.203}
\SG{-0.295}{-0.847}{-0.293}{-0.846}{0.206}
\SG{-0.293}{-0.846}{-0.290}{-0.845}{0.203}
\SG{-0.290}{-0.845}{-0.287}{-0.844}{0.200}
\SG{-0.287}{-0.844}{-0.284}{-0.843}{0.197}
\SG{-0.284}{-0.843}{-0.281}{-0.842}{0.194}
\SG{-0.281}{-0.842}{-0.279}{-0.841}{0.190}
\SG{-0.279}{-0.841}{-0.276}{-0.840}{0.187}
\SG{-0.276}{-0.840}{-0.273}{-0.839}{0.185}
\SG{-0.273}{-0.839}{-0.270}{-0.837}{0.183}
\SG{-0.270}{-0.837}{-0.268}{-0.836}{0.181}
\SG{-0.268}{-0.836}{-0.265}{-0.835}{0.179}
\SG{-0.265}{-0.835}{-0.262}{-0.834}{0.179}
\SG{-0.262}{-0.834}{-0.259}{-0.833}{0.181}
\SG{-0.259}{-0.833}{-0.256}{-0.832}{0.183}
\SG{-0.256}{-0.832}{-0.251}{-0.830}{0.187}
\SG{-0.251}{-0.830}{-0.246}{-0.827}{0.193}
\SG{-0.246}{-0.827}{-0.239}{-0.825}{0.189}
\SG{-0.239}{-0.825}{-0.233}{-0.822}{0.183}
\SG{-0.233}{-0.822}{-0.227}{-0.820}{0.175}
\SG{-0.227}{-0.820}{-0.220}{-0.817}{0.171}
\SG{-0.220}{-0.817}{-0.214}{-0.814}{0.167}
\SG{-0.214}{-0.814}{-0.208}{-0.812}{0.158}
\SG{-0.208}{-0.812}{-0.201}{-0.809}{0.147}
\SG{-0.201}{-0.809}{-0.195}{-0.806}{0.138}
\SG{-0.195}{-0.806}{-0.189}{-0.804}{0.132}
\SG{-0.189}{-0.804}{-0.182}{-0.801}{0.127}
\SG{-0.182}{-0.801}{-0.176}{-0.798}{0.119}
\SG{-0.176}{-0.798}{-0.170}{-0.796}{0.110}
\SG{-0.170}{-0.796}{-0.163}{-0.793}{0.106}
\SG{-0.163}{-0.793}{-0.157}{-0.791}{0.102}
\SG{-0.157}{-0.791}{-0.151}{-0.788}{0.104}
\SG{-0.151}{-0.788}{-0.144}{-0.785}{0.112}
\SG{-0.144}{-0.785}{-0.138}{-0.783}{0.118}
\SG{-0.138}{-0.783}{-0.132}{-0.780}{0.114}
\SG{-0.132}{-0.780}{-0.125}{-0.777}{0.108}
\SG{-0.125}{-0.777}{-0.119}{-0.775}{0.106}
\SG{-0.119}{-0.775}{-0.113}{-0.772}{0.104}
\SG{-0.113}{-0.772}{-0.106}{-0.770}{0.099}
\SG{-0.106}{-0.770}{-0.100}{-0.767}{0.092}
\SG{-0.100}{-0.767}{-0.093}{-0.764}{0.085}
\SG{-0.093}{-0.764}{-0.087}{-0.762}{0.081}
\SG{-0.087}{-0.762}{-0.081}{-0.759}{0.079}
\SG{-0.081}{-0.759}{-0.074}{-0.756}{0.085}
\SG{-0.074}{-0.756}{-0.068}{-0.754}{0.095}
\SG{-0.068}{-0.754}{-0.062}{-0.751}{0.099}
\SG{-0.062}{-0.751}{-0.055}{-0.749}{0.097}
\SG{-0.055}{-0.749}{-0.049}{-0.746}{0.091}
\SG{-0.049}{-0.746}{-0.043}{-0.743}{0.089}
\SG{-0.043}{-0.743}{-0.036}{-0.741}{0.089}
\SG{-0.036}{-0.741}{-0.030}{-0.738}{0.084}
\SG{-0.030}{-0.738}{-0.024}{-0.735}{0.076}
\SG{-0.024}{-0.735}{-0.017}{-0.733}{0.068}
\SG{-0.017}{-0.733}{-0.011}{-0.730}{0.064}
\SG{-0.011}{-0.730}{-0.005}{-0.727}{0.061}
\SG{-0.005}{-0.727}{0.002}{-0.725}{0.055}
\SG{0.002}{-0.725}{0.008}{-0.722}{0.048}
\SG{0.008}{-0.722}{0.014}{-0.720}{0.045}
\SG{0.014}{-0.720}{0.021}{-0.717}{0.042}
\SG{0.021}{-0.717}{0.027}{-0.714}{0.043}
\SG{0.027}{-0.714}{0.033}{-0.712}{0.051}
\SG{0.033}{-0.712}{0.040}{-0.709}{0.062}
\SG{0.040}{-0.709}{0.046}{-0.706}{0.060}
\SG{0.046}{-0.706}{0.052}{-0.704}{0.057}
\SG{0.052}{-0.704}{0.059}{-0.701}{0.054}
\SG{0.059}{-0.701}{0.065}{-0.699}{0.053}
\SG{0.065}{-0.699}{0.071}{-0.696}{0.052}
\SG{0.071}{-0.696}{0.078}{-0.693}{0.046}
\SG{0.078}{-0.693}{0.084}{-0.691}{0.039}
\SG{0.084}{-0.691}{0.091}{-0.688}{0.037}
\SG{0.091}{-0.688}{0.097}{-0.685}{0.035}
\SG{0.097}{-0.685}{0.103}{-0.683}{0.038}
\SG{0.103}{-0.683}{0.110}{-0.680}{0.047}
\SG{0.110}{-0.680}{0.116}{-0.677}{0.058}
\SG{0.116}{-0.677}{0.122}{-0.675}{0.056}
\SG{0.122}{-0.675}{0.129}{-0.672}{0.053}
\SG{0.129}{-0.672}{0.135}{-0.670}{0.050}
\SG{0.135}{-0.670}{0.141}{-0.667}{0.050}
\SG{0.141}{-0.667}{0.148}{-0.664}{0.048}
\SG{0.148}{-0.664}{0.154}{-0.662}{0.042}
\SG{0.154}{-0.662}{0.160}{-0.659}{0.034}
\SG{0.160}{-0.659}{0.167}{-0.656}{0.030}
\SG{0.167}{-0.656}{0.173}{-0.654}{0.028}
\SG{0.173}{-0.654}{0.179}{-0.651}{0.024}
\SG{0.179}{-0.651}{0.186}{-0.649}{0.018}
\SG{0.186}{-0.649}{0.192}{-0.646}{0.013}
\SG{0.192}{-0.646}{0.198}{-0.643}{0.012}
\SG{0.198}{-0.643}{0.205}{-0.641}{0.011}
\SG{0.205}{-0.641}{0.211}{-0.638}{0.017}
\SG{0.211}{-0.638}{0.217}{-0.635}{0.028}
\SG{0.217}{-0.635}{0.224}{-0.633}{0.032}
\SG{0.224}{-0.633}{0.230}{-0.630}{0.030}
\SG{0.230}{-0.630}{0.236}{-0.627}{0.027}
\SG{0.236}{-0.627}{0.243}{-0.625}{0.026}
\SG{0.243}{-0.625}{0.249}{-0.622}{0.026}
\SG{0.249}{-0.622}{0.256}{-0.620}{0.022}
\SG{0.256}{-0.620}{0.262}{-0.617}{0.016}
\SG{0.262}{-0.617}{0.268}{-0.614}{0.013}
\SG{0.268}{-0.614}{0.275}{-0.612}{0.011}
\SG{0.275}{-0.612}{0.281}{-0.609}{0.011}
\SG{0.281}{-0.609}{0.287}{-0.606}{0.019}
\SG{0.287}{-0.606}{0.294}{-0.604}{0.031}
\SG{0.294}{-0.604}{0.300}{-0.601}{0.033}
\SG{0.300}{-0.601}{0.306}{-0.599}{0.031}
\SG{0.306}{-0.599}{0.313}{-0.596}{0.027}
\SG{0.313}{-0.596}{0.319}{-0.593}{0.027}
\SG{0.319}{-0.593}{0.325}{-0.591}{0.027}
\SG{0.325}{-0.591}{0.332}{-0.588}{0.022}
\SG{0.332}{-0.588}{0.338}{-0.585}{0.016}
\SG{0.338}{-0.585}{0.344}{-0.583}{0.012}
\SG{0.344}{-0.583}{0.351}{-0.580}{0.011}
\SG{0.351}{-0.580}{0.357}{-0.577}{0.014}
\SG{0.357}{-0.577}{0.363}{-0.575}{0.024}
\SG{0.363}{-0.575}{0.370}{-0.572}{0.038}
\SG{0.370}{-0.572}{0.376}{-0.570}{0.037}
\SG{0.376}{-0.570}{0.382}{-0.567}{0.035}
\SG{0.382}{-0.567}{0.389}{-0.564}{0.032}
\SG{0.389}{-0.564}{0.395}{-0.562}{0.032}
\SG{0.395}{-0.562}{0.401}{-0.559}{0.032}
\SG{0.401}{-0.559}{0.408}{-0.556}{0.027}
\SG{0.408}{-0.556}{0.414}{-0.554}{0.020}
\SG{0.414}{-0.554}{0.421}{-0.551}{0.015}
\SG{0.421}{-0.551}{0.427}{-0.549}{0.014}
\SG{0.427}{-0.549}{0.433}{-0.546}{0.011}
\SG{0.433}{-0.546}{0.440}{-0.543}{0.006}
\SG{0.440}{-0.543}{0.446}{-0.541}{0.001}
\SG{0.446}{-0.541}{0.452}{-0.538}{0.000}
\SG{0.452}{-0.538}{0.459}{-0.535}{0.000}
\SG{0.459}{-0.535}{0.465}{-0.533}{0.006}
\SG{0.465}{-0.533}{0.471}{-0.530}{0.017}
\SG{0.471}{-0.530}{0.478}{-0.528}{0.024}
\SG{0.478}{-0.528}{0.484}{-0.525}{0.023}
\SG{0.484}{-0.525}{0.490}{-0.522}{0.020}
\SG{0.490}{-0.522}{0.497}{-0.520}{0.020}
\SG{0.497}{-0.520}{0.503}{-0.517}{0.021}
\SG{0.503}{-0.517}{0.509}{-0.514}{0.019}
\SG{0.509}{-0.514}{0.516}{-0.512}{0.014}
\SG{0.516}{-0.512}{0.522}{-0.509}{0.010}
\SG{0.522}{-0.509}{0.528}{-0.506}{0.009}
\SG{0.528}{-0.506}{0.535}{-0.504}{0.009}
\SG{0.535}{-0.504}{0.541}{-0.501}{0.018}
\SG{0.541}{-0.501}{0.547}{-0.499}{0.030}
\SG{0.547}{-0.499}{0.554}{-0.496}{0.036}
\SG{0.554}{-0.496}{0.560}{-0.493}{0.036}
\SG{0.560}{-0.493}{0.566}{-0.491}{0.032}
\SG{0.566}{-0.491}{0.573}{-0.488}{0.032}
\SG{0.573}{-0.488}{0.579}{-0.485}{0.034}
\SG{0.579}{-0.485}{0.586}{-0.483}{0.030}
\SG{0.586}{-0.483}{0.592}{-0.480}{0.025}
\SG{0.592}{-0.480}{0.598}{-0.478}{0.019}
\SG{0.598}{-0.478}{0.605}{-0.475}{0.017}
\SG{0.605}{-0.475}{0.611}{-0.472}{0.016}
\SG{0.611}{-0.472}{0.617}{-0.470}{0.012}
\SG{0.617}{-0.470}{0.624}{-0.467}{0.007}
\SG{0.624}{-0.467}{0.630}{-0.464}{0.006}
\SG{0.630}{-0.464}{0.636}{-0.462}{0.006}
\SG{0.636}{-0.462}{0.643}{-0.459}{0.009}
\SG{0.643}{-0.459}{0.649}{-0.456}{0.019}
\SG{0.649}{-0.456}{0.655}{-0.454}{0.031}
\SG{0.655}{-0.454}{0.662}{-0.451}{0.031}
\SG{0.662}{-0.451}{0.668}{-0.449}{0.030}
\SG{0.668}{-0.449}{0.674}{-0.446}{0.028}
\SG{0.674}{-0.446}{0.681}{-0.443}{0.029}
\SG{0.681}{-0.443}{0.687}{-0.441}{0.029}
\SG{0.687}{-0.441}{0.693}{-0.438}{0.026}
\SG{0.693}{-0.438}{0.700}{-0.435}{0.020}
\SG{0.700}{-0.435}{0.706}{-0.433}{0.020}
\SG{0.706}{-0.433}{0.712}{-0.430}{0.020}
\SG{0.712}{-0.430}{0.719}{-0.428}{0.026}
\SG{0.719}{-0.428}{0.725}{-0.425}{0.037}
\SG{0.725}{-0.425}{0.731}{-0.422}{0.048}
\SG{0.731}{-0.422}{0.738}{-0.420}{0.048}
\SG{0.738}{-0.420}{0.744}{-0.417}{0.046}
\SG{0.744}{-0.417}{0.751}{-0.414}{0.045}
\SG{0.751}{-0.414}{0.757}{-0.412}{0.047}
\SG{0.757}{-0.412}{0.763}{-0.409}{0.046}
\SG{0.763}{-0.409}{0.770}{-0.406}{0.042}
\SG{0.770}{-0.406}{0.776}{-0.404}{0.036}
\SG{0.776}{-0.404}{0.782}{-0.401}{0.033}
\SG{0.782}{-0.401}{0.789}{-0.399}{0.033}
\SG{0.789}{-0.399}{0.795}{-0.396}{0.030}
\SG{0.795}{-0.396}{0.801}{-0.393}{0.026}
\SG{0.801}{-0.393}{0.808}{-0.391}{0.024}
\SG{0.808}{-0.391}{0.814}{-0.388}{0.024}
\SG{0.814}{-0.388}{0.820}{-0.385}{0.025}
\SG{0.820}{-0.385}{0.827}{-0.383}{0.034}
\SG{0.827}{-0.383}{0.833}{-0.380}{0.046}
\SG{0.833}{-0.380}{0.839}{-0.378}{0.050}
\SG{0.839}{-0.378}{0.846}{-0.375}{0.050}
\SG{0.846}{-0.375}{0.852}{-0.372}{0.048}
\SG{0.852}{-0.372}{0.858}{-0.370}{0.049}
\SG{0.858}{-0.370}{0.865}{-0.367}{0.051}
\SG{0.865}{-0.367}{0.871}{-0.364}{0.048}
\SG{0.871}{-0.364}{0.877}{-0.362}{0.044}
\SG{0.877}{-0.362}{0.884}{-0.359}{0.042}
\SG{0.884}{-0.359}{0.890}{-0.356}{0.042}
\SG{0.890}{-0.356}{0.896}{-0.354}{0.045}
\SG{0.896}{-0.354}{0.903}{-0.351}{0.055}
\SG{0.903}{-0.351}{0.909}{-0.349}{0.069}
\SG{0.909}{-0.349}{0.916}{-0.346}{0.071}
\SG{0.916}{-0.346}{0.922}{-0.343}{0.070}
\SG{0.922}{-0.343}{0.928}{-0.341}{0.068}
\SG{0.928}{-0.341}{0.935}{-0.338}{0.070}
\SG{0.935}{-0.338}{0.941}{-0.335}{0.071}
\SG{0.941}{-0.335}{0.947}{-0.333}{0.068}
\SG{0.947}{-0.333}{0.954}{-0.330}{0.064}
\SG{0.954}{-0.330}{0.960}{-0.328}{0.062}
\SG{0.960}{-0.328}{0.966}{-0.325}{0.063}
\SG{0.966}{-0.325}{0.973}{-0.322}{0.068}
\SG{0.973}{-0.322}{0.979}{-0.320}{0.080}
\SG{0.979}{-0.320}{0.985}{-0.317}{0.095}
\SG{0.985}{-0.317}{0.992}{-0.314}{0.095}
\SG{0.992}{-0.314}{0.998}{-0.312}{0.094}
\SG{0.998}{-0.312}{1.004}{-0.309}{0.094}
\SG{1.004}{-0.309}{1.011}{-0.307}{0.096}
\SG{1.011}{-0.307}{1.017}{-0.304}{0.097}
\SG{1.017}{-0.304}{1.023}{-0.301}{0.094}
\SG{1.023}{-0.301}{1.030}{-0.299}{0.087}
\SG{1.030}{-0.299}{1.036}{-0.296}{0.086}
\SG{1.036}{-0.296}{1.042}{-0.293}{0.086}
\SG{1.042}{-0.293}{1.049}{-0.291}{0.085}
\SG{1.049}{-0.291}{1.055}{-0.288}{0.082}
\SG{1.055}{-0.288}{1.061}{-0.285}{0.079}
\SG{1.061}{-0.285}{1.068}{-0.283}{0.080}
\SG{1.068}{-0.283}{1.074}{-0.280}{0.082}
\SG{1.074}{-0.280}{1.081}{-0.278}{0.091}
\SG{1.081}{-0.278}{1.087}{-0.275}{0.104}
\SG{1.087}{-0.275}{1.093}{-0.272}{0.112}
\SG{1.093}{-0.272}{1.100}{-0.270}{0.112}
\SG{1.100}{-0.270}{1.106}{-0.267}{0.112}
\SG{1.106}{-0.267}{1.112}{-0.264}{0.113}
\SG{1.112}{-0.264}{1.119}{-0.262}{0.116}
\SG{1.119}{-0.262}{1.125}{-0.259}{0.116}
\SG{1.125}{-0.259}{1.131}{-0.257}{0.113}
\SG{1.131}{-0.257}{1.138}{-0.254}{0.111}
\SG{1.138}{-0.254}{1.144}{-0.251}{0.112}
\SG{1.144}{-0.251}{1.150}{-0.249}{0.115}
\SG{1.150}{-0.249}{1.157}{-0.246}{0.126}
\SG{1.157}{-0.246}{1.163}{-0.243}{0.141}
\SG{1.163}{-0.243}{1.169}{-0.241}{0.147}
\SG{1.169}{-0.241}{1.176}{-0.238}{0.149}
\SG{1.176}{-0.238}{1.182}{-0.235}{0.148}
\SG{1.182}{-0.235}{1.188}{-0.233}{0.151}
\SG{1.188}{-0.233}{1.195}{-0.230}{0.155}
\SG{1.195}{-0.230}{1.201}{-0.228}{0.152}
\SG{1.201}{-0.228}{1.207}{-0.225}{0.149}
\SG{1.207}{-0.225}{1.214}{-0.222}{0.146}
\SG{1.214}{-0.222}{1.220}{-0.220}{0.146}
\SG{1.220}{-0.220}{1.226}{-0.217}{0.147}
\SG{1.226}{-0.217}{1.233}{-0.214}{0.146}
\SG{1.233}{-0.214}{1.239}{-0.212}{0.143}
\SG{1.239}{-0.212}{1.245}{-0.209}{0.145}
\SG{1.245}{-0.209}{1.252}{-0.207}{0.147}
\SG{1.252}{-0.207}{1.258}{-0.204}{0.154}
\SG{1.258}{-0.204}{1.265}{-0.201}{0.166}
\SG{1.265}{-0.201}{1.271}{-0.199}{0.180}
\SG{1.271}{-0.199}{1.277}{-0.196}{0.183}
\SG{1.277}{-0.196}{1.284}{-0.193}{0.184}
\SG{1.284}{-0.193}{1.290}{-0.191}{0.185}
\SG{1.290}{-0.191}{1.296}{-0.188}{0.189}
\SG{1.296}{-0.188}{1.303}{-0.185}{0.191}
\SG{1.303}{-0.185}{1.309}{-0.183}{0.190}
\SG{1.309}{-0.183}{1.315}{-0.180}{0.189}
\SG{1.315}{-0.180}{1.322}{-0.178}{0.191}
\SG{1.322}{-0.178}{1.328}{-0.175}{0.194}
\SG{1.328}{-0.175}{1.334}{-0.172}{0.204}
\SG{1.334}{-0.172}{1.341}{-0.170}{0.219}
\SG{1.341}{-0.170}{1.347}{-0.167}{0.232}
\SG{1.347}{-0.167}{1.353}{-0.164}{0.235}
\SG{1.353}{-0.164}{1.360}{-0.162}{0.235}
\SG{1.360}{-0.162}{1.366}{-0.159}{0.239}
\SG{1.366}{-0.159}{1.372}{-0.157}{0.244}
\SG{1.372}{-0.157}{1.379}{-0.154}{0.246}
\SG{1.379}{-0.154}{1.385}{-0.151}{0.245}
\SG{1.385}{-0.151}{1.391}{-0.149}{0.243}
\SG{1.391}{-0.149}{1.398}{-0.146}{0.244}
\SG{1.398}{-0.146}{1.404}{-0.143}{0.247}
\SG{1.404}{-0.143}{1.410}{-0.141}{0.248}
\SG{1.410}{-0.141}{1.417}{-0.138}{0.248}
\SG{1.417}{-0.138}{1.423}{-0.135}{0.250}
\SG{1.423}{-0.135}{1.430}{-0.133}{0.254}
\SG{1.430}{-0.133}{1.436}{-0.130}{0.260}
\SG{1.436}{-0.130}{1.442}{-0.128}{0.273}
\SG{1.442}{-0.128}{1.449}{-0.125}{0.290}
\SG{1.449}{-0.125}{1.455}{-0.122}{0.298}
\SG{1.455}{-0.122}{1.461}{-0.120}{0.302}
\SG{1.461}{-0.120}{1.468}{-0.117}{0.305}
\SG{1.468}{-0.117}{1.474}{-0.114}{0.311}
\SG{1.474}{-0.114}{1.480}{-0.112}{0.317}
\SG{1.480}{-0.112}{1.487}{-0.109}{0.319}
\SG{1.487}{-0.109}{1.493}{-0.107}{0.320}
\SG{1.493}{-0.107}{1.499}{-0.104}{0.324}
\SG{1.499}{-0.104}{1.506}{-0.101}{0.330}
\SG{1.506}{-0.101}{1.512}{-0.099}{0.339}
\SG{1.512}{-0.099}{1.518}{-0.096}{0.356}
\SG{1.518}{-0.096}{1.525}{-0.093}{0.376}
\SG{1.525}{-0.093}{1.531}{-0.091}{0.383}
\SG{1.531}{-0.091}{1.537}{-0.088}{0.388}
\SG{1.537}{-0.088}{1.544}{-0.086}{0.393}
\SG{1.544}{-0.086}{1.550}{-0.083}{0.401}
\SG{1.550}{-0.083}{1.556}{-0.080}{0.410}
\SG{1.556}{-0.080}{1.563}{-0.078}{0.414}
\SG{1.563}{-0.078}{1.569}{-0.075}{0.417}
\SG{1.569}{-0.075}{1.575}{-0.072}{0.424}
\SG{1.575}{-0.072}{1.582}{-0.070}{0.433}
\SG{1.582}{-0.070}{1.588}{-0.067}{0.448}
\SG{1.588}{-0.067}{1.595}{-0.064}{0.470}
\SG{1.595}{-0.064}{1.601}{-0.062}{0.493}
\SG{1.601}{-0.062}{1.607}{-0.059}{0.503}
\SG{1.607}{-0.059}{1.614}{-0.057}{0.511}
\SG{1.614}{-0.057}{1.620}{-0.054}{0.522}
\SG{1.620}{-0.054}{1.626}{-0.051}{0.536}
\SG{1.626}{-0.051}{1.633}{-0.049}{0.548}
\SG{1.633}{-0.049}{1.639}{-0.046}{0.556}
\SG{1.639}{-0.046}{1.645}{-0.043}{0.562}
\SG{1.645}{-0.043}{1.652}{-0.041}{0.575}
\SG{1.652}{-0.041}{1.658}{-0.038}{0.589}
\SG{1.658}{-0.038}{1.664}{-0.036}{0.603}
\SG{1.664}{-0.036}{1.671}{-0.033}{0.615}
\SG{1.671}{-0.033}{1.677}{-0.030}{0.629}
\SG{1.677}{-0.030}{1.683}{-0.028}{0.648}
\SG{1.683}{-0.028}{1.690}{-0.025}{0.670}
\SG{1.690}{-0.025}{1.696}{-0.022}{0.701}
\SG{1.696}{-0.022}{1.702}{-0.020}{0.737}
\SG{1.702}{-0.020}{1.709}{-0.017}{0.764}
\SG{1.709}{-0.017}{1.715}{-0.014}{0.783}
\SG{1.715}{-0.014}{1.721}{-0.012}{0.803}
\SG{1.721}{-0.012}{1.728}{-0.009}{0.824}
\SG{1.728}{-0.009}{1.734}{-0.007}{0.850}
\SG{1.734}{-0.007}{1.740}{-0.004}{0.830}
\SG{1.740}{-0.004}{1.744}{-0.001}{0.869}
\SG{1.744}{-0.001}{1.744}{0.001}{1.000}
\SG{1.744}{0.001}{1.740}{0.004}{0.854}
\SG{1.740}{0.004}{1.734}{0.007}{0.792}
\SG{1.734}{0.007}{1.728}{0.009}{0.830}
\SG{1.728}{0.009}{1.721}{0.012}{0.850}
\SG{1.721}{0.012}{1.715}{0.014}{0.824}
\SG{1.715}{0.014}{1.709}{0.017}{0.803}
\SG{1.709}{0.017}{1.702}{0.020}{0.783}
\SG{1.702}{0.020}{1.696}{0.022}{0.764}
\SG{1.696}{0.022}{1.690}{0.025}{0.737}
\SG{1.690}{0.025}{1.683}{0.028}{0.701}
\SG{1.683}{0.028}{1.677}{0.030}{0.670}
\SG{1.677}{0.030}{1.671}{0.033}{0.648}
\SG{1.671}{0.033}{1.664}{0.036}{0.629}
\SG{1.664}{0.036}{1.658}{0.038}{0.615}
\SG{1.658}{0.038}{1.652}{0.041}{0.603}
\SG{1.652}{0.041}{1.645}{0.043}{0.589}
\SG{1.645}{0.043}{1.639}{0.046}{0.575}
\SG{1.639}{0.046}{1.633}{0.049}{0.562}
\SG{1.633}{0.049}{1.626}{0.051}{0.556}
\SG{1.626}{0.051}{1.620}{0.054}{0.548}
\SG{1.620}{0.054}{1.614}{0.057}{0.536}
\SG{1.614}{0.057}{1.607}{0.059}{0.522}
\SG{1.607}{0.059}{1.601}{0.062}{0.511}
\SG{1.601}{0.062}{1.595}{0.064}{0.503}
\SG{1.595}{0.064}{1.588}{0.067}{0.493}
\SG{1.588}{0.067}{1.582}{0.070}{0.470}
\SG{1.582}{0.070}{1.575}{0.072}{0.448}
\SG{1.575}{0.072}{1.569}{0.075}{0.433}
\SG{1.569}{0.075}{1.563}{0.078}{0.424}
\SG{1.563}{0.078}{1.556}{0.080}{0.417}
\SG{1.556}{0.080}{1.550}{0.083}{0.414}
\SG{1.550}{0.083}{1.544}{0.086}{0.410}
\SG{1.544}{0.086}{1.537}{0.088}{0.401}
\SG{1.537}{0.088}{1.531}{0.091}{0.393}
\SG{1.531}{0.091}{1.525}{0.093}{0.388}
\SG{1.525}{0.093}{1.518}{0.096}{0.383}
\SG{1.518}{0.096}{1.512}{0.099}{0.376}
\SG{1.512}{0.099}{1.506}{0.101}{0.356}
\SG{1.506}{0.101}{1.499}{0.104}{0.339}
\SG{1.499}{0.104}{1.493}{0.107}{0.330}
\SG{1.493}{0.107}{1.487}{0.109}{0.324}
\SG{1.487}{0.109}{1.480}{0.112}{0.320}
\SG{1.480}{0.112}{1.474}{0.114}{0.319}
\SG{1.474}{0.114}{1.468}{0.117}{0.317}
\SG{1.468}{0.117}{1.461}{0.120}{0.311}
\SG{1.461}{0.120}{1.455}{0.122}{0.305}
\SG{1.455}{0.122}{1.449}{0.125}{0.302}
\SG{1.449}{0.125}{1.442}{0.128}{0.298}
\SG{1.442}{0.128}{1.436}{0.130}{0.290}
\SG{1.436}{0.130}{1.430}{0.133}{0.273}
\SG{1.430}{0.133}{1.423}{0.135}{0.260}
\SG{1.423}{0.135}{1.417}{0.138}{0.254}
\SG{1.417}{0.138}{1.410}{0.141}{0.250}
\SG{1.410}{0.141}{1.404}{0.143}{0.248}
\SG{1.404}{0.143}{1.398}{0.146}{0.248}
\SG{1.398}{0.146}{1.391}{0.149}{0.247}
\SG{1.391}{0.149}{1.385}{0.151}{0.244}
\SG{1.385}{0.151}{1.379}{0.154}{0.243}
\SG{1.379}{0.154}{1.372}{0.157}{0.245}
\SG{1.372}{0.157}{1.366}{0.159}{0.246}
\SG{1.366}{0.159}{1.360}{0.162}{0.244}
\SG{1.360}{0.162}{1.353}{0.164}{0.239}
\SG{1.353}{0.164}{1.347}{0.167}{0.235}
\SG{1.347}{0.167}{1.341}{0.170}{0.235}
\SG{1.341}{0.170}{1.334}{0.172}{0.232}
\SG{1.334}{0.172}{1.328}{0.175}{0.219}
\SG{1.328}{0.175}{1.322}{0.178}{0.204}
\SG{1.322}{0.178}{1.315}{0.180}{0.194}
\SG{1.315}{0.180}{1.309}{0.183}{0.191}
\SG{1.309}{0.183}{1.303}{0.185}{0.189}
\SG{1.303}{0.185}{1.296}{0.188}{0.190}
\SG{1.296}{0.188}{1.290}{0.191}{0.191}
\SG{1.290}{0.191}{1.284}{0.193}{0.189}
\SG{1.284}{0.193}{1.277}{0.196}{0.185}
\SG{1.277}{0.196}{1.271}{0.199}{0.184}
\SG{1.271}{0.199}{1.265}{0.201}{0.183}
\SG{1.265}{0.201}{1.258}{0.204}{0.180}
\SG{1.258}{0.204}{1.252}{0.207}{0.166}
\SG{1.252}{0.207}{1.245}{0.209}{0.154}
\SG{1.245}{0.209}{1.239}{0.212}{0.147}
\SG{1.239}{0.212}{1.233}{0.214}{0.145}
\SG{1.233}{0.214}{1.226}{0.217}{0.143}
\SG{1.226}{0.217}{1.220}{0.220}{0.146}
\SG{1.220}{0.220}{1.214}{0.222}{0.147}
\SG{1.214}{0.222}{1.207}{0.225}{0.146}
\SG{1.207}{0.225}{1.201}{0.228}{0.146}
\SG{1.201}{0.228}{1.195}{0.230}{0.149}
\SG{1.195}{0.230}{1.188}{0.233}{0.152}
\SG{1.188}{0.233}{1.182}{0.235}{0.155}
\SG{1.182}{0.235}{1.176}{0.238}{0.151}
\SG{1.176}{0.238}{1.169}{0.241}{0.148}
\SG{1.169}{0.241}{1.163}{0.243}{0.149}
\SG{1.163}{0.243}{1.157}{0.246}{0.147}
\SG{1.157}{0.246}{1.150}{0.249}{0.141}
\SG{1.150}{0.249}{1.144}{0.251}{0.126}
\SG{1.144}{0.251}{1.138}{0.254}{0.115}
\SG{1.138}{0.254}{1.131}{0.257}{0.112}
\SG{1.131}{0.257}{1.125}{0.259}{0.111}
\SG{1.125}{0.259}{1.119}{0.262}{0.113}
\SG{1.119}{0.262}{1.112}{0.264}{0.116}
\SG{1.112}{0.264}{1.106}{0.267}{0.116}
\SG{1.106}{0.267}{1.100}{0.270}{0.113}
\SG{1.100}{0.270}{1.093}{0.272}{0.112}
\SG{1.093}{0.272}{1.087}{0.275}{0.112}
\SG{1.087}{0.275}{1.081}{0.278}{0.112}
\SG{1.081}{0.278}{1.074}{0.280}{0.104}
\SG{1.074}{0.280}{1.068}{0.283}{0.091}
\SG{1.068}{0.283}{1.061}{0.285}{0.082}
\SG{1.061}{0.285}{1.055}{0.288}{0.080}
\SG{1.055}{0.288}{1.049}{0.291}{0.079}
\SG{1.049}{0.291}{1.042}{0.293}{0.082}
\SG{1.042}{0.293}{1.036}{0.296}{0.085}
\SG{1.036}{0.296}{1.030}{0.299}{0.086}
\SG{1.030}{0.299}{1.023}{0.301}{0.086}
\SG{1.023}{0.301}{1.017}{0.304}{0.087}
\SG{1.017}{0.304}{1.011}{0.307}{0.094}
\SG{1.011}{0.307}{1.004}{0.309}{0.097}
\SG{1.004}{0.309}{0.998}{0.312}{0.096}
\SG{0.998}{0.312}{0.992}{0.314}{0.094}
\SG{0.992}{0.314}{0.985}{0.317}{0.094}
\SG{0.985}{0.317}{0.979}{0.320}{0.095}
\SG{0.979}{0.320}{0.973}{0.322}{0.095}
\SG{0.973}{0.322}{0.966}{0.325}{0.080}
\SG{0.966}{0.325}{0.960}{0.328}{0.068}
\SG{0.960}{0.328}{0.954}{0.330}{0.063}
\SG{0.954}{0.330}{0.947}{0.333}{0.062}
\SG{0.947}{0.333}{0.941}{0.335}{0.064}
\SG{0.941}{0.335}{0.935}{0.338}{0.068}
\SG{0.935}{0.338}{0.928}{0.341}{0.071}
\SG{0.928}{0.341}{0.922}{0.343}{0.070}
\SG{0.922}{0.343}{0.916}{0.346}{0.068}
\SG{0.916}{0.346}{0.909}{0.349}{0.070}
\SG{0.909}{0.349}{0.903}{0.351}{0.071}
\SG{0.903}{0.351}{0.896}{0.354}{0.069}
\SG{0.896}{0.354}{0.890}{0.356}{0.055}
\SG{0.890}{0.356}{0.884}{0.359}{0.045}
\SG{0.884}{0.359}{0.877}{0.362}{0.042}
\SG{0.877}{0.362}{0.871}{0.364}{0.042}
\SG{0.871}{0.364}{0.865}{0.367}{0.044}
\SG{0.865}{0.367}{0.858}{0.370}{0.048}
\SG{0.858}{0.370}{0.852}{0.372}{0.051}
\SG{0.852}{0.372}{0.846}{0.375}{0.049}
\SG{0.846}{0.375}{0.839}{0.378}{0.048}
\SG{0.839}{0.378}{0.833}{0.380}{0.050}
\SG{0.833}{0.380}{0.827}{0.383}{0.050}
\SG{0.827}{0.383}{0.820}{0.385}{0.046}
\SG{0.820}{0.385}{0.814}{0.388}{0.034}
\SG{0.814}{0.388}{0.808}{0.391}{0.025}
\SG{0.808}{0.391}{0.801}{0.393}{0.024}
\SG{0.801}{0.393}{0.795}{0.396}{0.024}
\SG{0.795}{0.396}{0.789}{0.399}{0.026}
\SG{0.789}{0.399}{0.782}{0.401}{0.030}
\SG{0.782}{0.401}{0.776}{0.404}{0.033}
\SG{0.776}{0.404}{0.770}{0.406}{0.033}
\SG{0.770}{0.406}{0.763}{0.409}{0.036}
\SG{0.763}{0.409}{0.757}{0.412}{0.042}
\SG{0.757}{0.412}{0.751}{0.414}{0.046}
\SG{0.751}{0.414}{0.744}{0.417}{0.047}
\SG{0.744}{0.417}{0.738}{0.420}{0.045}
\SG{0.738}{0.420}{0.731}{0.422}{0.046}
\SG{0.731}{0.422}{0.725}{0.425}{0.048}
\SG{0.725}{0.425}{0.719}{0.428}{0.048}
\SG{0.719}{0.428}{0.712}{0.430}{0.037}
\SG{0.712}{0.430}{0.706}{0.433}{0.026}
\SG{0.706}{0.433}{0.700}{0.435}{0.020}
\SG{0.700}{0.435}{0.693}{0.438}{0.020}
\SG{0.693}{0.438}{0.687}{0.441}{0.020}
\SG{0.687}{0.441}{0.681}{0.443}{0.026}
\SG{0.681}{0.443}{0.674}{0.446}{0.029}
\SG{0.674}{0.446}{0.668}{0.449}{0.029}
\SG{0.668}{0.449}{0.662}{0.451}{0.028}
\SG{0.662}{0.451}{0.655}{0.454}{0.030}
\SG{0.655}{0.454}{0.649}{0.456}{0.031}
\SG{0.649}{0.456}{0.643}{0.459}{0.031}
\SG{0.643}{0.459}{0.636}{0.462}{0.019}
\SG{0.636}{0.462}{0.630}{0.464}{0.009}
\SG{0.630}{0.464}{0.624}{0.467}{0.006}
\SG{0.624}{0.467}{0.617}{0.470}{0.006}
\SG{0.617}{0.470}{0.611}{0.472}{0.007}
\SG{0.611}{0.472}{0.605}{0.475}{0.012}
\SG{0.605}{0.475}{0.598}{0.478}{0.016}
\SG{0.598}{0.478}{0.592}{0.480}{0.017}
\SG{0.592}{0.480}{0.586}{0.483}{0.019}
\SG{0.586}{0.483}{0.579}{0.485}{0.025}
\SG{0.579}{0.485}{0.573}{0.488}{0.030}
\SG{0.573}{0.488}{0.566}{0.491}{0.034}
\SG{0.566}{0.491}{0.560}{0.493}{0.032}
\SG{0.560}{0.493}{0.554}{0.496}{0.032}
\SG{0.554}{0.496}{0.547}{0.499}{0.036}
\SG{0.547}{0.499}{0.541}{0.501}{0.036}
\SG{0.541}{0.501}{0.535}{0.504}{0.030}
\SG{0.535}{0.504}{0.528}{0.506}{0.018}
\SG{0.528}{0.506}{0.522}{0.509}{0.009}
\SG{0.522}{0.509}{0.516}{0.512}{0.009}
\SG{0.516}{0.512}{0.509}{0.514}{0.010}
\SG{0.509}{0.514}{0.503}{0.517}{0.014}
\SG{0.503}{0.517}{0.497}{0.520}{0.019}
\SG{0.497}{0.520}{0.490}{0.522}{0.021}
\SG{0.490}{0.522}{0.484}{0.525}{0.020}
\SG{0.484}{0.525}{0.478}{0.528}{0.020}
\SG{0.478}{0.528}{0.471}{0.530}{0.023}
\SG{0.471}{0.530}{0.465}{0.533}{0.024}
\SG{0.465}{0.533}{0.459}{0.535}{0.017}
\SG{0.459}{0.535}{0.452}{0.538}{0.006}
\SG{0.452}{0.538}{0.446}{0.541}{0.000}
\SG{0.446}{0.541}{0.440}{0.543}{0.000}
\SG{0.440}{0.543}{0.433}{0.546}{0.001}
\SG{0.433}{0.546}{0.427}{0.549}{0.006}
\SG{0.427}{0.549}{0.421}{0.551}{0.011}
\SG{0.421}{0.551}{0.414}{0.554}{0.014}
\SG{0.414}{0.554}{0.408}{0.556}{0.015}
\SG{0.408}{0.556}{0.401}{0.559}{0.020}
\SG{0.401}{0.559}{0.395}{0.562}{0.027}
\SG{0.395}{0.562}{0.389}{0.564}{0.032}
\SG{0.389}{0.564}{0.382}{0.567}{0.032}
\SG{0.382}{0.567}{0.376}{0.570}{0.032}
\SG{0.376}{0.570}{0.370}{0.572}{0.035}
\SG{0.370}{0.572}{0.363}{0.575}{0.037}
\SG{0.363}{0.575}{0.357}{0.577}{0.038}
\SG{0.357}{0.577}{0.351}{0.580}{0.024}
\SG{0.351}{0.580}{0.344}{0.583}{0.014}
\SG{0.344}{0.583}{0.338}{0.585}{0.011}
\SG{0.338}{0.585}{0.332}{0.588}{0.012}
\SG{0.332}{0.588}{0.325}{0.591}{0.016}
\SG{0.325}{0.591}{0.319}{0.593}{0.022}
\SG{0.319}{0.593}{0.313}{0.596}{0.027}
\SG{0.313}{0.596}{0.306}{0.599}{0.027}
\SG{0.306}{0.599}{0.300}{0.601}{0.027}
\SG{0.300}{0.601}{0.294}{0.604}{0.031}
\SG{0.294}{0.604}{0.287}{0.606}{0.033}
\SG{0.287}{0.606}{0.281}{0.609}{0.031}
\SG{0.281}{0.609}{0.275}{0.612}{0.019}
\SG{0.275}{0.612}{0.268}{0.614}{0.011}
\SG{0.268}{0.614}{0.262}{0.617}{0.011}
\SG{0.262}{0.617}{0.256}{0.620}{0.013}
\SG{0.256}{0.620}{0.249}{0.622}{0.016}
\SG{0.249}{0.622}{0.243}{0.625}{0.022}
\SG{0.243}{0.625}{0.236}{0.627}{0.026}
\SG{0.236}{0.627}{0.230}{0.630}{0.026}
\SG{0.230}{0.630}{0.224}{0.633}{0.027}
\SG{0.224}{0.633}{0.217}{0.635}{0.030}
\SG{0.217}{0.635}{0.211}{0.638}{0.032}
\SG{0.211}{0.638}{0.205}{0.641}{0.028}
\SG{0.205}{0.641}{0.198}{0.643}{0.017}
\SG{0.198}{0.643}{0.192}{0.646}{0.011}
\SG{0.192}{0.646}{0.186}{0.649}{0.012}
\SG{0.186}{0.649}{0.179}{0.651}{0.013}
\SG{0.179}{0.651}{0.173}{0.654}{0.018}
\SG{0.173}{0.654}{0.167}{0.656}{0.024}
\SG{0.167}{0.656}{0.160}{0.659}{0.028}
\SG{0.160}{0.659}{0.154}{0.662}{0.030}
\SG{0.154}{0.662}{0.148}{0.664}{0.034}
\SG{0.148}{0.664}{0.141}{0.667}{0.042}
\SG{0.141}{0.667}{0.135}{0.670}{0.048}
\SG{0.135}{0.670}{0.129}{0.672}{0.050}
\SG{0.129}{0.672}{0.122}{0.675}{0.050}
\SG{0.122}{0.675}{0.116}{0.677}{0.053}
\SG{0.116}{0.677}{0.110}{0.680}{0.056}
\SG{0.110}{0.680}{0.103}{0.683}{0.058}
\SG{0.103}{0.683}{0.097}{0.685}{0.047}
\SG{0.097}{0.685}{0.091}{0.688}{0.038}
\SG{0.091}{0.688}{0.084}{0.691}{0.035}
\SG{0.084}{0.691}{0.078}{0.693}{0.037}
\SG{0.078}{0.693}{0.071}{0.696}{0.039}
\SG{0.071}{0.696}{0.065}{0.699}{0.046}
\SG{0.065}{0.699}{0.059}{0.701}{0.052}
\SG{0.059}{0.701}{0.052}{0.704}{0.053}
\SG{0.052}{0.704}{0.046}{0.706}{0.054}
\SG{0.046}{0.706}{0.040}{0.709}{0.057}
\SG{0.040}{0.709}{0.033}{0.712}{0.060}
\SG{0.033}{0.712}{0.027}{0.714}{0.062}
\SG{0.027}{0.714}{0.021}{0.717}{0.051}
\SG{0.021}{0.717}{0.014}{0.720}{0.043}
\SG{0.014}{0.720}{0.008}{0.722}{0.042}
\SG{0.008}{0.722}{0.002}{0.725}{0.045}
\SG{0.002}{0.725}{-0.005}{0.727}{0.048}
\SG{-0.005}{0.727}{-0.011}{0.730}{0.055}
\SG{-0.011}{0.730}{-0.017}{0.733}{0.061}
\SG{-0.017}{0.733}{-0.024}{0.735}{0.064}
\SG{-0.024}{0.735}{-0.030}{0.738}{0.068}
\SG{-0.030}{0.738}{-0.036}{0.741}{0.076}
\SG{-0.036}{0.741}{-0.043}{0.743}{0.084}
\SG{-0.043}{0.743}{-0.049}{0.746}{0.089}
\SG{-0.049}{0.746}{-0.055}{0.749}{0.089}
\SG{-0.055}{0.749}{-0.062}{0.751}{0.091}
\SG{-0.062}{0.751}{-0.068}{0.754}{0.097}
\SG{-0.068}{0.754}{-0.074}{0.756}{0.099}
\SG{-0.074}{0.756}{-0.081}{0.759}{0.095}
\SG{-0.081}{0.759}{-0.087}{0.762}{0.085}
\SG{-0.087}{0.762}{-0.093}{0.764}{0.079}
\SG{-0.093}{0.764}{-0.100}{0.767}{0.081}
\SG{-0.100}{0.767}{-0.106}{0.770}{0.085}
\SG{-0.106}{0.770}{-0.113}{0.772}{0.092}
\SG{-0.113}{0.772}{-0.119}{0.775}{0.099}
\SG{-0.119}{0.775}{-0.125}{0.777}{0.104}
\SG{-0.125}{0.777}{-0.132}{0.780}{0.106}
\SG{-0.132}{0.780}{-0.138}{0.783}{0.108}
\SG{-0.138}{0.783}{-0.144}{0.785}{0.114}
\SG{-0.144}{0.785}{-0.151}{0.788}{0.118}
\SG{-0.151}{0.788}{-0.157}{0.791}{0.112}
\SG{-0.157}{0.791}{-0.163}{0.793}{0.104}
\SG{-0.163}{0.793}{-0.170}{0.796}{0.102}
\SG{-0.170}{0.796}{-0.176}{0.798}{0.106}
\SG{-0.176}{0.798}{-0.182}{0.801}{0.110}
\SG{-0.182}{0.801}{-0.189}{0.804}{0.119}
\SG{-0.189}{0.804}{-0.195}{0.806}{0.127}
\SG{-0.195}{0.806}{-0.201}{0.809}{0.132}
\SG{-0.201}{0.809}{-0.208}{0.812}{0.138}
\SG{-0.208}{0.812}{-0.214}{0.814}{0.147}
\SG{-0.214}{0.814}{-0.220}{0.817}{0.158}
\SG{-0.220}{0.817}{-0.227}{0.820}{0.167}
\SG{-0.227}{0.820}{-0.233}{0.822}{0.171}
\SG{-0.233}{0.822}{-0.239}{0.825}{0.175}
\SG{-0.239}{0.825}{-0.246}{0.827}{0.183}
\SG{-0.246}{0.827}{-0.251}{0.830}{0.189}
\SG{-0.251}{0.830}{-0.256}{0.832}{0.193}
\SG{-0.256}{0.832}{-0.259}{0.833}{0.187}

    \end{scope}
    \draw[azulclaro, line width=0.45pt, rounded corners=3pt] (-2.6,-2.0) rectangle (2.6,2.0);
    \node[anchor=north, align=center, text width=5.3cm, inner sep=0, font=\scriptsize,
          text=grispizarra] at (0,-2.2)
      {Líneas de campo y equipotenciales de un conductor cargado con punta;
       el color del contorno indica la densidad superficial $\sigma$.};
  \end{scope}

  % --- Palabras clave y clasificación -----------------------------------
  \node[anchor=north west, align=justify, inner sep=0, text width=16cm, font=\small] at ($(O)+(2.5,11.0)$)
    {{\scshape\color{azulmedio}\textls[60]{Palabras clave}}\enspace
     conductor perfecto · equilibrio electrostático · efecto punta · presión electrostática ·
     jaula de Faraday · coeficientes de potencial · coeficientes de capacidad e influencia ·
     influencia total · condensador\\[4pt]
     {\scshape\color{azulmedio}\textls[60]{Clasificación PACS}}\enspace
     41.20.Cv (Electrostática; ecuaciones de Poisson y Laplace, problemas de contorno)};

  % --- Expresiones centrales del tema ---------------------------------------
  \draw[azulclaro, line width=0.45pt] ($(O)+(2.5,8.55)$) -- ($(O)+(18.5,8.55)$);
  \node[anchor=north, align=center, inner sep=0, text width=16cm] at ($(O)+(10.5,8.1)$)
    {\color{azulprof}
     $\displaystyle\vE\big|_{\Sc^{+}}=\frac{\sigma}{\eo}\,\un\qquad
      \Pel=\frac{\sigma^{2}}{2\eo}\qquad
      \phi_i=\sum_{j=1}^{N}p_{ij}\,Q_j\qquad
      Q_i=\sum_{j=1}^{N}c_{ij}\,\phi_j\qquad
      C=\frac{1}{p_{11}+p_{22}-2p_{12}}$};
  \draw[azulclaro, line width=0.45pt] ($(O)+(2.5,6.25)$) -- ($(O)+(18.5,6.25)$);

  % --- Tipo de documento, lugar y curso ------------------------------------
  \node[anchor=south, align=center, inner sep=0, text width=16cm] at ($(O)+(10.5,2.2)$)
    {{\color{azulprof}\scshape\large\textls[100]{\tipotrabajo}}\\[5pt]
     {\color{grispizarra}\scshape\small\textls[80]{\miUbicacion}}};
\end{tikzpicture}
\tikzexternalenable
\end{titlepage}
